% ═══════════════════════════════════════════════════════════════════════════════
% COURS DE MATHÉMATIQUES — TERMINALE GÉNÉRALE SPÉCIALITÉ
% Programme officiel (BO spécial n°1 du 22 janvier 2019)
% Compilable directement sur Overleaf (pdflatex)
% ═══════════════════════════════════════════════════════════════════════════════

\documentclass[12pt,a4paper,twoside]{book}

% ── Encodage et langue ──
\usepackage[utf8]{inputenc}
\usepackage[T1]{fontenc}
\usepackage[french]{babel}
\usepackage{lmodern}

% ── Mise en page ──
\usepackage[top=2.5cm,bottom=2.5cm,left=2.5cm,right=2.5cm]{geometry}
\usepackage{fancyhdr}
\usepackage{titlesec}
\usepackage{enumitem}
\usepackage{multicol}
\usepackage{array}
\usepackage{booktabs}

% ── Mathématiques ──
\usepackage{amsmath,amssymb,amsfonts}
\usepackage{mathtools}
\usepackage{cancel}
\usepackage{stmaryrd} % pour \llbracket, \rrbracket

% ── Couleurs et encadrés ──
\usepackage[dvipsnames,svgnames]{xcolor}
\usepackage[most]{tcolorbox}

% ── Figures ──
\usepackage{tikz}
\usepackage{pgfplots}
\pgfplotsset{compat=1.18}
\usetikzlibrary{arrows.meta,calc,patterns,decorations.markings,positioning,angles,quotes}

% ── Navigation ──
\usepackage[hidelinks,bookmarksnumbered]{hyperref}

% ── Divers ──
\usepackage{marginnote}
\usepackage{etoolbox}

% ═══════════════════════════════════════════════════════════════════════════════
% COULEURS DU DOCUMENT
% ═══════════════════════════════════════════════════════════════════════════════

\definecolor{indigo}{HTML}{4f46e5}
\definecolor{indigobg}{HTML}{eef2ff}
\definecolor{indigoborder}{HTML}{a5b4fc}
\definecolor{vertdef}{HTML}{16a34a}
\definecolor{vertbg}{HTML}{f0fdf4}
\definecolor{vertborder}{HTML}{86efac}
\definecolor{orangemeth}{HTML}{ea580c}
\definecolor{orangebg}{HTML}{fff7ed}
\definecolor{orangeborder}{HTML}{fdba74}
\definecolor{rougeerr}{HTML}{dc2626}
\definecolor{rougebg}{HTML}{fef2f2}
\definecolor{rougeborder}{HTML}{fca5a5}
\definecolor{bleuclair}{HTML}{0284c7}
\definecolor{bleubg}{HTML}{f0f9ff}
\definecolor{bleuborder}{HTML}{7dd3fc}
\definecolor{violethist}{HTML}{7c3aed}
\definecolor{violetbg}{HTML}{f5f3ff}
\definecolor{violetborder}{HTML}{c4b5fd}
\definecolor{grisfond}{HTML}{f8fafc}

% ═══════════════════════════════════════════════════════════════════════════════
% ENCADRÉS TCOLORBOX
% ═══════════════════════════════════════════════════════════════════════════════

% Définition
\newtcolorbox{definition}[1][]{%
  enhanced, breakable,
  colback=indigobg, colframe=indigo,
  fonttitle=\bfseries\sffamily,
  title={Définition\ifstrempty{#1}{}{ — #1}},
  attach boxed title to top left={xshift=5mm,yshift=-3mm},
  boxed title style={colback=indigo,colframe=indigo},
  boxed title size=title,
  top=4mm, left=3mm, right=3mm,
  arc=2mm
}

% Propriété / Théorème
\newtcolorbox{propriete}[1][]{%
  enhanced, breakable,
  colback=vertbg, colframe=vertdef,
  fonttitle=\bfseries\sffamily,
  title={Propriété\ifstrempty{#1}{}{ — #1}},
  attach boxed title to top left={xshift=5mm,yshift=-3mm},
  boxed title style={colback=vertdef,colframe=vertdef},
  boxed title size=title,
  top=4mm, left=3mm, right=3mm,
  arc=2mm
}

\newtcolorbox{theoreme}[1][]{%
  enhanced, breakable,
  colback=vertbg, colframe=vertdef,
  fonttitle=\bfseries\sffamily,
  title={Théorème\ifstrempty{#1}{}{ — #1}},
  attach boxed title to top left={xshift=5mm,yshift=-3mm},
  boxed title style={colback=vertdef,colframe=vertdef},
  boxed title size=title,
  top=4mm, left=3mm, right=3mm,
  arc=2mm
}

% Méthode
\newtcolorbox{methode}[1][]{%
  enhanced, breakable,
  colback=orangebg, colframe=orangemeth,
  fonttitle=\bfseries\sffamily,
  title={Méthode\ifstrempty{#1}{}{ — #1}},
  attach boxed title to top left={xshift=5mm,yshift=-3mm},
  boxed title style={colback=orangemeth,colframe=orangemeth},
  boxed title size=title,
  top=4mm, left=3mm, right=3mm,
  arc=2mm
}

% Exemple
\newtcolorbox{exemple}[1][]{%
  enhanced, breakable,
  colback=bleubg, colframe=bleuclair,
  fonttitle=\bfseries\sffamily,
  title={Exemple\ifstrempty{#1}{}{ — #1}},
  attach boxed title to top left={xshift=5mm,yshift=-3mm},
  boxed title style={colback=bleuclair,colframe=bleuclair},
  boxed title size=title,
  top=4mm, left=3mm, right=3mm,
  arc=2mm
}

% Attention / Erreur fréquente
\newtcolorbox{attention}[1][]{%
  enhanced, breakable,
  colback=rougebg, colframe=rougeerr,
  fonttitle=\bfseries\sffamily,
  title={Attention\ifstrempty{#1}{}{ — #1}},
  attach boxed title to top left={xshift=5mm,yshift=-3mm},
  boxed title style={colback=rougeerr,colframe=rougeerr},
  boxed title size=title,
  top=4mm, left=3mm, right=3mm,
  arc=2mm
}

% Histoire des mathématiques
\newtcolorbox{histoire}[1][]{%
  enhanced, breakable,
  colback=violetbg, colframe=violethist,
  fonttitle=\bfseries\sffamily,
  title={Histoire des mathématiques\ifstrempty{#1}{}{ — #1}},
  attach boxed title to top left={xshift=5mm,yshift=-3mm},
  boxed title style={colback=violethist,colframe=violethist},
  boxed title size=title,
  top=4mm, left=3mm, right=3mm,
  arc=2mm
}

% À retenir (bilan)
\newtcolorbox{retenir}{%
  enhanced, breakable,
  colback=grisfond, colframe=indigo,
  fonttitle=\bfseries\sffamily\large,
  title={À retenir},
  attach boxed title to top center={yshift=-3mm},
  boxed title style={colback=indigo,colframe=indigo},
  boxed title size=title,
  top=6mm, left=4mm, right=4mm,
  arc=3mm,
  boxrule=2pt
}

% Démonstration
\newtcolorbox{demonstration}[1][]{%
  enhanced, breakable,
  colback=grisfond, colframe=gray!60,
  fonttitle=\bfseries\sffamily\itshape,
  title={Démonstration\ifstrempty{#1}{}{ — #1}},
  left=3mm, right=3mm, top=3mm,
  arc=1mm, boxrule=0.5pt
}

% ═══════════════════════════════════════════════════════════════════════════════
% EN-TÊTES ET PIEDS DE PAGE
% ═══════════════════════════════════════════════════════════════════════════════

\pagestyle{fancy}
\fancyhf{}
\fancyhead[LE]{\sffamily\small\leftmark}
\fancyhead[RO]{\sffamily\small\rightmark}
\fancyfoot[C]{\thepage}
\renewcommand{\headrulewidth}{0.4pt}

% ═══════════════════════════════════════════════════════════════════════════════
% TITRES
% ═══════════════════════════════════════════════════════════════════════════════

\titleformat{\chapter}[display]
  {\normalfont\sffamily\huge\bfseries\color{indigo}}
  {\chaptertitlename\ \thechapter}{20pt}{\Huge}
\titleformat{\section}
  {\normalfont\sffamily\Large\bfseries\color{indigo}}
  {\thesection}{1em}{}
\titleformat{\subsection}
  {\normalfont\sffamily\large\bfseries\color{indigo!80!black}}
  {\thesubsection}{1em}{}

% ═══════════════════════════════════════════════════════════════════════════════
% COMMANDES PERSONNALISÉES
% ═══════════════════════════════════════════════════════════════════════════════

\newcommand{\R}{\mathbb{R}}
\newcommand{\N}{\mathbb{N}}
\newcommand{\Z}{\mathbb{Z}}
\newcommand{\Q}{\mathbb{Q}}
\newcommand{\C}{\mathbb{C}}
\newcommand{\K}{\mathbb{K}}
\newcommand{\ds}{\displaystyle}
\newcommand{\vect}[1]{\overrightarrow{#1}}
\newcommand{\norme}[1]{\left\| #1 \right\|}
\newcommand{\abs}[1]{\left| #1 \right|}
\newcommand{\binom}[2]{\displaystyle\binom{#1}{#2}}
\DeclareMathOperator{\card}{card}

% ═══════════════════════════════════════════════════════════════════════════════
% DÉBUT DU DOCUMENT
% ═══════════════════════════════════════════════════════════════════════════════

\begin{document}

% ── Page de garde ──
\begin{titlepage}
\centering
\vspace*{3cm}
{\sffamily\Huge\bfseries\color{indigo} Mathématiques\\[0.3cm]}
{\sffamily\LARGE\color{indigo!80!black} Terminale générale — Spécialité\\[1.5cm]}
{\large\color{gray} Programme officiel\\Bulletin officiel spécial n°1 du 22 janvier 2019\\[2cm]}

\begin{tikzpicture}
  \draw[indigo,very thick] (0,0) -- (4,0) -- (4,3) -- (0,3) -- cycle;
  \draw[indigo,thick] (0,0) -- (1.5,1.5) -- (5.5,1.5) -- (4,0);
  \draw[indigo,thick] (4,3) -- (5.5,4.5) -- (5.5,1.5);
  \draw[indigo,thick] (0,3) -- (1.5,4.5) -- (5.5,4.5);
  \draw[indigo,dashed] (1.5,1.5) -- (1.5,4.5);
  \node[indigo] at (2,1.5) {\Large $\int_a^b f(x)\,\mathrm{d}x$};
  \node[indigo!70!black] at (3.5,4) {\large $\vec{u} \cdot \vec{v}$};
  \node[indigo!60!black] at (0.7,2.5) {\small $e^x$};
  \node[indigo!60!black] at (4.8,3) {\small $\ln x$};
  \node[indigo!50!black] at (3,0.5) {\small $\binom{n}{k}$};
\end{tikzpicture}

\vspace{2cm}
{\large\itshape Un cours complet, détaillé et illustré\\pour comprendre et réussir}
\vfill
{\small\color{gray} Année scolaire 2025–2026}
\end{titlepage}

% ── Table des matières ──
\tableofcontents
\newpage

% ── Avant-propos ──
\chapter*{Avant-propos}
\addcontentsline{toc}{chapter}{Avant-propos}

Ce livret couvre l'intégralité du programme de \textbf{spécialité mathématiques de Terminale générale}, tel que défini par le Bulletin officiel. Il est conçu pour accompagner les élèves tout au long de l'année et préparer efficacement le baccalauréat et les études supérieures.

\medskip

\textbf{Ce qui distingue ce livret :}
\begin{itemize}[leftmargin=*]
  \item Chaque notion est d'abord \textbf{motivée} : pourquoi l'étudie-t-on ? À quoi sert-elle ?
  \item Les définitions et théorèmes sont accompagnés d'\textbf{explications intuitives} avant toute formalisation.
  \item Des encadrés \textbf{historiques} éclairent l'origine des concepts.
  \item Les \textbf{erreurs fréquentes} sont signalées pour les éviter.
  \item Les \textbf{méthodes} sont présentées pas à pas, avec des exemples détaillés.
  \item Chaque chapitre se termine par un \textbf{bilan} synthétique.
\end{itemize}

\medskip

Le programme s'organise en quatre grandes parties :
\begin{enumerate}
  \item \textbf{Algèbre et géométrie} — Combinatoire, vecteurs et géométrie dans l'espace
  \item \textbf{Analyse} — Suites, limites, dérivation, continuité, logarithme, trigonométrie, primitives, intégrales
  \item \textbf{Probabilités} — Loi binomiale, variables aléatoires, loi des grands nombres
  \item \textbf{Algorithmique} — Notions transversales intégrées dans chaque chapitre
\end{enumerate}

\medskip

\textit{Bonne lecture et bon courage pour cette belle année de mathématiques.}

% ═══════════════════════════════════════════════════════════════════════════════
%
%  PARTIE I — ALGÈBRE ET GÉOMÉTRIE
%
% ═══════════════════════════════════════════════════════════════════════════════

\part{Algèbre et géométrie}

% ─────────────────────────────────────────────────────────────────────────────
% CHAPITRE 1 — COMBINATOIRE ET DÉNOMBREMENT
% ─────────────────────────────────────────────────────────────────────────────

\chapter{Combinatoire et dénombrement}

\section*{Introduction}

\textit{Combien de codes PIN à 4 chiffres peut-on former ? Combien de mains de 5 cartes peut-on distribuer dans un jeu de 52 cartes ? Combien d'équipes de 11 joueurs peut-on composer à partir d'un effectif de 23 ?}

\medskip

Ces questions, d'apparence simple, relèvent d'un domaine fondamental des mathématiques : la \textbf{combinatoire}, c'est-à-dire l'art de compter. Le dénombrement consiste à déterminer le nombre d'éléments d'un ensemble fini, souvent sans les lister tous. C'est un outil indispensable en probabilités, en informatique et dans de nombreuses sciences.

\medskip

\textbf{Objectifs du chapitre :}
\begin{itemize}[leftmargin=*]
  \item Utiliser les principes additif et multiplicatif pour dénombrer
  \item Connaître et manipuler les $k$-uplets, arrangements et permutations
  \item Calculer des coefficients binomiaux et utiliser le triangle de Pascal
  \item Résoudre des problèmes de dénombrement variés
\end{itemize}

\begin{histoire}[Des origines anciennes]
La combinatoire est l'une des branches les plus anciennes des mathématiques. Le \textbf{triangle de Pascal} était déjà connu des mathématiciens indiens (Pingala, II\ieme{} siècle av. J.-C.) et chinois (Yang Hui, XIII\ieme{} siècle). En Europe, c'est Blaise \textbf{Pascal} qui en systématise l'étude dans son \textit{Traité du triangle arithmétique} (1654), motivé par des questions de jeux de hasard. Au XIX\ieme{} siècle, des mathématiciens comme Lucas, Delannoy et Laisant développent les « mathématiques discrètes », aujourd'hui au cœur de l'informatique et de l'intelligence artificielle.
\end{histoire}

% ── I. Principes de dénombrement ──
\section{Principes de dénombrement}

\subsection{Ensembles finis et cardinal}

\begin{definition}[Cardinal d'un ensemble fini]
Un ensemble $E$ est \textbf{fini} s'il contient un nombre fini d'éléments. Ce nombre est appelé le \textbf{cardinal} de $E$, noté $\card(E)$ ou $|E|$.
\end{definition}

\begin{exemple}
L'ensemble $E = \{a, b, c, d\}$ a pour cardinal $\card(E) = 4$. L'ensemble des chiffres $\{0, 1, 2, \ldots, 9\}$ a pour cardinal $10$.
\end{exemple}

\subsection{Principe additif}

\begin{propriete}[Principe additif]
Si $A$ et $B$ sont deux ensembles finis \textbf{disjoints} (c'est-à-dire $A \cap B = \emptyset$), alors :
\[ \card(A \cup B) = \card(A) + \card(B) \]
Plus généralement, si $A_1, A_2, \ldots, A_n$ sont des ensembles finis deux à deux disjoints :
\[ \card(A_1 \cup A_2 \cup \cdots \cup A_n) = \card(A_1) + \card(A_2) + \cdots + \card(A_n) \]
\end{propriete}

\textit{Intuition :} si je dois choisir entre plusieurs options qui ne se chevauchent pas, je compte les options de chaque catégorie et j'additionne.

\begin{exemple}
Un club sportif propose 4 sports collectifs et 3 sports individuels. Un élève choisit un sport parmi ceux proposés. Nombre de choix possibles : $4 + 3 = 7$.
\end{exemple}

\begin{attention}[Ensembles non disjoints]
Si $A \cap B \neq \emptyset$, on ne peut pas simplement additionner. On utilise alors la formule :
\[ \card(A \cup B) = \card(A) + \card(B) - \card(A \cap B) \]
En soustrayant les éléments comptés deux fois.
\end{attention}

\subsection{Principe multiplicatif — Produit cartésien}

\begin{definition}[Produit cartésien]
Le \textbf{produit cartésien} de deux ensembles $A$ et $B$, noté $A \times B$, est l'ensemble de tous les \textbf{couples} $(a, b)$ avec $a \in A$ et $b \in B$ :
\[ A \times B = \{(a, b) \mid a \in A \text{ et } b \in B\} \]
\end{definition}

\begin{propriete}[Principe multiplicatif]
Si $A$ et $B$ sont deux ensembles finis, alors :
\[ \card(A \times B) = \card(A) \times \card(B) \]
Plus généralement :
\[ \card(A_1 \times A_2 \times \cdots \times A_n) = \card(A_1) \times \card(A_2) \times \cdots \times \card(A_n) \]
\end{propriete}

\textit{Intuition :} si une action se décompose en étapes successives indépendantes, le nombre total de possibilités est le produit du nombre de choix à chaque étape.

\begin{exemple}[Le menu du restaurant]
Un menu se compose d'une entrée parmi 3, d'un plat parmi 5 et d'un dessert parmi 4. Le nombre de menus possibles est :
\[ 3 \times 5 \times 4 = 60 \]
\end{exemple}

\begin{center}
\begin{tikzpicture}[level distance=1.5cm, sibling distance=2cm,
  every node/.style={font=\small},
  edge from parent/.style={draw,indigo,-{Stealth[length=3pt]}}]
  \node[indigo,font=\bfseries] {Menu}
    child {node {$E_1$}
      child {node {$P_1$}} child {node {$P_2$}} child {node {$P_3$}}
    }
    child {node {$E_2$}
      child {node {$P_1$}} child {node {$P_2$}} child {node {$P_3$}}
    }
    child {node {$E_3$}
      child {node {$P_1$}} child {node {$P_2$}} child {node {$P_3$}}
    };
\end{tikzpicture}
\captionof{figure}{Arbre des possibilités pour 3 entrées et 3 plats (sans les desserts)}
\end{center}

\subsection{$k$-uplets}

\begin{definition}[$k$-uplet]
Soient $A_1, A_2, \ldots, A_k$ des ensembles. Un \textbf{$k$-uplet} (ou \textbf{$k$-liste}) est un élément $(a_1, a_2, \ldots, a_k)$ du produit cartésien $A_1 \times A_2 \times \cdots \times A_k$.

L'\textbf{ordre compte} : le couple $(a, b)$ est différent du couple $(b, a)$ si $a \neq b$.
\end{definition}

\begin{propriete}[$k$-uplets d'un même ensemble]
Le nombre de $k$-uplets d'éléments d'un ensemble $A$ à $n$ éléments (avec répétitions possibles) est :
\[ \card(A^k) = n^k \]
\end{propriete}

\begin{exemple}[Code PIN]
Un code PIN est composé de 4 chiffres (de 0 à 9), chaque chiffre pouvant être répété. Nombre de codes possibles : $10^4 = 10\,000$.

\medskip

Un mot de 3 lettres dans un alphabet de 26 lettres (avec répétition) : $26^3 = 17\,576$ mots possibles.
\end{exemple}

% ── II. Nombre de parties d'un ensemble ──
\section{Nombre de parties d'un ensemble}

\begin{propriete}
Un ensemble à $n$ éléments possède $2^n$ parties (sous-ensembles).
\end{propriete}

\begin{methode}[Comprendre le $2^n$]
Pour construire une partie de $E = \{e_1, e_2, \ldots, e_n\}$, on décide pour chaque élément s'il est \textbf{dedans} ou \textbf{dehors}. Cela fait 2 choix par élément, soit $2^n$ parties au total.

On retrouve les $n$-uplets de $\{0, 1\}$ : chaque $n$-uplet code une partie (1 = l'élément est pris, 0 = il ne l'est pas).
\end{methode}

\begin{exemple}
L'ensemble $E = \{a, b, c\}$ possède $2^3 = 8$ parties :
\[ \emptyset,\ \{a\},\ \{b\},\ \{c\},\ \{a,b\},\ \{a,c\},\ \{b,c\},\ \{a,b,c\} \]
\end{exemple}

% ── III. Arrangements et permutations ──
\section{$k$-uplets d'éléments distincts et permutations}

\subsection{$k$-uplets d'éléments distincts}

\begin{definition}
Un \textbf{$k$-uplet d'éléments distincts} d'un ensemble $E$ à $n$ éléments est un $k$-uplet $(a_1, a_2, \ldots, a_k)$ où tous les $a_i$ sont \textbf{différents}. On parle aussi d'\textbf{arrangement} de $k$ éléments parmi $n$.
\end{definition}

\begin{propriete}
Le nombre de $k$-uplets d'éléments distincts d'un ensemble à $n$ éléments ($k \leqslant n$) est :
\[ n \times (n-1) \times (n-2) \times \cdots \times (n-k+1) = \frac{n!}{(n-k)!} \]
\end{propriete}

\begin{methode}[Raisonner étape par étape]
Pour le 1\textsuperscript{er} élément : $n$ choix. Pour le 2\textsuperscript{e} : $n-1$ choix (un élément déjà utilisé). Pour le 3\textsuperscript{e} : $n-2$ choix. Et ainsi de suite jusqu'au $k$\textsuperscript{e} : $n-k+1$ choix.
\end{methode}

\subsection{Factorielle}

\begin{definition}[Factorielle]
Pour tout entier naturel $n \geqslant 1$, on définit \textbf{$n$ factorielle} par :
\[ n! = n \times (n-1) \times (n-2) \times \cdots \times 2 \times 1 \]
Par convention, $0! = 1$.
\end{definition}

\begin{exemple}
$1! = 1$, \quad $3! = 6$, \quad $5! = 120$, \quad $10! = 3\,628\,800$.

\medskip

La croissance de la factorielle est extrêmement rapide : $20! \approx 2{,}4 \times 10^{18}$, ce qui dépasse le nombre de grains de sable sur Terre.
\end{exemple}

\subsection{Permutations}

\begin{definition}[Permutation]
Une \textbf{permutation} d'un ensemble $E$ à $n$ éléments est un \textbf{rangement complet} de tous les éléments de $E$. C'est un $n$-uplet où chaque élément apparaît exactement une fois.
\end{definition}

\begin{propriete}
Le nombre de permutations d'un ensemble à $n$ éléments est $n!$.
\end{propriete}

\begin{exemple}
De combien de façons peut-on ranger 6 livres sur une étagère ? $6! = 720$ façons.

\medskip

Combien d'anagrammes du mot MATHS ? Les 5 lettres sont distinctes : $5! = 120$ anagrammes.
\end{exemple}

% ── IV. Combinaisons ──
\section{Combinaisons — Coefficients binomiaux}

\subsection{Définition}

\begin{definition}[Combinaison]
Soit $E$ un ensemble à $n$ éléments et $k$ un entier tel que $0 \leqslant k \leqslant n$. Une \textbf{combinaison} de $k$ éléments de $E$ est une \textbf{partie} (sous-ensemble) de $E$ ayant exactement $k$ éléments.

Contrairement aux $k$-uplets, \textbf{l'ordre ne compte pas} : $\{a, b, c\} = \{c, a, b\}$.
\end{definition}

\begin{definition}[Coefficient binomial]
Le nombre de combinaisons de $k$ éléments parmi $n$ est le \textbf{coefficient binomial}, noté $\ds\binom{n}{k}$ (lire « $k$ parmi $n$ »).
\end{definition}

\begin{attention}[Ordre ou pas ?]
La question clé pour choisir entre arrangement et combinaison est : \textbf{l'ordre compte-t-il ?}
\begin{itemize}
  \item \textbf{L'ordre compte} $\rightarrow$ arrangement : $\ds\frac{n!}{(n-k)!}$
  \item \textbf{L'ordre ne compte pas} $\rightarrow$ combinaison : $\ds\binom{n}{k}$
\end{itemize}
\end{attention}

\subsection{Formule}

\begin{propriete}[Formule des coefficients binomiaux]
Pour $0 \leqslant k \leqslant n$ :
\[ \binom{n}{k} = \frac{n!}{k!\,(n-k)!} = \frac{n(n-1)\cdots(n-k+1)}{k!} \]
\end{propriete}

\begin{methode}[Comprendre la formule]
Le nombre de façons de choisir $k$ éléments \textbf{ordonnés} parmi $n$ est $\frac{n!}{(n-k)!}$. Mais dans une combinaison, l'ordre ne compte pas. Chaque combinaison de $k$ éléments donne $k!$ arrangements (permutations de ces $k$ éléments). On divise donc par $k!$.
\end{methode}

\begin{exemple}[Calculs]
\begin{itemize}
  \item $\ds\binom{5}{2} = \frac{5!}{2! \times 3!} = \frac{120}{2 \times 6} = 10$
  \item $\ds\binom{8}{3} = \frac{8 \times 7 \times 6}{3!} = \frac{336}{6} = 56$
  \item $\ds\binom{10}{4} = \frac{10 \times 9 \times 8 \times 7}{4!} = \frac{5040}{24} = 210$
\end{itemize}
\end{exemple}

\subsection{Valeurs particulières et symétrie}

\begin{propriete}[Cas particuliers]
Pour tout entier $n \geqslant 0$ :
\begin{itemize}
  \item $\ds\binom{n}{0} = 1$ \quad (une seule partie vide : $\emptyset$)
  \item $\ds\binom{n}{1} = n$ \quad ($n$ singletons)
  \item $\ds\binom{n}{n} = 1$ \quad (une seule partie contenant tous les éléments)
  \item $\ds\binom{n}{2} = \frac{n(n-1)}{2}$ \quad (nombre de paires)
\end{itemize}
\end{propriete}

\begin{propriete}[Symétrie]
Pour $0 \leqslant k \leqslant n$ :
\[ \binom{n}{k} = \binom{n}{n-k} \]
\end{propriete}

\textit{Interprétation :} choisir $k$ éléments \textbf{à prendre} revient à choisir $n-k$ éléments \textbf{à laisser}.

\begin{exemple}
$\ds\binom{10}{8} = \binom{10}{2} = \frac{10 \times 9}{2} = 45$. C'est bien plus rapide que de calculer directement.
\end{exemple}

% ── V. Triangle de Pascal ──
\section{Relation de Pascal et triangle de Pascal}

\subsection{Relation de Pascal}

\begin{propriete}[Relation de Pascal]
Pour tout $n \geqslant 1$ et $1 \leqslant k \leqslant n$ :
\[ \binom{n}{k} = \binom{n-1}{k-1} + \binom{n-1}{k} \]
\end{propriete}

\begin{methode}[Interprétation]
Pour former une partie de $k$ éléments de $\{1, 2, \ldots, n\}$, on distingue deux cas selon que l'élément $n$ est dans la partie ou non :
\begin{itemize}
  \item Si $n$ est dans la partie : il reste à choisir $k-1$ éléments parmi $\{1, \ldots, n-1\}$ $\rightarrow$ $\ds\binom{n-1}{k-1}$ choix
  \item Si $n$ n'est pas dans la partie : il faut choisir $k$ éléments parmi $\{1, \ldots, n-1\}$ $\rightarrow$ $\ds\binom{n-1}{k}$ choix
\end{itemize}
\end{methode}

\begin{demonstration}[de la relation de Pascal par le calcul]
\begin{align*}
\binom{n-1}{k-1} + \binom{n-1}{k} &= \frac{(n-1)!}{(k-1)!(n-k)!} + \frac{(n-1)!}{k!(n-1-k)!} \\
&= \frac{(n-1)!\,k}{k!(n-k)!} + \frac{(n-1)!\,(n-k)}{k!(n-k)!} \\
&= \frac{(n-1)!\,[k + (n-k)]}{k!(n-k)!} = \frac{n!}{k!(n-k)!} = \binom{n}{k} \qedhere
\end{align*}
\end{demonstration}

\subsection{Triangle de Pascal}

On construit le \textbf{triangle de Pascal} en plaçant les coefficients binomiaux ligne par ligne. Chaque nombre est la somme des deux nombres situés au-dessus de lui.

\begin{center}
\begin{tikzpicture}[scale=0.9]
  \def\rows{7}
  \foreach \n in {0,...,\rows} {
    \foreach \k in {0,...,\n} {
      \pgfmathtruncatemacro{\coeff}{factorial(\n)/(factorial(\k)*factorial(\n-\k))}
      \node[fill=indigobg,rounded corners=3pt,minimum size=7mm,inner sep=1pt,font=\small\bfseries]
        at (\k-\n/2, -\n*0.9) {\coeff};
    }
    \node[gray,font=\small\itshape] at (-\n/2-1.2, -\n*0.9) {$n=\n$};
  }
\end{tikzpicture}
\captionof{figure}{Triangle de Pascal (lignes $n=0$ à $n=7$)}
\end{center}

\subsection{Somme d'une ligne}

\begin{propriete}
Pour tout entier $n \geqslant 0$ :
\[ \sum_{k=0}^{n} \binom{n}{k} = 2^n \]
\end{propriete}

\begin{demonstration}[par dénombrement]
On compte l'ensemble $\mathcal{P}(E)$ de toutes les parties d'un ensemble $E$ à $n$ éléments de deux façons :
\begin{itemize}
  \item D'une part, $|\mathcal{P}(E)| = 2^n$ (chaque élément est pris ou non).
  \item D'autre part, il y a $\ds\binom{n}{k}$ parties de taille $k$, pour $k$ de 0 à $n$. Le total est $\ds\sum_{k=0}^{n}\binom{n}{k}$.
\end{itemize}
En identifiant : $\ds\sum_{k=0}^{n}\binom{n}{k} = 2^n$. \qed
\end{demonstration}

% ── VI. Applications ──
\section{Applications}

\begin{methode}[Choisir le bon outil]
Face à un problème de dénombrement, se poser les questions suivantes :
\begin{enumerate}
  \item \textbf{Y a-t-il répétition possible ?}
  \begin{itemize}
    \item Oui $\rightarrow$ $k$-uplets avec répétition : $n^k$
    \item Non $\rightarrow$ éléments distincts (questions suivantes)
  \end{itemize}
  \item \textbf{L'ordre compte-t-il ?}
  \begin{itemize}
    \item Oui $\rightarrow$ arrangement : $\ds\frac{n!}{(n-k)!}$
    \item Non $\rightarrow$ combinaison : $\ds\binom{n}{k}$
  \end{itemize}
  \item \textbf{Prend-on tous les éléments ?}
  \begin{itemize}
    \item Oui $\rightarrow$ permutation : $n!$
  \end{itemize}
\end{enumerate}
\end{methode}

\begin{exemple}[Le Loto]
Au Loto, on tire 5 numéros parmi 49 (sans remise, l'ordre ne compte pas).

Nombre de grilles : $\ds\binom{49}{5} = \frac{49 \times 48 \times 47 \times 46 \times 45}{120} = 1\,906\,884$.

La probabilité de trouver les 5 bons numéros est $\ds\frac{1}{1\,906\,884} \approx 5{,}24 \times 10^{-7}$, soit environ 1 chance sur 1,9 million.
\end{exemple}

\begin{exemple}[Mains de poker]
Un jeu de 52 cartes, 4 couleurs, 13 valeurs. On distribue 5 cartes.

\textbf{Nombre de mains :} $\ds\binom{52}{5} = 2\,598\,960$.

\textbf{Mains avec exactement 2 as :} on choisit 2 as parmi 4 ($\ds\binom{4}{2} = 6$) puis 3 cartes parmi les 48 non-as ($\ds\binom{48}{3} = 17\,296$). Total : $6 \times 17\,296 = 103\,776$.

\textbf{Couleur} (5 cartes de même couleur) : on choisit 1 couleur parmi 4, puis 5 cartes parmi 13. Total : $4 \times \ds\binom{13}{5} = 4 \times 1287 = 5148$.
\end{exemple}

% ── Bilan ──
\section*{Bilan du chapitre}

\begin{retenir}
\begin{itemize}[leftmargin=*]
  \item \textbf{Principe additif :} $\card(A \cup B) = \card(A) + \card(B)$ si $A \cap B = \emptyset$
  \item \textbf{Principe multiplicatif :} $\card(A \times B) = \card(A) \times \card(B)$
  \item \textbf{$k$-uplets avec répétition :} $n^k$
  \item \textbf{Factorielle :} $n! = n \times (n-1) \times \cdots \times 1$, \quad $0! = 1$
  \item \textbf{Arrangements} (ordre, sans répétition) : $\ds\frac{n!}{(n-k)!}$
  \item \textbf{Permutations :} $n!$
  \item \textbf{Combinaisons} (pas d'ordre) : $\ds\binom{n}{k} = \frac{n!}{k!(n-k)!}$
  \item \textbf{Symétrie :} $\ds\binom{n}{k} = \binom{n}{n-k}$
  \item \textbf{Pascal :} $\ds\binom{n}{k} = \binom{n-1}{k-1} + \binom{n-1}{k}$
  \item \textbf{Somme :} $\ds\sum_{k=0}^{n}\binom{n}{k} = 2^n$
\end{itemize}
\end{retenir}


% ─────────────────────────────────────────────────────────────────────────────
% CHAPITRES 2 À 15 — À COMPLÉTER
% ─────────────────────────────────────────────────────────────────────────────

% Les chapitres suivants seront ajoutés selon la même structure.
% Placeholder pour la compilation :

\chapter{Vecteurs, droites et plans de l'espace}

\section*{Introduction}

\textit{Comment décrire la position d'un satellite dans l'espace ? Comment un architecte modélise-t-il un bâtiment en trois dimensions ? Comment un physicien analyse-t-il les forces s'exerçant sur un objet en mouvement ?}

\medskip

En Seconde et en Première, nous avons étudié la géométrie du plan : vecteurs, droites, coordonnées dans un repère orthonormé. Ce chapitre étend ces outils à \textbf{l'espace à trois dimensions}. La géométrie dans l'espace est fondamentale en architecture, en physique, en ingénierie et en informatique graphique (jeux vidéo, modélisation 3D).

\medskip

\textbf{Objectifs du chapitre :}
\begin{itemize}[leftmargin=*]
  \item Manipuler les vecteurs de l'espace et la relation de Chasles
  \item Caractériser colinéarité et coplanarité
  \item Décrire droites et plans par leurs vecteurs directeurs
  \item Étudier les positions relatives de droites et plans
  \item Utiliser des repères et des coordonnées dans l'espace
\end{itemize}

\begin{histoire}[Des vecteurs dans l'espace]
La notion de vecteur dans l'espace s'est construite progressivement. Au XVII\ieme{} siècle, \textbf{Newton} utilise des grandeurs dirigées pour décrire forces et vitesses, sans les formaliser. Au XIX\ieme{} siècle, \textbf{Hermann Grassmann} (1809--1877) développe dans son \textit{Ausdehnungslehre} (1844) une théorie générale des espaces vectoriels de dimension quelconque, bien en avance sur son époque. Parallèlement, \textbf{William Rowan Hamilton} (1805--1865) invente les \textbf{quaternions}, système de nombres étendant les complexes à quatre dimensions, pour décrire les rotations dans l'espace. Ces travaux convergent vers la théorie moderne des espaces vectoriels, formalisée par Giuseppe Peano en 1888.
\end{histoire}

% ── I. Vecteurs de l'espace ──
\section{Vecteurs de l'espace}

\subsection{Rappels et extension à l'espace}

\begin{definition}[Vecteur de l'espace]
Un \textbf{vecteur de l'espace} est défini, comme dans le plan, par une direction, un sens et une norme. Deux vecteurs $\vect{AB}$ et $\vect{CD}$ sont \textbf{égaux} si et seulement si le quadrilatère $ABDC$ est un parallélogramme (éventuellement aplati).

Le vecteur nul est noté $\vec{0}$.
\end{definition}

\begin{propriete}[Opérations sur les vecteurs]
Les opérations sur les vecteurs du plan se généralisent à l'espace :
\begin{itemize}
  \item \textbf{Addition :} $\vec{u} + \vec{v}$ (règle du parallélogramme)
  \item \textbf{Multiplication par un réel :} $\lambda \vec{u}$ pour $\lambda \in \R$
  \item \textbf{Opposé :} $-\vec{u}$ (même direction, sens contraire)
\end{itemize}
\end{propriete}

\subsection{Relation de Chasles}

\begin{propriete}[Relation de Chasles]
Pour tous points $A$, $B$, $C$ de l'espace :
\[ \vect{AB} + \vect{BC} = \vect{AC} \]
\end{propriete}

\begin{exemple}
Soient $A$, $B$, $C$, $D$ quatre points de l'espace. On peut simplifier :
\[ \vect{AB} + \vect{BC} + \vect{CD} = \vect{AD} \]
En effet, par la relation de Chasles appliquée deux fois : $\vect{AB} + \vect{BC} = \vect{AC}$ puis $\vect{AC} + \vect{CD} = \vect{AD}$.
\end{exemple}

\subsection{Translation}

\begin{definition}[Translation]
La \textbf{translation} de vecteur $\vec{u}$ est la transformation qui, à tout point $M$ de l'espace, associe l'unique point $M'$ tel que $\vect{MM'} = \vec{u}$.
\end{definition}

% ── II. Combinaisons linéaires ──
\section{Combinaisons linéaires}

\subsection{Colinéarité}

\begin{definition}[Vecteurs colinéaires]
Deux vecteurs $\vec{u}$ et $\vec{v}$ sont \textbf{colinéaires} s'il existe un réel $\lambda$ tel que $\vec{v} = \lambda \vec{u}$, ou si $\vec{u} = \vec{0}$.

Deux vecteurs colinéaires ont la \textbf{même direction}.
\end{definition}

\begin{propriete}
Trois points $A$, $B$, $C$ sont \textbf{alignés} si et seulement si les vecteurs $\vect{AB}$ et $\vect{AC}$ sont colinéaires.
\end{propriete}

\subsection{Coplanarité}

\begin{definition}[Vecteurs coplanaires]
Trois vecteurs $\vec{u}$, $\vec{v}$, $\vec{w}$ sont \textbf{coplanaires} s'il existe des réels $a$, $b$, $c$ non tous nuls tels que :
\[ a\vec{u} + b\vec{v} + c\vec{w} = \vec{0} \]
De manière équivalente, $\vec{u}$, $\vec{v}$, $\vec{w}$ sont coplanaires si et seulement si l'un d'eux est \textbf{combinaison linéaire} des deux autres.
\end{definition}

\begin{propriete}[Caractérisation de la coplanarité]
Soient $\vec{u}$ et $\vec{v}$ deux vecteurs non colinéaires. Un vecteur $\vec{w}$ est coplanaire avec $\vec{u}$ et $\vec{v}$ si et seulement s'il existe des réels $a$ et $b$ tels que :
\[ \vec{w} = a\vec{u} + b\vec{v} \]
\end{propriete}

\begin{attention}[Coplanaire ne veut pas dire colinéaire]
Trois vecteurs coplanaires ne sont pas nécessairement colinéaires. Par exemple, $\vec{i}$ et $\vec{j}$ du plan ne sont pas colinéaires mais sont coplanaires (ils sont dans le même plan). En revanche, si $\vec{u}$, $\vec{v}$, $\vec{w}$ ne sont \textbf{pas} coplanaires, alors aucun des trois ne peut s'exprimer en fonction des deux autres : ils « engendrent » tout l'espace.
\end{attention}

\begin{center}
\begin{tikzpicture}[scale=1.2,>=Stealth]
  % Vecteurs coplanaires
  \begin{scope}[shift={(-4,0)}]
    \fill[indigobg,opacity=0.5] (-0.5,-0.3) -- (3,-0.3) -- (3.5,2.2) -- (0,2.2) -- cycle;
    \draw[->,thick,indigo] (0,0) -- (2,0) node[below] {$\vec{u}$};
    \draw[->,thick,vertdef] (0,0) -- (0.7,1.5) node[left] {$\vec{v}$};
    \draw[->,thick,rougeerr] (0,0) -- (2.7,1.5) node[right] {$\vec{w}=\vec{u}+\vec{v}$};
    \node[below,font=\small\itshape] at (1.5,-0.7) {Vecteurs coplanaires};
  \end{scope}
  % Vecteurs non coplanaires
  \begin{scope}[shift={(3,0)}]
    \draw[->,thick,indigo] (0,0) -- (2,0) node[below] {$\vec{u}$};
    \draw[->,thick,vertdef] (0,0) -- (0.7,1.5) node[left] {$\vec{v}$};
    \draw[->,thick,orangemeth] (0,0) -- (0.3,0.8) node[above right] {$\vec{w}$};
    \draw[dashed,gray] (0.3,0.8) -- (0.3,0) -- (0,0);
    \draw[dashed,gray] (0.3,0.8) -- (1,0.8);
    \node[below,font=\small\itshape] at (1.2,-0.7) {Vecteurs non coplanaires};
    \node[font=\small\itshape] at (1.2,-1.1) {($\vec{w}$ sort du plan)};
  \end{scope}
\end{tikzpicture}
\end{center}

\begin{methode}[Montrer que des vecteurs sont coplanaires]
Pour montrer que $\vec{w}$ est coplanaire avec $\vec{u}$ et $\vec{v}$ (non colinéaires), on cherche $a$ et $b$ tels que $\vec{w} = a\vec{u} + b\vec{v}$.

\medskip

\textbf{Étapes :}
\begin{enumerate}
  \item Exprimer $\vec{w}$ en fonction de $\vec{u}$ et $\vec{v}$ (si coordonnées : résoudre un système).
  \item Si le système a une solution : $\vec{w}$ est coplanaire avec $\vec{u}$ et $\vec{v}$.
  \item Si le système n'a pas de solution : $\vec{w}$ n'est pas coplanaire avec $\vec{u}$ et $\vec{v}$.
\end{enumerate}
\end{methode}

% ── III. Droites de l'espace ──
\section{Droites de l'espace}

\begin{definition}[Vecteur directeur d'une droite]
Un vecteur non nul $\vec{u}$ est un \textbf{vecteur directeur} d'une droite $d$ si $\vec{u}$ est colinéaire à $\vect{AB}$ pour tous points $A$, $B$ distincts de $d$.

Une droite est entièrement déterminée par un point et un vecteur directeur.
\end{definition}

\begin{definition}[Droites parallèles]
Deux droites sont \textbf{parallèles} si et seulement si leurs vecteurs directeurs sont \textbf{colinéaires}. On inclut le cas où les droites sont confondues.
\end{definition}

\begin{propriete}[Caractérisation d'une droite]
Soit $d$ une droite passant par $A$ de vecteur directeur $\vec{u}$. Un point $M$ appartient à $d$ si et seulement s'il existe un réel $t$ tel que :
\[ \vect{AM} = t\,\vec{u} \]
\end{propriete}

% ── IV. Plans de l'espace ──
\section{Plans de l'espace}

\subsection{Vecteurs directeurs d'un plan}

\begin{definition}[Couple de vecteurs directeurs d'un plan]
Deux vecteurs non colinéaires $\vec{u}$ et $\vec{v}$ sont des \textbf{vecteurs directeurs} d'un plan $\mathcal{P}$ si $\vec{u}$ et $\vec{v}$ sont parallèles à $\mathcal{P}$ (c'est-à-dire colinéaires à des vecteurs tracés dans $\mathcal{P}$).
\end{definition}

\subsection{Caractérisations d'un plan}

\begin{propriete}[Différentes façons de définir un plan]
Un plan de l'espace est déterminé par :
\begin{itemize}
  \item \textbf{Trois points non alignés} $A$, $B$, $C$
  \item \textbf{Un point et deux vecteurs non colinéaires} : $A$, $\vec{u}$, $\vec{v}$
  \item \textbf{Un point et une droite ne passant pas par ce point}
  \item \textbf{Deux droites sécantes}
  \item \textbf{Deux droites strictement parallèles}
\end{itemize}
\end{propriete}

\begin{propriete}[Appartenance à un plan]
Soit $\mathcal{P}$ le plan passant par $A$ et de vecteurs directeurs $\vec{u}$ et $\vec{v}$ (non colinéaires). Un point $M$ appartient à $\mathcal{P}$ si et seulement s'il existe des réels $s$ et $t$ tels que :
\[ \vect{AM} = s\,\vec{u} + t\,\vec{v} \]
\end{propriete}

\begin{definition}[Plans parallèles]
Deux plans sont \textbf{parallèles} si et seulement si tout couple de vecteurs directeurs de l'un est coplanaire avec les vecteurs directeurs de l'autre. En pratique, les vecteurs directeurs de l'un s'expriment comme combinaisons linéaires de ceux de l'autre.
\end{definition}

% ── V. Positions relatives ──
\section{Positions relatives}

\subsection{Positions relatives de deux droites}

\begin{propriete}[Quatre cas possibles]
Deux droites de l'espace sont dans l'une des situations suivantes :

\medskip

\begin{center}
\renewcommand{\arraystretch}{1.4}
\begin{tabular}{|l|c|c|}
\hline
\textbf{Position} & \textbf{Point commun ?} & \textbf{Coplanaires ?} \\
\hline
Confondues & $\infty$ & Oui \\
Strictement parallèles & Non & Oui \\
Sécantes & 1 & Oui \\
Non coplanaires & Non & Non \\
\hline
\end{tabular}
\end{center}
\end{propriete}

\begin{attention}[Le cas « non coplanaires »]
Contrairement au plan, deux droites de l'espace peuvent ne pas se couper \textbf{sans être parallèles}. C'est le cas des droites \textbf{non coplanaires} (aussi appelées « gauches »). Par exemple, dans un cube, les droites $(AB)$ et $(DH)$ ne se coupent pas et ne sont pas parallèles.
\end{attention}

\subsection{Positions relatives d'une droite et d'un plan}

\begin{propriete}[Trois cas possibles]
\begin{center}
\renewcommand{\arraystretch}{1.4}
\begin{tabular}{|l|c|}
\hline
\textbf{Position} & \textbf{Points communs} \\
\hline
La droite est incluse dans le plan & Tous (la droite entière) \\
La droite est strictement parallèle au plan & Aucun \\
La droite est sécante au plan & Exactement 1 \\
\hline
\end{tabular}
\end{center}
\end{propriete}

\begin{propriete}[Parallélisme droite-plan]
Une droite $d$ de vecteur directeur $\vec{u}$ est parallèle à un plan $\mathcal{P}$ de vecteurs directeurs $\vec{v}_1$ et $\vec{v}_2$ si et seulement si $\vec{u}$, $\vec{v}_1$, $\vec{v}_2$ sont \textbf{coplanaires} (c'est-à-dire $\vec{u} = a\vec{v}_1 + b\vec{v}_2$ pour certains réels $a$, $b$).
\end{propriete}

\subsection{Positions relatives de deux plans}

\begin{propriete}[Trois cas possibles]
\begin{center}
\renewcommand{\arraystretch}{1.4}
\begin{tabular}{|l|c|}
\hline
\textbf{Position} & \textbf{Intersection} \\
\hline
Confondus & Le plan entier \\
Strictement parallèles & Vide \\
Sécants & Une droite \\
\hline
\end{tabular}
\end{center}
\end{propriete}

\begin{propriete}[Théorème du toit]
Si deux plans $\mathcal{P}_1$ et $\mathcal{P}_2$ sont sécants selon une droite $\Delta$, et si une droite $d_1 \subset \mathcal{P}_1$ et une droite $d_2 \subset \mathcal{P}_2$ sont toutes deux parallèles à $\Delta$, alors $d_1 \parallel d_2$.
\end{propriete}

% ── VI. Bases et repères ──
\section{Bases et repères de l'espace}

\begin{definition}[Base de l'espace]
Trois vecteurs $\vec{u}$, $\vec{v}$, $\vec{w}$ forment une \textbf{base} de l'espace s'ils sont \textbf{non coplanaires}, c'est-à-dire si aucun des trois ne s'écrit comme combinaison linéaire des deux autres.
\end{definition}

\begin{propriete}[Décomposition unique]
Si $(\vec{u}, \vec{v}, \vec{w})$ est une base de l'espace, alors tout vecteur $\vec{t}$ de l'espace s'écrit de manière \textbf{unique} :
\[ \vec{t} = a\vec{u} + b\vec{v} + c\vec{w} \qquad \text{avec } a, b, c \in \R \]
Les réels $a$, $b$, $c$ sont les \textbf{composantes} (ou coordonnées) de $\vec{t}$ dans la base $(\vec{u}, \vec{v}, \vec{w})$.
\end{propriete}

\begin{demonstration}[Unicité de la décomposition]
Supposons que $\vec{t} = a\vec{u} + b\vec{v} + c\vec{w}$ et $\vec{t} = a'\vec{u} + b'\vec{v} + c'\vec{w}$. Par soustraction :
\[ (a-a')\vec{u} + (b-b')\vec{v} + (c-c')\vec{w} = \vec{0} \]
Comme $\vec{u}$, $\vec{v}$, $\vec{w}$ sont non coplanaires, la seule combinaison linéaire donnant $\vec{0}$ est la combinaison triviale : $a - a' = 0$, $b - b' = 0$, $c - c' = 0$. Donc $a = a'$, $b = b'$, $c = c'$. \qed
\end{demonstration}

\begin{definition}[Repère de l'espace]
Un \textbf{repère} de l'espace est la donnée d'un point $O$ (origine) et d'une base $(\vec{i}, \vec{j}, \vec{k})$. On le note $(O\,;\,\vec{i}, \vec{j}, \vec{k})$.

Les \textbf{coordonnées} d'un point $M$ dans ce repère sont les composantes $(x, y, z)$ du vecteur $\vect{OM}$ :
\[ \vect{OM} = x\vec{i} + y\vec{j} + z\vec{k} \]
\end{definition}

\begin{center}
\begin{tikzpicture}[scale=1.5,>=Stealth]
  \coordinate (O) at (0,0);
  \draw[->,thick,indigo] (O) -- (2,0) node[below right] {$\vec{i}$};
  \draw[->,thick,vertdef] (O) -- (0,2) node[above left] {$\vec{k}$};
  \draw[->,thick,orangemeth] (O) -- (1.2,-0.6) node[below] {$\vec{j}$};
  \node[below left] at (O) {$O$};
  % Point M
  \coordinate (M) at (1.6,1.5);
  \fill[rougeerr] (M) circle (1.5pt) node[above right] {$M(x,y,z)$};
  \draw[dashed,gray] (M) -- (1.6,0) -- (0,0);
  \draw[dashed,gray] (1.6,0) -- (1.6+0.48,-0.24);
\end{tikzpicture}
\end{center}

\begin{propriete}[Coordonnées du milieu]
Si $A(x_A, y_A, z_A)$ et $B(x_B, y_B, z_B)$, le milieu $I$ de $[AB]$ a pour coordonnées :
\[ I\left(\frac{x_A + x_B}{2}\,;\,\frac{y_A + y_B}{2}\,;\,\frac{z_A + z_B}{2}\right) \]
\end{propriete}

\begin{propriete}[Coordonnées d'un vecteur]
Si $A(x_A, y_A, z_A)$ et $B(x_B, y_B, z_B)$, alors :
\[ \vect{AB} \begin{pmatrix} x_B - x_A \\ y_B - y_A \\ z_B - z_A \end{pmatrix} \]
\end{propriete}

% ── VII. Applications ──
\section{Applications}

\begin{exemple}[Travail dans un cube]
On considère le cube $ABCDEFGH$ où $ABCD$ est la face inférieure et $EFGH$ la face supérieure, avec $E$ au-dessus de $A$, $F$ au-dessus de $B$, etc.

On munit l'espace du repère $(A\,;\,\vect{AB}, \vect{AD}, \vect{AE})$.

\medskip

\begin{center}
\begin{tikzpicture}[scale=2,>=Stealth]
  % Face inférieure
  \coordinate (A) at (0,0);
  \coordinate (B) at (2,0);
  \coordinate (C) at (2.8,0.6);
  \coordinate (D) at (0.8,0.6);
  % Face supérieure
  \coordinate (E) at (0,1.8);
  \coordinate (F) at (2,1.8);
  \coordinate (G) at (2.8,2.4);
  \coordinate (H) at (0.8,2.4);
  % Arêtes
  \draw[thick] (A) -- (B) -- (C) -- (D) -- cycle;
  \draw[thick] (E) -- (F) -- (G) -- (H) -- cycle;
  \draw[thick] (A) -- (E);
  \draw[thick] (B) -- (F);
  \draw[thick] (C) -- (G);
  \draw[dashed] (D) -- (H);
  % Labels
  \node[below left] at (A) {$A$};
  \node[below right] at (B) {$B$};
  \node[right] at (C) {$C$};
  \node[above left] at (D) {$D$};
  \node[left] at (E) {$E$};
  \node[right] at (F) {$F$};
  \node[above right] at (G) {$G$};
  \node[above] at (H) {$H$};
  % Vecteurs repère
  \draw[->,very thick,indigo] (A) -- (1,0) node[below,font=\small] {$\vec{i}$};
  \draw[->,very thick,vertdef] (A) -- (0.4,0.3) node[below right,font=\small] {$\vec{j}$};
  \draw[->,very thick,orangemeth] (A) -- (0,0.9) node[left,font=\small] {$\vec{k}$};
\end{tikzpicture}
\end{center}

\textbf{Coordonnées des sommets :}
\[ A(0,0,0) \quad B(1,0,0) \quad C(1,1,0) \quad D(0,1,0) \]
\[ E(0,0,1) \quad F(1,0,1) \quad G(1,1,1) \quad H(0,1,1) \]

\textbf{1.} Calculons $\vect{AG}$ :
\[ \vect{AG} = \begin{pmatrix} 1-0 \\ 1-0 \\ 1-0 \end{pmatrix} = \begin{pmatrix} 1 \\ 1 \\ 1 \end{pmatrix} = \vec{i} + \vec{j} + \vec{k} \]

\textbf{2.} Le milieu $I$ de $[AG]$ a pour coordonnées :
\[ I\left(\frac{0+1}{2}\,;\,\frac{0+1}{2}\,;\,\frac{0+1}{2}\right) = I\left(\frac{1}{2}, \frac{1}{2}, \frac{1}{2}\right) \]
C'est le centre du cube.

\textbf{3.} Les vecteurs $\vect{AB}$ et $\vect{EG}$ sont-ils colinéaires ?
\[ \vect{AB} = \begin{pmatrix} 1 \\ 0 \\ 0 \end{pmatrix} \qquad \vect{EG} = \begin{pmatrix} 1 \\ 1 \\ 0 \end{pmatrix} \]
Il n'existe pas de réel $\lambda$ tel que $(1,1,0) = \lambda(1,0,0)$. Donc $\vect{AB}$ et $\vect{EG}$ ne sont pas colinéaires : les droites $(AB)$ et $(EG)$ ne sont pas parallèles.
\end{exemple}

\begin{methode}[Montrer que trois points sont alignés dans l'espace]
Pour montrer que $A$, $B$, $C$ sont alignés :
\begin{enumerate}
  \item Calculer les vecteurs $\vect{AB}$ et $\vect{AC}$.
  \item Chercher un réel $\lambda$ tel que $\vect{AC} = \lambda\,\vect{AB}$.
  \item Si un tel $\lambda$ existe, les vecteurs sont colinéaires et les points sont alignés.
\end{enumerate}
\end{methode}

\begin{exemple}[Alignement dans le cube]
Dans le cube $ABCDEFGH$ précédent, le milieu $M$ de $[EG]$ a pour coordonnées $M\left(\frac{1}{2}, \frac{1}{2}, 1\right)$.

Montrons que $D$, $I$ (centre du cube) et $F$ sont alignés.
\[ \vect{DF} = \begin{pmatrix} 1 \\ -1 \\ 1 \end{pmatrix} \qquad \vect{DI} = \begin{pmatrix} \frac{1}{2} \\ -\frac{1}{2} \\ \frac{1}{2} \end{pmatrix} = \frac{1}{2}\vect{DF} \]
Les vecteurs sont colinéaires, donc $D$, $I$, $F$ sont alignés, et $I$ est le milieu de $[DF]$.
\end{exemple}

\begin{exemple}[Coplanarité dans le cube]
Dans le même cube, les vecteurs $\vect{AB}$, $\vect{AD}$, $\vect{AE}$ sont-ils coplanaires ?

Si $a\,\vect{AB} + b\,\vect{AD} + c\,\vect{AE} = \vec{0}$, alors :
\[ a\begin{pmatrix}1\\0\\0\end{pmatrix} + b\begin{pmatrix}0\\1\\0\end{pmatrix} + c\begin{pmatrix}0\\0\\1\end{pmatrix} = \begin{pmatrix}0\\0\\0\end{pmatrix} \]
Ce qui donne $a=0$, $b=0$, $c=0$. La seule solution est la triviale : les vecteurs sont \textbf{non coplanaires}. Ils forment donc une base de l'espace.
\end{exemple}

% ── Bilan ──
\section*{Bilan du chapitre}

\begin{retenir}
\begin{itemize}[leftmargin=*]
  \item \textbf{Vecteurs de l'espace :} mêmes propriétés que dans le plan (addition, multiplication par un réel, relation de Chasles)
  \item \textbf{Colinéarité :} $\vec{v} = \lambda\vec{u}$ — même direction
  \item \textbf{Coplanarité :} trois vecteurs $\vec{u}$, $\vec{v}$, $\vec{w}$ sont coplanaires si l'un est combinaison linéaire des deux autres
  \item \textbf{Droite :} définie par un point $A$ et un vecteur directeur $\vec{u}$ ; $M \in d \Leftrightarrow \vect{AM} = t\vec{u}$
  \item \textbf{Plan :} défini par un point $A$ et deux vecteurs directeurs non colinéaires $\vec{u}$, $\vec{v}$ ; $M \in \mathcal{P} \Leftrightarrow \vect{AM} = s\vec{u} + t\vec{v}$
  \item \textbf{Deux droites :} confondues, parallèles, sécantes, ou non coplanaires
  \item \textbf{Droite et plan :} incluse, parallèle, ou sécante
  \item \textbf{Deux plans :} confondus, parallèles, ou sécants (intersection = droite)
  \item \textbf{Base :} trois vecteurs non coplanaires ; tout vecteur se décompose de façon unique
  \item \textbf{Repère :} $(O\,;\,\vec{i}, \vec{j}, \vec{k})$ ; coordonnées $(x,y,z)$ avec $\vect{OM} = x\vec{i}+y\vec{j}+z\vec{k}$
  \item \textbf{Milieu :} $I\left(\frac{x_A+x_B}{2}, \frac{y_A+y_B}{2}, \frac{z_A+z_B}{2}\right)$
\end{itemize}
\end{retenir}

\chapter{Orthogonalité et distances dans l'espace}

\section*{Introduction}

\textit{Comment mesurer la distance entre un satellite et la surface de la Terre ? Comment calculer l'angle entre deux murs d'un bâtiment ? Comment déterminer la hauteur d'une pyramide sans la mesurer directement ?}

\medskip

Après avoir décrit les objets de l'espace (vecteurs, droites, plans) au chapitre précédent, nous disposons maintenant des outils pour \textbf{mesurer} : distances, angles, perpendicularité. L'outil central est le \textbf{produit scalaire}, qui étend à l'espace le produit scalaire du plan étudié en Première.

\medskip

\textbf{Objectifs du chapitre :}
\begin{itemize}[leftmargin=*]
  \item Calculer le produit scalaire de deux vecteurs de l'espace
  \item Déterminer des normes et des distances
  \item Caractériser l'orthogonalité de vecteurs, droites et plans
  \item Calculer le projeté orthogonal d'un point sur un plan ou une droite
  \item Utiliser la formule de distance d'un point à un plan
\end{itemize}

\begin{histoire}[Du produit scalaire]
Le produit scalaire apparaît implicitement dès les travaux de \textbf{Lagrange} (1736--1813) sur la mécanique analytique. C'est \textbf{Hamilton} qui, dans sa théorie des quaternions (1843), distingue pour la première fois la partie « scalaire » et la partie « vectorielle » du produit de deux vecteurs. \textbf{Josiah Willard Gibbs} (1839--1903), physicien américain, isole le produit scalaire (\textit{dot product}) et le produit vectoriel (\textit{cross product}) dans ses \textit{Elements of Vector Analysis} (1881), créant ainsi le formalisme vectoriel utilisé aujourd'hui en physique et en ingénierie.
\end{histoire}

% ── I. Produit scalaire dans l'espace ──
\section{Produit scalaire dans l'espace}

\subsection{Définition}

\begin{definition}[Produit scalaire]
Soient $\vec{u}$ et $\vec{v}$ deux vecteurs de l'espace. Le \textbf{produit scalaire} de $\vec{u}$ et $\vec{v}$, noté $\vec{u} \cdot \vec{v}$, est le nombre réel défini par :
\begin{itemize}
  \item Si $\vec{u} = \vec{0}$ ou $\vec{v} = \vec{0}$ : $\vec{u} \cdot \vec{v} = 0$
  \item Sinon : $\vec{u} \cdot \vec{v} = \norme{\vec{u}} \times \norme{\vec{v}} \times \cos\theta$
\end{itemize}
où $\theta$ est l'angle entre les vecteurs $\vec{u}$ et $\vec{v}$ (angle $(\vec{u}, \vec{v})$ compris entre $0$ et $\pi$).
\end{definition}

\begin{propriete}[Expression dans une base orthonormée]
Si $\vec{u}\begin{pmatrix}x\\y\\z\end{pmatrix}$ et $\vec{v}\begin{pmatrix}x'\\y'\\z'\end{pmatrix}$ dans un repère \textbf{orthonormé}, alors :
\[ \vec{u} \cdot \vec{v} = xx' + yy' + zz' \]
\end{propriete}

\begin{attention}[Repère orthonormé obligatoire]
La formule $\vec{u} \cdot \vec{v} = xx' + yy' + zz'$ n'est valable que dans un repère \textbf{orthonormé}, c'est-à-dire un repère $(O\,;\,\vec{i},\vec{j},\vec{k})$ où :
\[ \norme{\vec{i}} = \norme{\vec{j}} = \norme{\vec{k}} = 1 \qquad \text{et} \qquad \vec{i}\cdot\vec{j} = \vec{j}\cdot\vec{k} = \vec{i}\cdot\vec{k} = 0 \]
Dans un repère quelconque (par exemple le repère du cube), cette formule est \textbf{fausse}.
\end{attention}

\subsection{Propriétés du produit scalaire}

\begin{propriete}[Propriétés algébriques]
Pour tous vecteurs $\vec{u}$, $\vec{v}$, $\vec{w}$ et tout réel $\lambda$ :
\begin{enumerate}
  \item \textbf{Symétrie :} $\vec{u} \cdot \vec{v} = \vec{v} \cdot \vec{u}$
  \item \textbf{Bilinéarité :}
  \begin{itemize}
    \item $\vec{u} \cdot (\vec{v} + \vec{w}) = \vec{u}\cdot\vec{v} + \vec{u}\cdot\vec{w}$
    \item $(\lambda\vec{u}) \cdot \vec{v} = \lambda(\vec{u}\cdot\vec{v})$
  \end{itemize}
  \item \textbf{Carré scalaire :} $\vec{u} \cdot \vec{u} = \norme{\vec{u}}^2 \geqslant 0$, avec égalité si et seulement si $\vec{u} = \vec{0}$
\end{enumerate}
\end{propriete}

\begin{exemple}[Calcul dans un repère orthonormé]
Soient $\vec{u}\begin{pmatrix}2\\-1\\3\end{pmatrix}$ et $\vec{v}\begin{pmatrix}1\\4\\-2\end{pmatrix}$ dans un repère orthonormé.

\[ \vec{u}\cdot\vec{v} = 2\times 1 + (-1)\times 4 + 3\times(-2) = 2 - 4 - 6 = -8 \]
\end{exemple}

% ── II. Norme et distance ──
\section{Norme et distance}

\begin{definition}[Norme d'un vecteur]
La \textbf{norme} (ou longueur) d'un vecteur $\vec{u}$ est le nombre réel positif :
\[ \norme{\vec{u}} = \sqrt{\vec{u}\cdot\vec{u}} \]
\end{definition}

\begin{propriete}[Norme dans un repère orthonormé]
Si $\vec{u}\begin{pmatrix}x\\y\\z\end{pmatrix}$ dans un repère orthonormé :
\[ \norme{\vec{u}} = \sqrt{x^2 + y^2 + z^2} \]
\end{propriete}

\begin{propriete}[Distance entre deux points]
Si $A(x_A, y_A, z_A)$ et $B(x_B, y_B, z_B)$ dans un repère orthonormé :
\[ AB = \norme{\vect{AB}} = \sqrt{(x_B - x_A)^2 + (y_B - y_A)^2 + (z_B - z_A)^2} \]
\end{propriete}

\begin{exemple}
Soient $A(1, 2, -1)$ et $B(3, 0, 2)$ dans un repère orthonormé.

\[ AB = \sqrt{(3-1)^2 + (0-2)^2 + (2-(-1))^2} = \sqrt{4 + 4 + 9} = \sqrt{17} \]
\end{exemple}

% ── III. Orthogonalité des vecteurs ──
\section{Orthogonalité}

\subsection{Vecteurs orthogonaux}

\begin{definition}[Vecteurs orthogonaux]
Deux vecteurs $\vec{u}$ et $\vec{v}$ sont \textbf{orthogonaux} si et seulement si $\vec{u} \cdot \vec{v} = 0$. On note $\vec{u} \perp \vec{v}$.

\medskip

Par convention, le vecteur nul est orthogonal à tout vecteur.
\end{definition}

\begin{propriete}[Caractérisation par les coordonnées]
Dans un repère orthonormé, $\vec{u}\begin{pmatrix}x\\y\\z\end{pmatrix}$ et $\vec{v}\begin{pmatrix}x'\\y'\\z'\end{pmatrix}$ sont orthogonaux si et seulement si :
\[ xx' + yy' + zz' = 0 \]
\end{propriete}

\subsection{Formules de polarisation}

\begin{propriete}[Formules de polarisation]
Pour tous vecteurs $\vec{u}$ et $\vec{v}$ :
\begin{align}
  \vec{u}\cdot\vec{v} &= \frac{1}{2}\left(\norme{\vec{u}+\vec{v}}^2 - \norme{\vec{u}}^2 - \norme{\vec{v}}^2\right) \label{eq:pol1}\\
  \vec{u}\cdot\vec{v} &= \frac{1}{2}\left(\norme{\vec{u}}^2 + \norme{\vec{v}}^2 - \norme{\vec{u}-\vec{v}}^2\right) \label{eq:pol2}\\
  \vec{u}\cdot\vec{v} &= \frac{1}{4}\left(\norme{\vec{u}+\vec{v}}^2 - \norme{\vec{u}-\vec{v}}^2\right) \label{eq:pol3}
\end{align}
\end{propriete}

\begin{demonstration}[de la formule \eqref{eq:pol1}]
Développons $\norme{\vec{u}+\vec{v}}^2$ en utilisant les propriétés du produit scalaire :
\begin{align*}
\norme{\vec{u}+\vec{v}}^2 &= (\vec{u}+\vec{v})\cdot(\vec{u}+\vec{v}) \\
&= \vec{u}\cdot\vec{u} + \vec{u}\cdot\vec{v} + \vec{v}\cdot\vec{u} + \vec{v}\cdot\vec{v} \\
&= \norme{\vec{u}}^2 + 2\,\vec{u}\cdot\vec{v} + \norme{\vec{v}}^2
\end{align*}
D'où :
\[ \vec{u}\cdot\vec{v} = \frac{1}{2}\left(\norme{\vec{u}+\vec{v}}^2 - \norme{\vec{u}}^2 - \norme{\vec{v}}^2\right) \qed \]
\end{demonstration}

\begin{methode}[Calculer un produit scalaire sans repère orthonormé]
Quand on ne dispose pas d'un repère orthonormé (par exemple dans un cube avec le repère $(A\,;\,\vect{AB},\vect{AD},\vect{AE})$), on utilise les formules de polarisation en calculant des normes (longueurs) dans la figure.

\medskip

\textbf{Étapes :}
\begin{enumerate}
  \item Décomposer les vecteurs en fonction de vecteurs dont on connaît les normes.
  \item Utiliser la formule $\vec{u}\cdot\vec{v} = \frac{1}{2}(\norme{\vec{u}+\vec{v}}^2 - \norme{\vec{u}}^2 - \norme{\vec{v}}^2)$.
  \item Ou bien utiliser la bilinéarité pour développer.
\end{enumerate}
\end{methode}

% ── IV. Base orthonormée ──
\section{Base orthonormée}

\begin{definition}[Base orthonormée]
Une base $(\vec{i}, \vec{j}, \vec{k})$ de l'espace est \textbf{orthonormée} (ou orthonormale) si :
\[ \vec{i}\cdot\vec{j} = 0, \qquad \vec{j}\cdot\vec{k} = 0, \qquad \vec{i}\cdot\vec{k} = 0 \]
et
\[ \norme{\vec{i}} = \norme{\vec{j}} = \norme{\vec{k}} = 1 \]
\end{definition}

\begin{attention}[Quand utiliser quelle formule ?]
\begin{itemize}
  \item \textbf{Repère orthonormé disponible :} utiliser $\vec{u}\cdot\vec{v} = xx'+yy'+zz'$ (le plus simple).
  \item \textbf{Pas de repère orthonormé :} utiliser la définition avec l'angle $\norme{\vec{u}}\norme{\vec{v}}\cos\theta$, les formules de polarisation, ou la bilinéarité.
\end{itemize}
\end{attention}

% ── V. Orthogonalité droites et plans ──
\section{Orthogonalité de droites et de plans}

\subsection{Droites orthogonales}

\begin{definition}[Droites orthogonales]
Deux droites de l'espace sont \textbf{orthogonales} si leurs vecteurs directeurs sont orthogonaux.

\medskip

\textbf{Remarque :} des droites orthogonales ne sont pas nécessairement sécantes (elles peuvent être non coplanaires).
\end{definition}

\begin{definition}[Droites perpendiculaires]
Deux droites sont \textbf{perpendiculaires} si elles sont orthogonales \textbf{et} sécantes.
\end{definition}

\subsection{Vecteur normal à un plan}

\begin{definition}[Vecteur normal]
Un vecteur $\vec{n}$ non nul est \textbf{normal} (ou orthogonal) à un plan $\mathcal{P}$ s'il est orthogonal à \textbf{tout} vecteur directeur de $\mathcal{P}$.
\end{definition}

\begin{propriete}
Un vecteur $\vec{n}$ non nul est normal à un plan $\mathcal{P}$ de vecteurs directeurs $\vec{u}$ et $\vec{v}$ si et seulement si :
\[ \vec{n}\cdot\vec{u} = 0 \quad \text{et} \quad \vec{n}\cdot\vec{v} = 0 \]
\end{propriete}

\begin{center}
\begin{tikzpicture}[scale=1.3,>=Stealth]
  % Plan
  \fill[indigobg,opacity=0.5] (-2,-0.5) -- (2,-0.5) -- (2.5,1.5) -- (-1.5,1.5) -- cycle;
  \node[indigo,font=\small] at (2.2,1.2) {$\mathcal{P}$};
  % Vecteurs directeurs
  \draw[->,thick,vertdef] (0,0.3) -- (1.5,0.3) node[below right] {$\vec{u}$};
  \draw[->,thick,vertdef] (0,0.3) -- (0.3,1.2) node[above left] {$\vec{v}$};
  % Vecteur normal
  \draw[->,very thick,rougeerr] (0,0.3) -- (0,2.3) node[right] {$\vec{n}$};
  % Point A
  \fill (0,0.3) circle (1.5pt) node[below left] {$A$};
  % Angle droit
  \draw[rougeerr] (0.15,0.3) -- (0.15,0.45) -- (0,0.45);
\end{tikzpicture}
\end{center}

\subsection{Droite orthogonale à un plan}

\begin{definition}[Droite orthogonale à un plan]
Une droite $d$ est \textbf{orthogonale} (ou perpendiculaire) à un plan $\mathcal{P}$ si son vecteur directeur est un vecteur normal à $\mathcal{P}$.
\end{definition}

\begin{propriete}[Caractérisation]
Une droite $d$ est orthogonale à un plan $\mathcal{P}$ si et seulement si $d$ est orthogonale à \textbf{deux droites sécantes} de $\mathcal{P}$.
\end{propriete}

\begin{propriete}[Plan passant par un point de normale donnée]
Soit $A$ un point et $\vec{n}$ un vecteur non nul. Il existe un unique plan $\mathcal{P}$ passant par $A$ et de vecteur normal $\vec{n}$. Un point $M$ appartient à $\mathcal{P}$ si et seulement si :
\[ \vect{AM} \cdot \vec{n} = 0 \]
\end{propriete}

\begin{demonstration}
Soit $\mathcal{P}$ le plan passant par $A$ et orthogonal à $\vec{n}$.

$(\Rightarrow)$ Si $M \in \mathcal{P}$, alors $\vect{AM}$ est un vecteur de $\mathcal{P}$, donc orthogonal à $\vec{n}$, d'où $\vect{AM}\cdot\vec{n} = 0$.

$(\Leftarrow)$ Si $\vect{AM}\cdot\vec{n} = 0$, alors $\vect{AM}$ est orthogonal à $\vec{n}$. Soient $\vec{u}$, $\vec{v}$ deux vecteurs directeurs de $\mathcal{P}$. Tout vecteur orthogonal à $\vec{n}$ est combinaison linéaire de $\vec{u}$ et $\vec{v}$ (car $\vec{n}$, $\vec{u}$, $\vec{v}$ forment une base). Donc $\vect{AM} = s\vec{u}+t\vec{v}$ et $M \in \mathcal{P}$. \qed
\end{demonstration}

% ── VI. Projeté orthogonal ──
\section{Projeté orthogonal}

\begin{definition}[Projeté orthogonal sur un plan]
Soit $\mathcal{P}$ un plan et $M$ un point de l'espace. Le \textbf{projeté orthogonal} de $M$ sur $\mathcal{P}$ est le point $H$ de $\mathcal{P}$ tel que la droite $(MH)$ est orthogonale à $\mathcal{P}$.

Le point $H$ est le point de $\mathcal{P}$ \textbf{le plus proche} de $M$.
\end{definition}

\begin{definition}[Projeté orthogonal sur une droite]
Soit $d$ une droite et $M$ un point de l'espace. Le \textbf{projeté orthogonal} de $M$ sur $d$ est le point $H$ de $d$ tel que $\vect{MH} \cdot \vec{u} = 0$, où $\vec{u}$ est un vecteur directeur de $d$.
\end{definition}

\begin{methode}[Calculer le projeté orthogonal sur une droite]
Soit $d$ passant par $A$ de vecteur directeur $\vec{u}$ et $M$ un point. Le projeté $H$ de $M$ sur $d$ vérifie $\vect{AH} = t\vec{u}$ pour un certain réel $t$.

\medskip

La condition $\vect{MH}\cdot\vec{u} = 0$ donne :
\[ (\vect{MA} + t\vec{u})\cdot\vec{u} = 0 \implies t = \frac{\vect{AM}\cdot\vec{u}}{\norme{\vec{u}}^2} \]

Alors $H = A + t\vec{u}$ et $MH = \norme{\vect{MH}}$.
\end{methode}

\begin{exemple}[Projeté dans un repère orthonormé]
Dans un repère orthonormé, soit la droite $d$ passant par $A(1,0,2)$ de vecteur directeur $\vec{u}\begin{pmatrix}1\\2\\-1\end{pmatrix}$ et le point $M(3,1,0)$.

Calculons $t$ :
\[ \vect{AM} = \begin{pmatrix}2\\1\\-2\end{pmatrix} \qquad \vect{AM}\cdot\vec{u} = 2+2+2 = 6 \qquad \norme{\vec{u}}^2 = 1+4+1 = 6 \]
\[ t = \frac{6}{6} = 1 \]

Donc $H = A + 1 \cdot \vec{u}$, soit $H(2, 2, 1)$.

Vérification : $\vect{MH} = \begin{pmatrix}-1\\1\\1\end{pmatrix}$ et $\vect{MH}\cdot\vec{u} = -1+2-1 = 0$. \checkmark
\end{exemple}

% ── VII. Distance d'un point à un plan ──
\section{Distance d'un point à un plan}

\begin{definition}[Distance d'un point à un plan]
La \textbf{distance} d'un point $M$ à un plan $\mathcal{P}$ est la longueur $MH$, où $H$ est le projeté orthogonal de $M$ sur $\mathcal{P}$.

C'est la plus courte distance entre $M$ et un point quelconque de $\mathcal{P}$.
\end{definition}

\begin{propriete}[Formule de distance]
Soit $\mathcal{P}$ un plan passant par un point $A$ et de vecteur normal $\vec{n}$. La distance du point $M$ au plan $\mathcal{P}$ est :
\[ d(M, \mathcal{P}) = \frac{\abs{\vect{AM}\cdot\vec{n}}}{\norme{\vec{n}}} \]
\end{propriete}

\begin{demonstration}
Soit $H$ le projeté orthogonal de $M$ sur $\mathcal{P}$. Alors $\vect{MH}$ est colinéaire à $\vec{n}$ : il existe $\lambda \in \R$ tel que $\vect{MH} = \lambda\vec{n}$.

On décompose : $\vect{AM} = \vect{AH} + \vect{HM}$, donc $\vect{AM}\cdot\vec{n} = \vect{AH}\cdot\vec{n} + \vect{HM}\cdot\vec{n}$.

Comme $H \in \mathcal{P}$, le vecteur $\vect{AH}$ est dans le plan $\mathcal{P}$, donc $\vect{AH}\cdot\vec{n} = 0$.

D'où $\vect{AM}\cdot\vec{n} = \vect{HM}\cdot\vec{n} = -\lambda\vec{n}\cdot\vec{n} = -\lambda\norme{\vec{n}}^2$.

Ainsi $\lambda = -\dfrac{\vect{AM}\cdot\vec{n}}{\norme{\vec{n}}^2}$ et :
\[ MH = \norme{\vect{MH}} = \abs{\lambda}\norme{\vec{n}} = \frac{\abs{\vect{AM}\cdot\vec{n}}}{\norme{\vec{n}}^2}\times\norme{\vec{n}} = \frac{\abs{\vect{AM}\cdot\vec{n}}}{\norme{\vec{n}}} \qed \]
\end{demonstration}

\begin{exemple}[Distance dans un cube]
Dans le cube $ABCDEFGH$ d'arête 1 muni du repère orthonormé $(A\,;\,\vect{AB},\vect{AD},\vect{AE})$, calculons la distance de $A$ au plan $(BDH)$.

Le plan $(BDH)$ passe par $B(1,0,0)$. On cherche un vecteur normal.
\[ \vect{BD} = \begin{pmatrix}-1\\1\\0\end{pmatrix} \qquad \vect{BH} = \begin{pmatrix}-1\\1\\1\end{pmatrix} \]

Un vecteur $\vec{n}\begin{pmatrix}a\\b\\c\end{pmatrix}$ normal au plan vérifie $\vec{n}\cdot\vect{BD}=0$ et $\vec{n}\cdot\vect{BH}=0$ :
\[ -a+b = 0 \implies a=b \qquad -a+b+c = 0 \implies c = 0 \]
Prenons $\vec{n} = \begin{pmatrix}1\\1\\1\end{pmatrix}$ (on vérifie : $-1+1+0=0$ et $-1+1+1 = 1 \neq 0$).

Corrigeons : $\vec{n}\cdot\vect{BH} = -1+1+1 = 1 \neq 0$. Reprenons. De $-a+b=0$, on a $a=b$. De $-a+b+c=0$, on a $c=0$. Donc $\vec{n} = \begin{pmatrix}1\\1\\0\end{pmatrix}$.

Vérifions : $\vec{n}\cdot\vect{BD} = -1+1+0 = 0$ \checkmark \quad $\vec{n}\cdot\vect{BH} = -1+1+0 = 0$ \checkmark

\medskip

Distance de $A(0,0,0)$ au plan :
\[ d(A, (BDH)) = \frac{\abs{\vect{BA}\cdot\vec{n}}}{\norme{\vec{n}}} = \frac{\abs{\begin{pmatrix}-1\\0\\0\end{pmatrix}\cdot\begin{pmatrix}1\\1\\0\end{pmatrix}}}{\sqrt{1+1+0}} = \frac{\abs{-1}}{\sqrt{2}} = \frac{1}{\sqrt{2}} = \frac{\sqrt{2}}{2} \]
\end{exemple}

% ── VIII. Applications ──
\section{Applications}

\begin{exemple}[Angle entre deux diagonales du cube]
Dans le cube $ABCDEFGH$ d'arête 1, calculons l'angle entre les diagonales $[AG]$ et $[BH]$.

Dans le repère orthonormé $(A\,;\,\vect{AB},\vect{AD},\vect{AE})$ :
\[ \vect{AG} = \begin{pmatrix}1\\1\\1\end{pmatrix} \qquad \vect{BH} = \begin{pmatrix}-1\\1\\1\end{pmatrix} \]

\[ \vect{AG}\cdot\vect{BH} = -1+1+1 = 1 \qquad \norme{\vect{AG}} = \sqrt{3} \qquad \norme{\vect{BH}} = \sqrt{3} \]

\[ \cos\theta = \frac{\vect{AG}\cdot\vect{BH}}{\norme{\vect{AG}}\times\norme{\vect{BH}}} = \frac{1}{3} \]

Donc $\theta = \arccos\left(\frac{1}{3}\right) \approx 70{,}5°$.
\end{exemple}

\begin{exemple}[Sphère et plan]
Dans un repère orthonormé, la sphère $\mathcal{S}$ de centre $\Omega(1, 2, 3)$ et de rayon $R = 5$ a pour équation :
\[ (x-1)^2 + (y-2)^2 + (z-3)^2 = 25 \]

Le plan $\mathcal{P}$ d'équation $x + 2y + 2z - 1 = 0$ a pour vecteur normal $\vec{n}\begin{pmatrix}1\\2\\2\end{pmatrix}$.

Distance de $\Omega$ au plan :
\[ d(\Omega, \mathcal{P}) = \frac{\abs{1 + 4 + 6 - 1}}{\sqrt{1+4+4}} = \frac{10}{3} \approx 3{,}33 \]

Comme $d(\Omega, \mathcal{P}) = \frac{10}{3} < 5 = R$, le plan coupe la sphère selon un \textbf{cercle} de rayon :
\[ r = \sqrt{R^2 - d^2} = \sqrt{25 - \frac{100}{9}} = \sqrt{\frac{125}{9}} = \frac{5\sqrt{5}}{3} \]
\end{exemple}

% ── Bilan ──
\section*{Bilan du chapitre}

\begin{retenir}
\begin{itemize}[leftmargin=*]
  \item \textbf{Produit scalaire :} $\vec{u}\cdot\vec{v} = \norme{\vec{u}}\norme{\vec{v}}\cos\theta$
  \item \textbf{En repère orthonormé :} $\vec{u}\cdot\vec{v} = xx'+yy'+zz'$
  \item \textbf{Norme :} $\norme{\vec{u}} = \sqrt{x^2+y^2+z^2}$
  \item \textbf{Distance :} $AB = \sqrt{(x_B-x_A)^2+(y_B-y_A)^2+(z_B-z_A)^2}$
  \item \textbf{Orthogonalité :} $\vec{u} \perp \vec{v} \Leftrightarrow \vec{u}\cdot\vec{v} = 0$
  \item \textbf{Formules de polarisation :} permettent de calculer $\vec{u}\cdot\vec{v}$ à partir des normes
  \item \textbf{Vecteur normal à un plan :} orthogonal à tout vecteur du plan
  \item \textbf{Plan de normale $\vec{n}$ passant par $A$ :} $M \in \mathcal{P} \Leftrightarrow \vect{AM}\cdot\vec{n} = 0$
  \item \textbf{Distance point-plan :} $d(M, \mathcal{P}) = \frac{\abs{\vect{AM}\cdot\vec{n}}}{\norme{\vec{n}}}$
  \item \textbf{Projeté orthogonal} sur une droite : $t = \frac{\vect{AM}\cdot\vec{u}}{\norme{\vec{u}}^2}$
\end{itemize}
\end{retenir}

\chapter{Représentations paramétriques et équations cartésiennes}

\section*{Introduction}

\textit{Comment décrire une droite de l'espace par des équations ? Comment traduire algébriquement qu'un point appartient à un plan ? Comment déterminer l'intersection de deux plans ?}

\medskip

Les chapitres précédents décrivent droites et plans par des points et des vecteurs. Ce chapitre introduit les outils \textbf{algébriques} permettant de travailler avec des \textbf{équations} : représentations paramétriques pour les droites, équations cartésiennes pour les plans. Ce passage de la géométrie à l'algèbre est l'un des ponts les plus puissants des mathématiques : il transforme des questions géométriques (intersection, parallélisme, distance) en résolutions de systèmes d'équations.

\medskip

\textbf{Objectifs du chapitre :}
\begin{itemize}[leftmargin=*]
  \item Écrire et exploiter la représentation paramétrique d'une droite
  \item Déterminer l'équation cartésienne d'un plan
  \item Étudier les positions relatives de droites et plans par le calcul
  \item Résoudre des systèmes d'équations linéaires à 3 inconnues
  \item Appliquer ces outils à des problèmes de géométrie (projeté, symétrique, angle)
\end{itemize}

\begin{histoire}[Géométrie analytique dans l'espace]
La géométrie analytique, initiée par \textbf{Descartes} et \textbf{Fermat} au XVII\ieme{} siècle dans le plan, est étendue à l'espace par \textbf{Euler} (1707--1783) et \textbf{Monge} (1746--1818). C'est Euler qui, dans ses \textit{Introductio in analysin infinitorum} (1748), systématise l'usage des trois coordonnées $(x,y,z)$ et des équations de plans et de droites dans l'espace. \textbf{Gaspard Monge}, fondateur de la géométrie descriptive, développe les méthodes de projection et de représentation qui restent à la base du dessin technique et de la CAO (conception assistée par ordinateur).
\end{histoire}

% ── I. Représentation paramétrique d'une droite ──
\section{Représentation paramétrique d'une droite}

\begin{definition}[Représentation paramétrique]
Soit $d$ la droite passant par le point $A(x_A, y_A, z_A)$ de vecteur directeur $\vec{u}\begin{pmatrix}a\\b\\c\end{pmatrix}$.

Un point $M(x, y, z)$ appartient à $d$ si et seulement s'il existe un réel $t$ (le \textbf{paramètre}) tel que :
\[ \begin{cases} x = x_A + ta \\ y = y_A + tb \\ z = z_A + tc \end{cases} \qquad t \in \R \]
Ce système est une \textbf{représentation paramétrique} de $d$.
\end{definition}

\begin{center}
\begin{tikzpicture}[scale=1.2,>=Stealth]
  \coordinate (A) at (0,0);
  \coordinate (u) at (2,1);
  % Droite
  \draw[thick,indigo] ($(A)-1.5*(u)$) -- ($(A)+2*(u)$);
  % Point A
  \fill[indigo] (A) circle (2pt) node[below left] {$A(x_A,y_A,z_A)$};
  % Vecteur directeur
  \draw[->,very thick,rougeerr] (A) -- (u) node[above left] {$\vec{u}\begin{pmatrix}a\\b\\c\end{pmatrix}$};
  % Point M
  \coordinate (M) at ($(A)+1.5*(u)$);
  \fill[vertdef] (M) circle (2pt) node[above left] {$M = A + t\vec{u}$};
  \draw[->,thick,vertdef,dashed] (A) -- (M);
  % t annotation
  \node[below right,font=\small\itshape,vertdef] at ($(A)!0.5!(M)$) {$t\vec{u}$};
\end{tikzpicture}
\end{center}

\begin{attention}[Non-unicité de la représentation]
La représentation paramétrique d'une droite n'est \textbf{pas unique}. En changeant le point de départ ou le vecteur directeur (en le multipliant par un réel non nul), on obtient une représentation différente de la \textbf{même droite}.
\end{attention}

\begin{exemple}[Écrire une représentation paramétrique]
La droite $d$ passe par $A(1, -2, 3)$ avec le vecteur directeur $\vec{u}\begin{pmatrix}2\\1\\-1\end{pmatrix}$.

Représentation paramétrique :
\[ \begin{cases} x = 1 + 2t \\ y = -2 + t \\ z = 3 - t \end{cases} \qquad t \in \R \]

\textbf{Vérification :} pour $t = 0$, on retrouve $A(1, -2, 3)$ \checkmark. Pour $t = 1$ : $M(3, -1, 2)$.
\end{exemple}

\begin{methode}[Lire point et vecteur directeur]
Étant donnée la représentation paramétrique $\begin{cases} x = 3 - t \\ y = 1 + 2t \\ z = -4 + 5t \end{cases}$, on lit directement :
\begin{itemize}
  \item Le point $A$ : coordonnées obtenues pour $t = 0$, soit $A(3, 1, -4)$.
  \item Le vecteur directeur $\vec{u}$ : coefficients de $t$, soit $\vec{u}\begin{pmatrix}-1\\2\\5\end{pmatrix}$.
\end{itemize}
\end{methode}

\begin{methode}[Tester l'appartenance d'un point à une droite]
Pour vérifier si $M(x_0, y_0, z_0)$ appartient à la droite $d : \begin{cases} x = x_A + ta \\ y = y_A + tb \\ z = z_A + tc \end{cases}$ :
\begin{enumerate}
  \item Résoudre $x_0 = x_A + ta$ pour trouver $t$.
  \item Vérifier que ce même $t$ satisfait les deux autres équations.
  \item Si oui : $M \in d$. Si non : $M \notin d$.
\end{enumerate}
\end{methode}

\begin{exemple}[Test d'appartenance]
Soit $d : \begin{cases} x = 2 + 3t \\ y = -1 + t \\ z = 4 - 2t \end{cases}$. Le point $P(8, 1, 0)$ appartient-il à $d$ ?

De $x = 2 + 3t = 8$, on tire $t = 2$.

Vérification : $y = -1 + 2 = 1$ \checkmark \quad $z = 4 - 4 = 0$ \checkmark

Donc $P \in d$.

\medskip

Le point $Q(5, 0, 3)$ appartient-il à $d$ ? De $x = 2 + 3t = 5$, on tire $t = 1$. Vérification : $y = -1 + 1 = 0$ \checkmark \quad $z = 4 - 2 = 2 \neq 3$ $\times$

Donc $Q \notin d$.
\end{exemple}

% ── II. Équation cartésienne d'un plan ──
\section{Équation cartésienne d'un plan}

\begin{propriete}[Équation cartésienne]
Dans un repère orthonormé, tout plan $\mathcal{P}$ admet une \textbf{équation cartésienne} de la forme :
\[ ax + by + cz + d = 0 \]
où $a$, $b$, $c$ ne sont pas tous nuls.

Réciproquement, l'ensemble des points $M(x,y,z)$ vérifiant $ax + by + cz + d = 0$ (avec $(a,b,c) \neq (0,0,0)$) est un plan.
\end{propriete}

\begin{propriete}[Lien avec le vecteur normal]
Le plan d'équation $ax + by + cz + d = 0$ a pour vecteur normal :
\[ \vec{n} = \begin{pmatrix}a\\b\\c\end{pmatrix} \]
Les coefficients de $x$, $y$, $z$ sont les coordonnées du vecteur normal.
\end{propriete}

\begin{demonstration}
Le plan $\mathcal{P}$ passe par un point $A(x_A, y_A, z_A)$ et a pour vecteur normal $\vec{n}\begin{pmatrix}a\\b\\c\end{pmatrix}$. Un point $M(x,y,z)$ appartient à $\mathcal{P}$ si et seulement si $\vect{AM}\cdot\vec{n} = 0$ :
\begin{align*}
\vect{AM}\cdot\vec{n} &= a(x - x_A) + b(y - y_A) + c(z - z_A) = 0 \\
&\Leftrightarrow ax + by + cz + \underbrace{(-ax_A - by_A - cz_A)}_{= d} = 0 \qed
\end{align*}
\end{demonstration}

\begin{center}
\begin{tikzpicture}[scale=1.2,>=Stealth]
  % Plan
  \fill[indigobg,opacity=0.4] (-2.5,-0.8) -- (2,-0.8) -- (2.8,1.8) -- (-1.7,1.8) -- cycle;
  \node[indigo] at (2.5,1.5) {$\mathcal{P} : ax+by+cz+d=0$};
  % Point A
  \fill (0,0.3) circle (1.5pt) node[below left] {$A$};
  % Vecteur normal
  \draw[->,very thick,rougeerr] (0,0.3) -- (0,2.5) node[right] {$\vec{n}\begin{pmatrix}a\\b\\c\end{pmatrix}$};
  % Angle droit
  \draw[rougeerr] (0.15,0.3) -- (0.15,0.5) -- (0,0.5);
\end{tikzpicture}
\end{center}

\begin{methode}[Déterminer l'équation d'un plan connaissant un vecteur normal et un point]
On connaît un point $A(x_A, y_A, z_A)$ du plan et un vecteur normal $\vec{n}\begin{pmatrix}a\\b\\c\end{pmatrix}$.
\begin{enumerate}
  \item L'équation est de la forme $ax + by + cz + d = 0$.
  \item On détermine $d$ en substituant les coordonnées de $A$ : $d = -ax_A - by_A - cz_A$.
\end{enumerate}
\end{methode}

\begin{exemple}[Équation à partir d'un point et d'un vecteur normal]
Le plan $\mathcal{P}$ passe par $A(1, -2, 3)$ et a pour vecteur normal $\vec{n}\begin{pmatrix}2\\-1\\4\end{pmatrix}$.

Équation : $2x - y + 4z + d = 0$.

On substitue $A$ : $2(1) - (-2) + 4(3) + d = 0 \implies 2 + 2 + 12 + d = 0 \implies d = -16$.

Donc $\mathcal{P} : 2x - y + 4z - 16 = 0$.
\end{exemple}

\begin{methode}[Déterminer l'équation d'un plan passant par trois points]
Soient $A$, $B$, $C$ trois points non alignés. Pour trouver l'équation du plan $(ABC)$ :
\begin{enumerate}
  \item Calculer $\vect{AB}$ et $\vect{AC}$.
  \item Chercher $\vec{n}\begin{pmatrix}a\\b\\c\end{pmatrix}$ tel que $\vec{n}\cdot\vect{AB} = 0$ et $\vec{n}\cdot\vect{AC} = 0$ (système de 2 équations à 3 inconnues : on fixe une inconnue).
  \item Déterminer $d$ avec les coordonnées de $A$.
\end{enumerate}
\end{methode}

\begin{exemple}[Plan passant par trois points]
Trouver l'équation du plan $(ABC)$ avec $A(1,0,0)$, $B(0,2,0)$, $C(0,0,3)$.

\[ \vect{AB} = \begin{pmatrix}-1\\2\\0\end{pmatrix} \qquad \vect{AC} = \begin{pmatrix}-1\\0\\3\end{pmatrix} \]

On cherche $\vec{n}\begin{pmatrix}a\\b\\c\end{pmatrix}$ tel que :
\[ \begin{cases} -a + 2b = 0 \\ -a + 3c = 0 \end{cases} \implies \begin{cases} a = 2b \\ a = 3c \end{cases} \]
On pose $a = 6$ : $b = 3$, $c = 2$, soit $\vec{n}\begin{pmatrix}6\\3\\2\end{pmatrix}$.

Équation : $6x + 3y + 2z + d = 0$. Avec $A(1,0,0)$ : $6 + d = 0$, donc $d = -6$.

\[ \boxed{\mathcal{P} : 6x + 3y + 2z - 6 = 0} \]

Vérification : $B(0,2,0)$ : $0 + 6 + 0 - 6 = 0$ \checkmark \quad $C(0,0,3)$ : $0 + 0 + 6 - 6 = 0$ \checkmark
\end{exemple}

% ── III. Positions relatives par le calcul ──
\section{Positions relatives par le calcul}

\subsection{Intersection de deux plans}

\begin{methode}[Deux plans]
Soient $\mathcal{P}_1 : a_1x + b_1y + c_1z + d_1 = 0$ et $\mathcal{P}_2 : a_2x + b_2y + c_2z + d_2 = 0$.

On résout le système :
\[ \begin{cases} a_1x + b_1y + c_1z + d_1 = 0 \\ a_2x + b_2y + c_2z + d_2 = 0 \end{cases} \]

\begin{itemize}
  \item \textbf{Si les vecteurs normaux sont colinéaires} ($\vec{n}_1 = \lambda\vec{n}_2$) :
  \begin{itemize}
    \item $d_1 = \lambda d_2$ : plans \textbf{confondus}
    \item $d_1 \neq \lambda d_2$ : plans \textbf{strictement parallèles}
  \end{itemize}
  \item \textbf{Si les vecteurs normaux ne sont pas colinéaires :} les plans sont \textbf{sécants} selon une droite. Le système admet une infinité de solutions dépendant d'un paramètre (la droite d'intersection).
\end{itemize}
\end{methode}

\begin{exemple}[Intersection de deux plans]
Soient $\mathcal{P}_1 : x + y - z + 1 = 0$ et $\mathcal{P}_2 : 2x - y + z - 3 = 0$.

Les vecteurs normaux $\vec{n}_1\begin{pmatrix}1\\1\\-1\end{pmatrix}$ et $\vec{n}_2\begin{pmatrix}2\\-1\\1\end{pmatrix}$ ne sont pas colinéaires : les plans sont sécants.

On résout le système en posant $z = t$ (paramètre) :
\[ \begin{cases} x + y = z - 1 = t - 1 \\ 2x - y = -z + 3 = -t + 3 \end{cases} \]
En additionnant : $3x = 2$, donc $x = \frac{2}{3}$.

Hmm, cela ne dépend pas de $t$, reprenons. Additionnons les deux équations :
\[ 3x = (t-1) + (-t+3) = 2 \implies x = \frac{2}{3} \]
Puis $y = t - 1 - \frac{2}{3} = t - \frac{5}{3}$.

La droite d'intersection a pour représentation paramétrique :
\[ \begin{cases} x = \frac{2}{3} \\ y = -\frac{5}{3} + t \\ z = t \end{cases} \qquad t \in \R \]

Vecteur directeur : $\vec{u}\begin{pmatrix}0\\1\\1\end{pmatrix}$, point : $\left(\frac{2}{3}, -\frac{5}{3}, 0\right)$.
\end{exemple}

\subsection{Intersection d'une droite et d'un plan}

\begin{methode}[Droite et plan]
Soit $d : \begin{cases} x = x_A + ta \\ y = y_A + tb \\ z = z_A + tc \end{cases}$ et $\mathcal{P} : \alpha x + \beta y + \gamma z + \delta = 0$.

On substitue les expressions paramétriques dans l'équation du plan :
\[ \alpha(x_A + ta) + \beta(y_A + tb) + \gamma(z_A + tc) + \delta = 0 \]

C'est une équation en $t$ :
\begin{itemize}
  \item \textbf{Un coefficient de $t$ non nul :} une solution unique $\rightarrow$ droite \textbf{sécante} au plan (un point d'intersection).
  \item \textbf{Coefficient de $t$ nul, reste nul :} identité $\rightarrow$ droite \textbf{incluse} dans le plan.
  \item \textbf{Coefficient de $t$ nul, reste non nul :} contradiction $\rightarrow$ droite \textbf{parallèle} au plan (pas d'intersection).
\end{itemize}
\end{methode}

\begin{exemple}[Intersection droite-plan]
Soit $d : \begin{cases} x = 1 + 2t \\ y = -1 + t \\ z = 3 - t \end{cases}$ et $\mathcal{P} : x + 3y - z + 2 = 0$.

Substitution : $(1+2t) + 3(-1+t) - (3-t) + 2 = 0$

$1 + 2t - 3 + 3t - 3 + t + 2 = 0 \implies 6t - 3 = 0 \implies t = \frac{1}{2}$

Point d'intersection : $M\left(2, -\frac{1}{2}, \frac{5}{2}\right)$.
\end{exemple}

\subsection{Intersection de deux droites}

\begin{methode}[Deux droites]
Pour étudier l'intersection de $d_1 : \begin{cases} x = x_1 + t a_1 \\ y = y_1 + t b_1 \\ z = z_1 + t c_1 \end{cases}$ et $d_2 : \begin{cases} x = x_2 + s a_2 \\ y = y_2 + s b_2 \\ z = z_2 + s c_2 \end{cases}$ :

On résout le système (3 équations, 2 inconnues $t$ et $s$) :
\[ \begin{cases} x_1 + t a_1 = x_2 + s a_2 \\ y_1 + t b_1 = y_2 + s b_2 \\ z_1 + t c_1 = z_2 + s c_2 \end{cases} \]

\begin{itemize}
  \item \textbf{Solution unique $(t, s)$ :} les droites sont \textbf{sécantes}.
  \item \textbf{Infinité de solutions :} les droites sont \textbf{confondues} (si vecteurs directeurs colinéaires) ou la même droite.
  \item \textbf{Pas de solution :} les droites sont \textbf{parallèles} (vecteurs colinéaires) ou \textbf{non coplanaires} (vecteurs non colinéaires).
\end{itemize}
\end{methode}

% ── IV. Systèmes d'équations linéaires ──
\section{Systèmes d'équations linéaires}

\subsection{Intersection de trois plans}

\begin{propriete}
L'intersection de trois plans se ramène à la résolution d'un système de trois équations à trois inconnues :
\[ \begin{cases} a_1 x + b_1 y + c_1 z = d_1 \\ a_2 x + b_2 y + c_2 z = d_2 \\ a_3 x + b_3 y + c_3 z = d_3 \end{cases} \]

Les situations possibles sont :
\begin{itemize}
  \item \textbf{Solution unique :} les trois plans se coupent en un seul point.
  \item \textbf{Infinité de solutions (1 paramètre) :} les trois plans se coupent selon une droite.
  \item \textbf{Infinité de solutions (2 paramètres) :} au moins deux plans sont confondus.
  \item \textbf{Aucune solution :} configuration incompatible (plans parallèles, ou en « prisme »).
\end{itemize}
\end{propriete}

\begin{exemple}[Trois plans sécants en un point]
Résoudre :
\[ \begin{cases} x + y + z = 6 \\ 2x - y + z = 3 \\ x + 2y - z = 3 \end{cases} \]

$L_2 - 2L_1$ : $-3y - z = -9$, soit $3y + z = 9$.

$L_3 - L_1$ : $y - 2z = -3$.

Du système réduit :
\[ \begin{cases} 3y + z = 9 \\ y - 2z = -3 \end{cases} \]

$3(y-2z) = -9$, donc $3y - 6z = -9$. Par soustraction avec $3y + z = 9$ : $-7z = -18$, soit $z = \frac{18}{7}$.

Puis $y = -3 + 2z = -3 + \frac{36}{7} = \frac{15}{7}$.

Et $x = 6 - y - z = 6 - \frac{15}{7} - \frac{18}{7} = 6 - \frac{33}{7} = \frac{9}{7}$.

Solution unique : $\left(\frac{9}{7}, \frac{15}{7}, \frac{18}{7}\right)$. Les trois plans se coupent en ce point.
\end{exemple}

\subsection{Vérifier qu'une famille de vecteurs forme une base}

\begin{methode}[Test de base]
Trois vecteurs $\vec{u}\begin{pmatrix}a_1\\b_1\\c_1\end{pmatrix}$, $\vec{v}\begin{pmatrix}a_2\\b_2\\c_2\end{pmatrix}$, $\vec{w}\begin{pmatrix}a_3\\b_3\\c_3\end{pmatrix}$ forment une base si et seulement si le système :
\[ \begin{cases} a_1 x + a_2 y + a_3 z = 0 \\ b_1 x + b_2 y + b_3 z = 0 \\ c_1 x + c_2 y + c_3 z = 0 \end{cases} \]
n'admet que la solution triviale $(0, 0, 0)$.

\textit{Équivalent :} les trois vecteurs ne sont pas coplanaires.
\end{methode}

% ── V. Applications ──
\section{Applications}

\begin{exemple}[Projeté orthogonal par les équations]
Dans un repère orthonormé, trouver le projeté orthogonal $H$ de $M(3, 1, -1)$ sur le plan $\mathcal{P} : 2x - y + 2z - 1 = 0$.

Le vecteur normal est $\vec{n}\begin{pmatrix}2\\-1\\2\end{pmatrix}$.

La droite $(MH)$ est perpendiculaire à $\mathcal{P}$, donc dirigée par $\vec{n}$. Sa représentation paramétrique :
\[ \begin{cases} x = 3 + 2t \\ y = 1 - t \\ z = -1 + 2t \end{cases} \]

On substitue dans l'équation de $\mathcal{P}$ :
\[ 2(3+2t) - (1-t) + 2(-1+2t) - 1 = 0 \]
\[ 6 + 4t - 1 + t - 2 + 4t - 1 = 0 \implies 9t + 2 = 0 \implies t = -\frac{2}{9} \]

Donc $H\left(3 - \frac{4}{9},\; 1 + \frac{2}{9},\; -1 - \frac{4}{9}\right) = H\left(\frac{23}{9},\; \frac{11}{9},\; -\frac{13}{9}\right)$.
\end{exemple}

\begin{methode}[Trouver le symétrique d'un point par rapport à un plan]
Le symétrique $M'$ de $M$ par rapport au plan $\mathcal{P}$ vérifie : $H$ est le milieu de $[MM']$, où $H$ est le projeté de $M$ sur $\mathcal{P}$.

\begin{enumerate}
  \item Calculer $H$ (projeté orthogonal, méthode ci-dessus).
  \item Utiliser la relation du milieu : $M' = 2H - M$, soit :
  \[ x_{M'} = 2x_H - x_M, \qquad y_{M'} = 2y_H - y_M, \qquad z_{M'} = 2z_H - z_M \]
\end{enumerate}
\end{methode}

\begin{exemple}[Symétrique par rapport à un plan]
Reprenons l'exemple précédent : $M(3, 1, -1)$, $H\left(\frac{23}{9}, \frac{11}{9}, -\frac{13}{9}\right)$.

\[ M' = 2H - M = \left(\frac{46}{9} - 3,\; \frac{22}{9} - 1,\; -\frac{26}{9} + 1\right) = \left(\frac{19}{9},\; \frac{13}{9},\; -\frac{17}{9}\right) \]

Vérification : $H$ est bien le milieu de $[MM']$ :
\[ \frac{3 + \frac{19}{9}}{2} = \frac{\frac{27+19}{9}}{2} = \frac{46}{18} = \frac{23}{9} \checkmark \]
\end{exemple}

\begin{exemple}[Angle entre une droite et un plan]
Soit la droite $d$ de vecteur directeur $\vec{u}\begin{pmatrix}1\\2\\2\end{pmatrix}$ et le plan $\mathcal{P}$ de vecteur normal $\vec{n}\begin{pmatrix}2\\-1\\2\end{pmatrix}$.

L'angle $\alpha$ entre $d$ et $\mathcal{P}$ est le complémentaire de l'angle entre $\vec{u}$ et $\vec{n}$ :
\[ \sin\alpha = \abs{\cos(\vec{u}, \vec{n})} = \frac{\abs{\vec{u}\cdot\vec{n}}}{\norme{\vec{u}} \times \norme{\vec{n}}} \]

\[ \vec{u}\cdot\vec{n} = 2 - 2 + 4 = 4 \qquad \norme{\vec{u}} = 3 \qquad \norme{\vec{n}} = 3 \]

\[ \sin\alpha = \frac{4}{9} \implies \alpha = \arcsin\left(\frac{4}{9}\right) \approx 26{,}4° \]
\end{exemple}

\begin{attention}[Angle droite-plan vs angle entre vecteurs]
L'angle entre une droite et un plan est compris entre $0°$ et $90°$. C'est le \textbf{sinus} de l'angle qui fait intervenir le produit scalaire avec le vecteur normal, et non le cosinus. Ne pas confondre avec l'angle entre deux droites ou entre deux plans.
\end{attention}

\begin{propriete}[Distance d'un point à un plan — formule avec l'équation cartésienne]
Si $\mathcal{P} : ax + by + cz + d = 0$ et $M(x_0, y_0, z_0)$ :
\[ d(M, \mathcal{P}) = \frac{\abs{ax_0 + by_0 + cz_0 + d}}{\sqrt{a^2 + b^2 + c^2}} \]
\end{propriete}

\begin{demonstration}
Soit $A$ un point de $\mathcal{P}$, donc $ax_A + by_A + cz_A + d = 0$. Le vecteur normal est $\vec{n}\begin{pmatrix}a\\b\\c\end{pmatrix}$.

D'après la formule du chapitre précédent :
\begin{align*}
d(M, \mathcal{P}) &= \frac{\abs{\vect{AM}\cdot\vec{n}}}{\norme{\vec{n}}} = \frac{\abs{a(x_0-x_A)+b(y_0-y_A)+c(z_0-z_A)}}{\sqrt{a^2+b^2+c^2}} \\
&= \frac{\abs{ax_0+by_0+cz_0 - (ax_A+by_A+cz_A)}}{\sqrt{a^2+b^2+c^2}} = \frac{\abs{ax_0+by_0+cz_0+d}}{\sqrt{a^2+b^2+c^2}} \qed
\end{align*}
\end{demonstration}

% ── Bilan ──
\section*{Bilan du chapitre}

\begin{retenir}
\begin{itemize}[leftmargin=*]
  \item \textbf{Représentation paramétrique d'une droite :} $\begin{cases} x = x_A + ta \\ y = y_A + tb \\ z = z_A + tc \end{cases}$, $t \in \R$
  \item Le point se lit pour $t=0$, le vecteur directeur est $(a,b,c)$
  \item La représentation n'est \textbf{pas unique}
  \item \textbf{Équation cartésienne d'un plan :} $ax + by + cz + d = 0$, vecteur normal $\vec{n}\begin{pmatrix}a\\b\\c\end{pmatrix}$
  \item \textbf{Méthode pour trouver l'équation :} vecteur normal $+$ un point, ou trois points non alignés
  \item \textbf{Intersection droite-plan :} substituer la param. dans l'équation $\rightarrow$ 1 solution, 0 solution, ou $\infty$
  \item \textbf{Intersection de deux plans :} résoudre le système $\rightarrow$ droite, vide, ou plan
  \item \textbf{Projeté orthogonal :} droite perpendiculaire au plan passant par $M$, puis intersection
  \item \textbf{Symétrique :} $M' = 2H - M$
  \item \textbf{Angle droite-plan :} $\sin\alpha = \frac{\abs{\vec{u}\cdot\vec{n}}}{\norme{\vec{u}}\norme{\vec{n}}}$
  \item \textbf{Distance point-plan :} $d(M, \mathcal{P}) = \frac{\abs{ax_0+by_0+cz_0+d}}{\sqrt{a^2+b^2+c^2}}$
\end{itemize}
\end{retenir}

\part{Analyse}

% ─────────────────────────────────────────────────────────────────────────────
% CHAPITRE 5 — SUITES : LIMITES ET CONVERGENCE
% ─────────────────────────────────────────────────────────────────────────────

\chapter{Suites : limites et convergence}

\section*{Introduction}

\textit{En Première, on a étudié les suites arithmétiques et géométriques, leur sens de variation et leurs sommes. Mais une question fondamentale restait ouverte : vers quoi « tend » une suite quand $n$ devient très grand ? La suite $\ds\frac{1}{n}$ finit-elle par s'annuler ? La suite $0{,}9^n$ atteint-elle zéro ?}

\medskip

La notion de \textbf{limite} permet de formaliser ce comportement à l'infini. Elle est au cœur de l'analyse et constitue le pont entre le discret (les suites) et le continu (les fonctions).

\medskip

\textbf{Objectifs du chapitre :}
\begin{itemize}[leftmargin=*]
  \item Définir la limite d'une suite (finie ou infinie)
  \item Connaître les limites de référence ($q^n$, $\ds\frac{1}{n^\alpha}$)
  \item Calculer des limites à l'aide de théorèmes d'opérations
  \item Lever des formes indéterminées
  \item Utiliser les théorèmes de comparaison (encadrement, gendarmes)
  \item Démontrer par récurrence
  \item Appliquer le théorème de convergence monotone
\end{itemize}

\begin{histoire}[Cauchy et la rigueur des limites]
Pendant des siècles, les mathématiciens ont raisonné sur les limites de manière intuitive. C'est \textbf{Augustin-Louis Cauchy} (1789--1857) qui, dans son \textit{Cours d'analyse} (1821), donne la première définition rigoureuse de la limite d'une suite à l'aide de la formulation $\varepsilon$-$N$ : « la suite $(u_n)$ a pour limite $\ell$ si, pour tout $\varepsilon > 0$, il existe un rang $N$ tel que, pour tout $n \geqslant N$, $|u_n - \ell| < \varepsilon$ ». Cette formalisation a révolutionné l'analyse et reste le fondement de tout le calcul moderne.
\end{histoire}

% ── I. Limite d'une suite ──
\section{Limite d'une suite}

\subsection{Suite convergente}

\begin{definition}[Limite finie]
On dit que la suite $(u_n)$ \textbf{converge} vers le réel $\ell$ (ou a pour limite $\ell$) si les termes $u_n$ se rapprochent aussi près que l'on veut de $\ell$ quand $n$ devient suffisamment grand.

On note : $\ds\lim_{n \to +\infty} u_n = \ell$.

On dit alors que la suite est \textbf{convergente}.
\end{definition}

\textit{Formulation rigoureuse (pour les curieux) :} pour tout $\varepsilon > 0$, il existe un entier $N$ tel que, pour tout $n \geqslant N$, on a $|u_n - \ell| < \varepsilon$.

\begin{exemple}
La suite $u_n = \ds\frac{1}{n}$ converge vers $0$ : plus $n$ est grand, plus $\ds\frac{1}{n}$ est proche de $0$.

\medskip

La suite $u_n = \ds\frac{2n+1}{n} = 2 + \frac{1}{n}$ converge vers $2$.
\end{exemple}

\begin{center}
\begin{tikzpicture}
\begin{axis}[
  width=10cm, height=6cm,
  xlabel={$n$}, ylabel={$u_n$},
  xmin=0, xmax=20, ymin=-0.1, ymax=1.2,
  xtick={0,5,...,20},
  grid=major, grid style={gray!20},
  axis lines=middle,
  title={Suite $u_n = \frac{1}{n}$ : convergence vers $0$},
  title style={font=\sffamily\bfseries}
]
\addplot[only marks, mark=*, mark size=2pt, indigo] coordinates {
  (1,1) (2,0.5) (3,0.333) (4,0.25) (5,0.2)
  (6,0.167) (7,0.143) (8,0.125) (9,0.111) (10,0.1)
  (11,0.091) (12,0.083) (13,0.077) (14,0.071) (15,0.067)
  (16,0.0625) (17,0.059) (18,0.056) (19,0.053) (20,0.05)
};
\addplot[rougeerr, dashed, thick, domain=0:20] {0};
\node[rougeerr, anchor=west] at (axis cs:15,0.15) {$\ell = 0$};
\end{axis}
\end{tikzpicture}
\captionof{figure}{Les termes de la suite $u_n = \frac{1}{n}$ se rapprochent de $0$.}
\end{center}

\subsection{Suite divergente vers $+\infty$ ou $-\infty$}

\begin{definition}[Limite infinie]
\begin{itemize}
  \item On dit que $(u_n)$ \textbf{diverge vers $+\infty$} si, pour tout réel $A$, il existe un rang $N$ tel que, pour tout $n \geqslant N$, $u_n > A$.

  On note : $\ds\lim_{n \to +\infty} u_n = +\infty$.

  \item On dit que $(u_n)$ \textbf{diverge vers $-\infty$} si, pour tout réel $A$, il existe un rang $N$ tel que, pour tout $n \geqslant N$, $u_n < A$.

  On note : $\ds\lim_{n \to +\infty} u_n = -\infty$.
\end{itemize}
\end{definition}

\begin{exemple}
$\ds\lim_{n \to +\infty} n^2 = +\infty$, \quad $\ds\lim_{n \to +\infty} (-n) = -\infty$.
\end{exemple}

\begin{attention}[Suite divergente sans limite]
Certaines suites divergent sans tendre ni vers $+\infty$ ni vers $-\infty$. Par exemple, $u_n = (-1)^n$ n'a pas de limite : elle oscille entre $1$ et $-1$.
\end{attention}

\subsection{Unicité de la limite}

\begin{propriete}
Si une suite $(u_n)$ converge, sa limite est \textbf{unique}.
\end{propriete}

% ── II. Limites de référence ──
\section{Limites de référence}

\begin{propriete}[Limites de référence]
\begin{center}
\renewcommand{\arraystretch}{1.8}
\begin{tabular}{cc}
\toprule
\textbf{Suite} & \textbf{Limite} \\
\midrule
$\ds\frac{1}{n}$ & $0$ \\
$\ds\frac{1}{n^2}$ & $0$ \\
$\ds\frac{1}{\sqrt{n}}$ & $0$ \\
$n$ & $+\infty$ \\
$n^2$ & $+\infty$ \\
$\sqrt{n}$ & $+\infty$ \\
\bottomrule
\end{tabular}
\end{center}
\end{propriete}

\subsection{Comportement de $q^n$}

\begin{propriete}[Limite de $q^n$]
Soit $q$ un réel.
\begin{center}
\renewcommand{\arraystretch}{1.6}
\begin{tabular}{ccc}
\toprule
\textbf{Valeur de $q$} & \textbf{Comportement de $q^n$} & \textbf{Limite} \\
\midrule
$q > 1$ & croissance vers $+\infty$ & $+\infty$ \\
$q = 1$ & constante égale à $1$ & $1$ \\
$-1 < q < 1$ & décroissance vers $0$ & $0$ \\
$q \leqslant -1$ & oscillation & pas de limite \\
\bottomrule
\end{tabular}
\end{center}
\end{propriete}

\begin{exemple}
\begin{itemize}
  \item $\ds\lim_{n \to +\infty} 1{,}05^n = +\infty$ \quad (croissance exponentielle)
  \item $\ds\lim_{n \to +\infty} 0{,}9^n = 0$ \quad (décroissance vers 0)
  \item $\ds\lim_{n \to +\infty} \left(\frac{2}{3}\right)^{\!n} = 0$ \quad (car $\left|\frac{2}{3}\right| < 1$)
  \item $(-1)^n$ n'a pas de limite \quad (oscille entre $-1$ et $1$)
\end{itemize}
\end{exemple}

\begin{attention}
Le résultat $q^n \to 0$ quand $|q| < 1$ est fondamental en probabilités (loi géométrique) et en modélisation (décroissance radioactive, amortissement).
\end{attention}

% ── III. Opérations sur les limites ──
\section{Opérations sur les limites}

\subsection{Somme de limites}

\begin{propriete}[Limite d'une somme]
Si $\ds\lim u_n = \ell$ et $\ds\lim v_n = \ell'$, alors $\ds\lim (u_n + v_n) = \ell + \ell'$.

\medskip

\begin{center}
\renewcommand{\arraystretch}{1.6}
\begin{tabular}{c|cccc}
$\lim u_n$ & $\ell$ & $\ell$ & $+\infty$ & $+\infty$ \\
$\lim v_n$ & $\ell'$ & $+\infty$ & $+\infty$ & $-\infty$ \\
\hline
$\lim (u_n+v_n)$ & $\ell+\ell'$ & $+\infty$ & $+\infty$ & \textbf{FI}
\end{tabular}
\end{center}

\textbf{FI} = Forme Indéterminée : $+\infty - \infty$.
\end{propriete}

\subsection{Produit de limites}

\begin{propriete}[Limite d'un produit]
Si $\ds\lim u_n = \ell$ et $\ds\lim v_n = \ell'$, alors $\ds\lim (u_n \cdot v_n) = \ell \cdot \ell'$.

\medskip

\begin{center}
\renewcommand{\arraystretch}{1.6}
\begin{tabular}{c|cccc}
$\lim u_n$ & $\ell \neq 0$ & $\ell$ & $+\infty$ & $0$ \\
$\lim v_n$ & $\ell'$ & $\pm\infty$ & $+\infty$ & $\pm\infty$ \\
\hline
$\lim (u_n \cdot v_n)$ & $\ell \cdot \ell'$ & $\pm\infty$ & $+\infty$ & \textbf{FI}
\end{tabular}
\end{center}

\textbf{FI} = Forme Indéterminée : $0 \times \infty$.
\end{propriete}

\subsection{Quotient de limites}

\begin{propriete}[Limite d'un quotient]
Si $\ds\lim u_n = \ell$ et $\ds\lim v_n = \ell' \neq 0$, alors $\ds\lim \frac{u_n}{v_n} = \frac{\ell}{\ell'}$.

\medskip

\begin{center}
\renewcommand{\arraystretch}{1.6}
\begin{tabular}{c|cccc}
$\lim u_n$ & $\ell$ & $\pm\infty$ & $\pm\infty$ & $0$ \\
$\lim v_n$ & $0^\pm$ & $\ell' \neq 0$ & $\pm\infty$ & $0$ \\
\hline
$\ds\lim \frac{u_n}{v_n}$ & $\pm\infty$ & $\pm\infty$ & \textbf{FI} & \textbf{FI}
\end{tabular}
\end{center}

\textbf{FI} = Formes Indéterminées : $\ds\frac{\infty}{\infty}$ et $\ds\frac{0}{0}$.
\end{propriete}

\begin{methode}[Lever une forme indéterminée]
Les quatre formes indéterminées sont : $+\infty - \infty$, \; $0 \times \infty$, \; $\ds\frac{\infty}{\infty}$, \; $\ds\frac{0}{0}$.

Pour lever une FI, on peut :
\begin{enumerate}
  \item \textbf{Factoriser par le terme dominant.} Pour un quotient de polynômes, factoriser numérateur et dénominateur par la plus haute puissance de $n$.
  \item \textbf{Multiplier par l'expression conjuguée} (si racine carrée).
  \item \textbf{Utiliser les croissances comparées} (section suivante).
\end{enumerate}
\end{methode}

\begin{exemple}[Forme $\frac{\infty}{\infty}$]
Calculons $\ds\lim_{n \to +\infty} \frac{3n^2 - 5n + 1}{2n^2 + n}$.

On factorise par $n^2$ :
\[ \frac{3n^2 - 5n + 1}{2n^2 + n} = \frac{n^2\!\left(3 - \frac{5}{n} + \frac{1}{n^2}\right)}{n^2\!\left(2 + \frac{1}{n}\right)} = \frac{3 - \frac{5}{n} + \frac{1}{n^2}}{2 + \frac{1}{n}} \xrightarrow[n \to +\infty]{} \frac{3}{2} \]
\end{exemple}

\begin{exemple}[Forme $+\infty - \infty$]
Calculons $\ds\lim_{n \to +\infty} \bigl(\sqrt{n+1} - \sqrt{n}\bigr)$.

On multiplie par l'expression conjuguée :
\[ \sqrt{n+1} - \sqrt{n} = \frac{(\sqrt{n+1} - \sqrt{n})(\sqrt{n+1} + \sqrt{n})}{\sqrt{n+1} + \sqrt{n}} = \frac{1}{\sqrt{n+1} + \sqrt{n}} \xrightarrow[n \to +\infty]{} 0 \]
\end{exemple}

% ── IV. Théorèmes de comparaison ──
\section{Théorèmes de comparaison}

\subsection{Théorème de comparaison}

\begin{theoreme}[Comparaison]
Soient $(u_n)$ et $(v_n)$ deux suites telles que, à partir d'un certain rang, $u_n \leqslant v_n$.
\begin{itemize}
  \item Si $\ds\lim_{n \to +\infty} u_n = +\infty$, alors $\ds\lim_{n \to +\infty} v_n = +\infty$.
  \item Si $\ds\lim_{n \to +\infty} v_n = -\infty$, alors $\ds\lim_{n \to +\infty} u_n = -\infty$.
\end{itemize}
\end{theoreme}

\textit{Intuition :} si une suite plus petite tend vers $+\infty$, la plus grande aussi.

\subsection{Théorème des gendarmes (encadrement)}

\begin{theoreme}[Théorème des gendarmes]
Soient $(u_n)$, $(v_n)$ et $(w_n)$ trois suites telles que, à partir d'un certain rang :
\[ v_n \leqslant u_n \leqslant w_n \]
Si $\ds\lim_{n \to +\infty} v_n = \lim_{n \to +\infty} w_n = \ell$, alors $\ds\lim_{n \to +\infty} u_n = \ell$.
\end{theoreme}

\textit{Intuition :} si $(u_n)$ est « coincée » entre deux suites qui convergent vers la même limite, elle converge aussi vers cette limite.

\begin{exemple}
Montrons que $\ds\lim_{n \to +\infty} \frac{\cos n}{n} = 0$.

On sait que $-1 \leqslant \cos n \leqslant 1$ pour tout $n$. En divisant par $n > 0$ :
\[ -\frac{1}{n} \leqslant \frac{\cos n}{n} \leqslant \frac{1}{n} \]
Comme $\ds\lim_{n \to +\infty} \frac{-1}{n} = \lim_{n \to +\infty} \frac{1}{n} = 0$, par le théorème des gendarmes : $\ds\lim_{n \to +\infty} \frac{\cos n}{n} = 0$.
\end{exemple}

\begin{exemple}
Montrons que $\ds\lim_{n \to +\infty} \frac{\sin(n^2)}{n} = 0$.

On encadre : $-1 \leqslant \sin(n^2) \leqslant 1$, donc $\ds -\frac{1}{n} \leqslant \frac{\sin(n^2)}{n} \leqslant \frac{1}{n}$.

Par le théorème des gendarmes, $\ds\lim_{n \to +\infty} \frac{\sin(n^2)}{n} = 0$.
\end{exemple}

% ── V. Croissances comparées ──
\section{Croissances comparées}

\begin{propriete}[Croissances comparées pour les suites]
Pour tout entier $p \geqslant 1$ et tout réel $q > 1$ :
\[ \lim_{n \to +\infty} \frac{n^p}{q^n} = 0 \]
Autrement dit, toute suite géométrique de raison $q > 1$ \textbf{l'emporte} sur toute suite polynomiale.
\end{propriete}

\textit{Interprétation :} la croissance exponentielle est \textbf{toujours} plus rapide que la croissance polynomiale, quelle que soit la puissance.

\begin{exemple}
\begin{itemize}
  \item $\ds\lim_{n \to +\infty} \frac{n^3}{2^n} = 0$ \quad (car $2 > 1$)
  \item $\ds\lim_{n \to +\infty} \frac{n^{100}}{1{,}01^n} = 0$ \quad (même $1{,}01$ finit par l'emporter sur $n^{100}$)
  \item $\ds\lim_{n \to +\infty} n \cdot 0{,}99^n = 0$ \quad (écrire $\ds\frac{n}{(1/0{,}99)^n}$ avec $\ds\frac{1}{0{,}99} > 1$)
\end{itemize}
\end{exemple}

\begin{attention}
Ne pas confondre $2^n$ et $n^2$. La suite $2^n$ croît \textbf{exponentiellement} (doublement à chaque étape), alors que $n^2$ croît \textbf{polynomialement}. Pour $n = 10$ : $n^2 = 100$ mais $2^n = 1024$. Pour $n = 20$ : $n^2 = 400$ mais $2^n \approx 10^6$.
\end{attention}

% ── VI. Suite croissante majorée ──
\section{Théorème de convergence monotone}

\begin{theoreme}[Suite croissante majorée]
Toute suite \textbf{croissante et majorée} est \textbf{convergente}.

De même, toute suite \textbf{décroissante et minorée} est convergente.
\end{theoreme}

\textit{Ce théorème est admis. Il est fondamental car il permet d'affirmer l'existence d'une limite sans la calculer.}

\begin{methode}[Utiliser le théorème de convergence monotone]
Pour montrer qu'une suite $(u_n)$ converge :
\begin{enumerate}
  \item Montrer que $(u_n)$ est \textbf{monotone} (croissante ou décroissante).
  \item Montrer que $(u_n)$ est \textbf{bornée} (majorée si croissante, minorée si décroissante).
  \item Conclure par le théorème : la suite converge vers une limite $\ell$.
  \item \textbf{Chercher la limite $\ell$} : si $u_{n+1} = f(u_n)$, alors $\ell = f(\ell)$ (point fixe).
\end{enumerate}
\end{methode}

\begin{exemple}
Soit $(u_n)$ définie par $u_0 = 1$ et $u_{n+1} = \ds\frac{u_n + 3}{2}$.

\textbf{1. Monotonie :} $u_{n+1} - u_n = \ds\frac{u_n + 3}{2} - u_n = \frac{3 - u_n}{2}$. Si $u_n < 3$, alors $u_{n+1} - u_n > 0$ : la suite est croissante (tant que $u_n < 3$).

\textbf{2. Majoration :} montrons par récurrence que $u_n < 3$ pour tout $n$.
\begin{itemize}
  \item $u_0 = 1 < 3$ : vrai.
  \item Si $u_n < 3$, alors $u_{n+1} = \ds\frac{u_n + 3}{2} < \frac{3 + 3}{2} = 3$ : vrai.
\end{itemize}

\textbf{3. Conclusion :} la suite est croissante et majorée par $3$, donc elle converge vers une limite $\ell$.

\textbf{4. Calcul de $\ell$ :} en passant à la limite dans $u_{n+1} = \ds\frac{u_n + 3}{2}$, on obtient $\ell = \ds\frac{\ell + 3}{2}$, soit $2\ell = \ell + 3$, d'où $\ell = 3$.
\end{exemple}

\begin{exemple}[Suite $u_{n+1} = \sqrt{2 + u_n}$]
Soit $(u_n)$ définie par $u_0 = 0$ et $u_{n+1} = \sqrt{2 + u_n}$.

\textbf{Monotonie :} On montre par récurrence que $0 \leqslant u_n \leqslant u_{n+1}$. La suite est croissante.

\textbf{Majoration :} Montrons que $u_n \leqslant 2$ pour tout $n$. Initialisation : $u_0 = 0 \leqslant 2$. Hérédité : si $u_n \leqslant 2$, alors $u_{n+1} = \sqrt{2 + u_n} \leqslant \sqrt{2 + 2} = 2$.

\textbf{Convergence :} La suite est croissante et majorée par $2$, donc elle converge vers $\ell$.

\textbf{Limite :} $\ell = \sqrt{2 + \ell}$, donc $\ell^2 = 2 + \ell$, soit $\ell^2 - \ell - 2 = 0$. On trouve $\ell = 2$ (la racine $\ell = -1$ est exclue car $\ell \geqslant 0$).
\end{exemple}

% ── VII. Raisonnement par récurrence ──
\section{Raisonnement par récurrence}

\begin{definition}[Principe de récurrence]
Le \textbf{raisonnement par récurrence} est une méthode de démonstration qui permet de prouver qu'une propriété $\mathcal{P}(n)$ est vraie pour tout entier naturel $n \geqslant n_0$.

Il se déroule en \textbf{deux étapes} :
\begin{enumerate}
  \item \textbf{Initialisation :} on vérifie que $\mathcal{P}(n_0)$ est vraie.
  \item \textbf{Hérédité :} on suppose que $\mathcal{P}(n)$ est vraie pour un certain $n \geqslant n_0$ (\textbf{hypothèse de récurrence}), et on montre que $\mathcal{P}(n+1)$ est alors vraie.
\end{enumerate}

\textbf{Conclusion :} $\mathcal{P}(n)$ est vraie pour tout $n \geqslant n_0$.
\end{definition}

\textit{Analogie :} c'est le principe des dominos. Si le premier tombe (initialisation) et si chaque domino fait tomber le suivant (hérédité), alors tous les dominos tombent.

\begin{methode}[Rédiger une récurrence]
\begin{enumerate}
  \item Écrire clairement la propriété $\mathcal{P}(n)$.
  \item \textbf{Initialisation :} vérifier $\mathcal{P}(n_0)$ par un calcul direct.
  \item \textbf{Hérédité :} écrire « Supposons $\mathcal{P}(n)$ vraie pour un $n \geqslant n_0$ fixé. Montrons $\mathcal{P}(n+1)$. »
  \item Utiliser explicitement l'hypothèse de récurrence dans la preuve.
  \item \textbf{Conclusion :} « Par le principe de récurrence, $\mathcal{P}(n)$ est vraie pour tout $n \geqslant n_0$. »
\end{enumerate}
\end{methode}

\begin{exemple}[Somme des entiers]
Montrons que pour tout $n \geqslant 1$ : $\ds\sum_{k=1}^{n} k = \frac{n(n+1)}{2}$.

\textbf{Propriété :} $\mathcal{P}(n)$ : « $1 + 2 + \cdots + n = \ds\frac{n(n+1)}{2}$ ».

\textbf{Initialisation ($n = 1$) :} $\text{Gauche} = 1$ et $\text{Droite} = \ds\frac{1 \times 2}{2} = 1$. $\mathcal{P}(1)$ est vraie.

\textbf{Hérédité :} Supposons $\mathcal{P}(n)$ vraie pour un $n \geqslant 1$ fixé, c'est-à-dire $1 + 2 + \cdots + n = \ds\frac{n(n+1)}{2}$.

Alors :
\begin{align*}
1 + 2 + \cdots + n + (n+1) &= \frac{n(n+1)}{2} + (n+1) \quad \text{(par hypothèse de récurrence)} \\
&= (n+1)\!\left(\frac{n}{2} + 1\right) = (n+1) \cdot \frac{n+2}{2} = \frac{(n+1)(n+2)}{2}
\end{align*}
C'est exactement $\mathcal{P}(n+1)$.

\textbf{Conclusion :} Par le principe de récurrence, $\mathcal{P}(n)$ est vraie pour tout $n \geqslant 1$.
\end{exemple}

\begin{exemple}[Inégalité de Bernoulli]
Montrons que pour tout $n \geqslant 0$ et tout $x \geqslant -1$ : $(1 + x)^n \geqslant 1 + nx$.

\textbf{Initialisation ($n = 0$) :} $(1+x)^0 = 1 \geqslant 1 + 0 = 1$. Vrai.

\textbf{Hérédité :} Supposons $(1+x)^n \geqslant 1 + nx$ pour un certain $n \geqslant 0$.
\begin{align*}
(1+x)^{n+1} &= (1+x)^n \cdot (1+x) \\
&\geqslant (1 + nx)(1+x) \quad \text{(car $1+x \geqslant 0$ par hypothèse)} \\
&= 1 + nx + x + nx^2 = 1 + (n+1)x + nx^2 \\
&\geqslant 1 + (n+1)x \quad \text{(car $nx^2 \geqslant 0$)}
\end{align*}

\textbf{Conclusion :} Par récurrence, $(1+x)^n \geqslant 1 + nx$ pour tout $n \geqslant 0$.
\end{exemple}

\begin{attention}[Erreurs fréquentes en récurrence]
\begin{itemize}
  \item \textbf{Oublier l'initialisation :} sans elle, la preuve est incomplète. La propriété « $n = n + 1$ pour tout $n$ » est héréditaire mais jamais vraie !
  \item \textbf{Montrer $\mathcal{P}(n)$ au lieu de $\mathcal{P}(n+1)$ :} dans l'hérédité, c'est $\mathcal{P}(n+1)$ qu'il faut démontrer, en utilisant $\mathcal{P}(n)$.
  \item \textbf{Ne pas utiliser l'hypothèse de récurrence :} si la preuve de l'hérédité n'utilise pas $\mathcal{P}(n)$, c'est suspect --- la propriété est probablement démontrable directement.
\end{itemize}
\end{attention}

% ── VIII. Applications ──
\section{Applications}

\begin{exemple}[Placement financier]
On place un capital $C_0 = 1000$~€ à un taux annuel de $3\,\%$. Le capital au bout de $n$ années est :
\[ C_n = 1000 \times 1{,}03^n \]
C'est une suite géométrique de raison $q = 1{,}03 > 1$, donc $\ds\lim_{n \to +\infty} C_n = +\infty$ : le capital croît indéfiniment (en théorie).

Au bout de $n = 24$ ans : $C_{24} = 1000 \times 1{,}03^{24} \approx 2032{,}79$~€ (le capital a doublé).
\end{exemple}

\begin{exemple}[Décroissance radioactive]
La masse d'un échantillon radioactif est modélisée par $m_n = m_0 \times 0{,}5^n$ où $n$ est le nombre de demi-vies écoulées. Comme $|0{,}5| < 1$, on a $\ds\lim_{n \to +\infty} m_n = 0$ : la masse tend vers zéro sans jamais l'atteindre.
\end{exemple}

\begin{exemple}[Algorithme de Héron]
Pour calculer $\sqrt{a}$ ($a > 0$), on utilise la suite $u_0 > 0$ et $u_{n+1} = \ds\frac{1}{2}\!\left(u_n + \frac{a}{u_n}\right)$.

On peut montrer que cette suite est décroissante et minorée par $\sqrt{a}$ (à partir du rang 1), donc elle converge. La limite vérifie $\ell = \ds\frac{1}{2}\!\left(\ell + \frac{a}{\ell}\right)$, soit $2\ell^2 = \ell^2 + a$, d'où $\ell = \sqrt{a}$.

Avec $a = 2$ et $u_0 = 2$ : $u_1 = 1{,}5$ ; $u_2 = 1{,}4167$ ; $u_3 = 1{,}41422$ (déjà 5 décimales exactes !).
\end{exemple}

% ── Bilan ──
\section*{Bilan du chapitre}

\begin{retenir}
\begin{itemize}[leftmargin=*]
  \item \textbf{Convergence :} $(u_n)$ converge vers $\ell$ si $u_n$ se rapproche de $\ell$ quand $n \to +\infty$
  \item \textbf{Limites de $q^n$ :} $q > 1 \Rightarrow +\infty$ ; $|q| < 1 \Rightarrow 0$ ; $q = 1 \Rightarrow 1$ ; $q \leqslant -1 \Rightarrow$ pas de limite
  \item \textbf{Opérations :} somme, produit, quotient des limites (avec tableaux)
  \item \textbf{Formes indéterminées :} $+\infty - \infty$, \; $0 \times \infty$, \; $\ds\frac{\infty}{\infty}$, \; $\ds\frac{0}{0}$ --- lever par factorisation, conjuguée ou croissances comparées
  \item \textbf{Gendarmes :} si $v_n \leqslant u_n \leqslant w_n$ et $v_n, w_n \to \ell$, alors $u_n \to \ell$
  \item \textbf{Croissances comparées :} $\ds\frac{n^p}{q^n} \to 0$ pour $q > 1$
  \item \textbf{Convergence monotone :} suite croissante majorée $\Rightarrow$ convergente (admis)
  \item \textbf{Récurrence :} initialisation + hérédité $\Rightarrow$ propriété vraie pour tout $n \geqslant n_0$
\end{itemize}
\end{retenir}


% ─────────────────────────────────────────────────────────────────────────────
% CHAPITRE 6 — LIMITES DES FONCTIONS
% ─────────────────────────────────────────────────────────────────────────────

\chapter{Limites des fonctions}

\section*{Introduction}

\textit{On sait maintenant ce que signifie la limite d'une suite. Mais les suites ne sont que des fonctions définies sur $\N$. Comment étendre la notion de limite aux fonctions définies sur $\R$ ? Que signifie « $f(x)$ tend vers $+\infty$ quand $x$ tend vers $2$ » ou « $f(x)$ se rapproche de $3$ quand $x$ tend vers $+\infty$ » ?}

\medskip

Ce chapitre étend systématiquement les outils du chapitre précédent au cadre continu. Les limites de fonctions permettent de décrire le \textbf{comportement global} d'une fonction (asymptotes) et son \textbf{comportement local} au voisinage de certains points.

\medskip

\textbf{Objectifs du chapitre :}
\begin{itemize}[leftmargin=*]
  \item Déterminer la limite d'une fonction en $+\infty$, $-\infty$ ou en un point
  \item Interpréter graphiquement les limites (asymptotes)
  \item Utiliser les opérations sur les limites et lever les formes indéterminées
  \item Appliquer les théorèmes de comparaison (gendarmes, croissances comparées)
  \item Déterminer les asymptotes horizontales, verticales et obliques
\end{itemize}

% ── I. Limite en l'infini ──
\section{Limite d'une fonction en l'infini}

\subsection{Limite finie en $+\infty$ --- Asymptote horizontale}

\begin{definition}[Limite finie en $+\infty$]
On dit que $f(x)$ a pour limite $\ell$ quand $x$ tend vers $+\infty$ si $f(x)$ peut être rendu aussi proche que l'on veut de $\ell$ pour $x$ suffisamment grand.

On note : $\ds\lim_{x \to +\infty} f(x) = \ell$.
\end{definition}

\begin{propriete}[Asymptote horizontale]
Si $\ds\lim_{x \to +\infty} f(x) = \ell$, alors la droite $y = \ell$ est \textbf{asymptote horizontale} à la courbe de $f$ en $+\infty$. De même en $-\infty$.
\end{propriete}

\begin{exemple}
$\ds\lim_{x \to +\infty} \frac{1}{x} = 0$ : la droite $y = 0$ (axe des abscisses) est asymptote horizontale.

\medskip

$\ds\lim_{x \to +\infty} \frac{2x + 1}{x - 3} = \lim_{x \to +\infty} \frac{2 + \frac{1}{x}}{1 - \frac{3}{x}} = 2$ : la droite $y = 2$ est asymptote horizontale.
\end{exemple}

\subsection{Limite infinie en l'infini}

\begin{definition}[Limite infinie en $+\infty$]
On dit que $f(x)$ a pour limite $+\infty$ quand $x$ tend vers $+\infty$ si $f(x)$ devient aussi grand que l'on veut pour $x$ suffisamment grand.

On note : $\ds\lim_{x \to +\infty} f(x) = +\infty$.
\end{definition}

\begin{exemple}
$\ds\lim_{x \to +\infty} x^2 = +\infty$, \quad $\ds\lim_{x \to +\infty} e^x = +\infty$, \quad $\ds\lim_{x \to -\infty} x^3 = -\infty$.
\end{exemple}

% ── II. Limite en un point ──
\section{Limite d'une fonction en un point}

\subsection{Limite finie en un point}

\begin{definition}[Limite en un point]
On dit que $f(x)$ a pour limite $\ell$ quand $x$ tend vers $a$ si $f(x)$ se rapproche de $\ell$ lorsque $x$ se rapproche de $a$ (sans nécessairement être définie en $a$).

On note : $\ds\lim_{x \to a} f(x) = \ell$.
\end{definition}

\begin{exemple}
$\ds\lim_{x \to 0} \frac{\sin x}{x} = 1$ \quad (limite fondamentale, $f$ n'est pas définie en $0$ mais la limite existe).
\end{exemple}

\subsection{Limite infinie en un point --- Asymptote verticale}

\begin{definition}[Asymptote verticale]
Si $\ds\lim_{x \to a} f(x) = +\infty$ ou $\ds\lim_{x \to a} f(x) = -\infty$, alors la droite $x = a$ est \textbf{asymptote verticale} à la courbe de $f$.
\end{definition}

\begin{exemple}
$\ds\lim_{x \to 0^+} \frac{1}{x} = +\infty$ et $\ds\lim_{x \to 0^-} \frac{1}{x} = -\infty$ : la droite $x = 0$ est asymptote verticale à $\ds\frac{1}{x}$.
\end{exemple}

\subsection{Limites à gauche et à droite}

\begin{definition}[Limites latérales]
On peut étudier la limite de $f$ en $a$ \textbf{par la gauche} (notée $\ds\lim_{x \to a^-} f(x)$) et \textbf{par la droite} (notée $\ds\lim_{x \to a^+} f(x)$).

$f$ admet une limite en $a$ si et seulement si les limites à gauche et à droite existent et sont égales.
\end{definition}

\begin{attention}[Limites latérales différentes]
Quand les limites à gauche et à droite diffèrent, la fonction n'a pas de limite en ce point.

Exemple : la fonction partie entière $E(x)$ en tout entier $n$ : $\ds\lim_{x \to n^-} E(x) = n-1$ et $\ds\lim_{x \to n^+} E(x) = n$.
\end{attention}

% ── III. Opérations sur les limites ──
\section{Opérations sur les limites de fonctions}

Les théorèmes d'opérations sont les mêmes que pour les suites.

\begin{propriete}[Opérations]
Si $\ds\lim_{x \to a} f(x) = \ell$ et $\ds\lim_{x \to a} g(x) = \ell'$ (avec $a \in \R$ ou $a = \pm\infty$), alors :
\begin{itemize}
  \item $\ds\lim_{x \to a} \bigl(f(x) + g(x)\bigr) = \ell + \ell'$ \quad (sauf FI $+\infty - \infty$)
  \item $\ds\lim_{x \to a} \bigl(f(x) \cdot g(x)\bigr) = \ell \cdot \ell'$ \quad (sauf FI $0 \times \infty$)
  \item $\ds\lim_{x \to a} \frac{f(x)}{g(x)} = \frac{\ell}{\ell'}$ si $\ell' \neq 0$ \quad (sauf FI $\ds\frac{\infty}{\infty}$ ou $\ds\frac{0}{0}$)
\end{itemize}
\end{propriete}

\begin{methode}[Lever une forme indéterminée pour les fonctions]
\begin{enumerate}
  \item \textbf{Fonctions rationnelles en $\pm\infty$ :} factoriser par la plus haute puissance de $x$.
  \item \textbf{Forme $\frac{0}{0}$ en un point :} factoriser, simplifier, ou utiliser la quantité conjuguée.
  \item \textbf{Avec exponentielle ou logarithme :} utiliser les croissances comparées.
\end{enumerate}
\end{methode}

\begin{exemple}[Forme $\frac{0}{0}$]
$\ds\lim_{x \to 1} \frac{x^2 - 1}{x - 1} = \lim_{x \to 1} \frac{(x-1)(x+1)}{x-1} = \lim_{x \to 1} (x+1) = 2$.
\end{exemple}

\begin{exemple}[Forme $\frac{\infty}{\infty}$]
$\ds\lim_{x \to +\infty} \frac{3x^3 - x}{x^3 + 2x^2} = \lim_{x \to +\infty} \frac{3 - \frac{1}{x^2}}{1 + \frac{2}{x}} = 3$.
\end{exemple}

% ── IV. Théorèmes de comparaison ──
\section{Théorèmes de comparaison pour les fonctions}

\begin{theoreme}[Comparaison]
Si $f(x) \leqslant g(x)$ au voisinage de $+\infty$ et si $\ds\lim_{x \to +\infty} f(x) = +\infty$, alors $\ds\lim_{x \to +\infty} g(x) = +\infty$.
\end{theoreme}

\begin{theoreme}[Théorème des gendarmes pour les fonctions]
Si, au voisinage de $a$ (fini ou infini) :
\[ g(x) \leqslant f(x) \leqslant h(x) \]
et si $\ds\lim_{x \to a} g(x) = \lim_{x \to a} h(x) = \ell$, alors $\ds\lim_{x \to a} f(x) = \ell$.
\end{theoreme}

\begin{exemple}
Montrons que $\ds\lim_{x \to +\infty} \frac{\sin x}{x} = 0$.

Pour tout $x > 0$ : $\ds -\frac{1}{x} \leqslant \frac{\sin x}{x} \leqslant \frac{1}{x}$. Les deux bornes tendent vers $0$, donc par les gendarmes : $\ds\lim_{x \to +\infty} \frac{\sin x}{x} = 0$.
\end{exemple}

% ── V. Croissances comparées ──
\section{Croissances comparées}

\begin{propriete}[Croissances comparées avec l'exponentielle]
Pour tout entier $n \geqslant 1$ :
\[ \lim_{x \to +\infty} \frac{x^n}{e^x} = 0 \qquad \text{et} \qquad \lim_{x \to -\infty} x^n \, e^x = 0 \]
L'exponentielle l'emporte sur toute puissance de $x$ en $+\infty$.
\end{propriete}

\begin{propriete}[Croissances comparées avec le logarithme]
Pour tout entier $n \geqslant 1$ :
\[ \lim_{x \to +\infty} \frac{\ln x}{x^n} = 0 \qquad \text{et} \qquad \lim_{x \to 0^+} x^n \ln x = 0 \]
Toute puissance de $x$ l'emporte sur le logarithme.
\end{propriete}

\textit{Résumé de la hiérarchie :} $\ln x \ll x^n \ll e^x$ quand $x \to +\infty$.

\begin{exemple}
\begin{itemize}
  \item $\ds\lim_{x \to +\infty} \frac{x^5}{e^x} = 0$
  \item $\ds\lim_{x \to +\infty} \frac{\ln x}{\sqrt{x}} = \lim_{x \to +\infty} \frac{\ln x}{x^{1/2}} = 0$
  \item $\ds\lim_{x \to 0^+} x \ln x = 0$ \quad (le $x$ l'emporte et « ramène » vers 0)
  \item $\ds\lim_{x \to +\infty} x^2 e^{-x} = \lim_{x \to +\infty} \frac{x^2}{e^x} = 0$
\end{itemize}
\end{exemple}

\begin{attention}[FI typiques avec $e^x$ et $\ln x$]
Les formes $x^n \cdot e^{-x}$ en $+\infty$ ou $x^n \cdot \ln x$ en $0^+$ ressemblent à des formes indéterminées $\infty \times 0$ mais se résolvent systématiquement par les croissances comparées : elles tendent toutes vers $0$.
\end{attention}

% ── VI. Composition de limites ──
\section{Limite d'une composition}

\begin{propriete}[Composition de limites]
Si $\ds\lim_{x \to a} g(x) = b$ et $\ds\lim_{X \to b} f(X) = \ell$, alors :
\[ \lim_{x \to a} f\bigl(g(x)\bigr) = \ell \]
\end{propriete}

\begin{exemple}
\begin{itemize}
  \item $\ds\lim_{x \to +\infty} e^{-x^2} = \lim_{X \to -\infty} e^X = 0$ \quad (en posant $X = -x^2 \to -\infty$)
  \item $\ds\lim_{x \to +\infty} \ln(x^2 + 1) = +\infty$ \quad (car $x^2 + 1 \to +\infty$ et $\ln X \to +\infty$)
  \item $\ds\lim_{x \to 0^+} \sqrt{\frac{1}{x}} = +\infty$ \quad (car $\frac{1}{x} \to +\infty$ et $\sqrt{X} \to +\infty$)
\end{itemize}
\end{exemple}

% ── VII. Asymptotes --- synthèse ──
\section{Asymptotes --- synthèse}

\subsection{Récapitulatif}

\begin{propriete}[Les trois types d'asymptotes]
Soit $\mathcal{C}_f$ la courbe de $f$.
\begin{enumerate}
  \item \textbf{Asymptote horizontale} $y = \ell$ : si $\ds\lim_{x \to +\infty} f(x) = \ell$ ou $\ds\lim_{x \to -\infty} f(x) = \ell$.
  \item \textbf{Asymptote verticale} $x = a$ : si $\ds\lim_{x \to a} f(x) = \pm\infty$.
  \item \textbf{Asymptote oblique} $y = ax + b$ : si $\ds\lim_{x \to +\infty} \bigl[f(x) - (ax+b)\bigr] = 0$.
\end{enumerate}
\end{propriete}

\subsection{Détermination d'une asymptote oblique}

\begin{methode}[Trouver une asymptote oblique]
Pour chercher une asymptote oblique $y = ax + b$ en $+\infty$ :
\begin{enumerate}
  \item Calculer $a = \ds\lim_{x \to +\infty} \frac{f(x)}{x}$. Si cette limite est finie et non nulle, une asymptote oblique est possible.
  \item Calculer $b = \ds\lim_{x \to +\infty} \bigl[f(x) - ax\bigr]$. Si cette limite est finie, $y = ax + b$ est asymptote oblique.
\end{enumerate}
\end{methode}

\begin{exemple}
Soit $f(x) = \ds\frac{x^2 + 1}{x} = x + \frac{1}{x}$ pour $x \neq 0$.

$f(x) - x = \ds\frac{1}{x} \xrightarrow[x \to +\infty]{} 0$. Donc $y = x$ est asymptote oblique en $+\infty$ (et aussi en $-\infty$).

Par ailleurs, $\ds\lim_{x \to 0^+} f(x) = +\infty$ et $\ds\lim_{x \to 0^-} f(x) = -\infty$ : la droite $x = 0$ est asymptote verticale.
\end{exemple}

\begin{center}
\begin{tikzpicture}
\begin{axis}[
  width=10cm, height=8cm,
  xlabel={$x$}, ylabel={$y$},
  xmin=-6, xmax=6, ymin=-8, ymax=8,
  xtick={-5,...,5},
  grid=major, grid style={gray!15},
  axis lines=middle,
  samples=200,
  title={$f(x) = x + \frac{1}{x}$ : AV en $x=0$, AO $y = x$},
  title style={font=\sffamily\bfseries}
]
\addplot[indigo, thick, domain=-6:-0.15] {x + 1/x};
\addplot[indigo, thick, domain=0.15:6] {x + 1/x};
\addplot[rougeerr, dashed, thick, domain=-6:6] {x};
\draw[rougeerr, dashed, thick] (axis cs:0,-8) -- (axis cs:0,8);
\node[rougeerr, anchor=south west] at (axis cs:3.5,4.5) {$y = x$};
\node[rougeerr, anchor=west] at (axis cs:0.3,-6) {$x = 0$};
\end{axis}
\end{tikzpicture}
\captionof{figure}{La courbe de $f(x) = x + \frac{1}{x}$ avec ses asymptotes.}
\end{center}

% ── Bilan ──
\section*{Bilan du chapitre}

\begin{retenir}
\begin{itemize}[leftmargin=*]
  \item \textbf{Limite en $+\infty$ :} $\ds\lim_{x \to +\infty} f(x) = \ell$ signifie $f(x)$ se rapproche de $\ell$ ; $y = \ell$ est AH
  \item \textbf{Limite en un point :} $\ds\lim_{x \to a} f(x) = \pm\infty$ donne une AV $x = a$
  \item \textbf{Opérations :} mêmes règles que les suites, avec mêmes FI
  \item \textbf{Gendarmes :} $g \leqslant f \leqslant h$ et $g, h \to \ell$ $\Rightarrow$ $f \to \ell$
  \item \textbf{Croissances comparées :} $\ln x \ll x^n \ll e^x$ en $+\infty$
  \item \textbf{Composition :} $\ds\lim f\bigl(g(x)\bigr) = f\!\bigl(\lim g(x)\bigr)$ (sous conditions)
  \item \textbf{AO $y = ax + b$ :} si $\ds\lim \bigl[f(x) - (ax+b)\bigr] = 0$
\end{itemize}
\end{retenir}


% ─────────────────────────────────────────────────────────────────────────────
% CHAPITRE 7 — DÉRIVATION : COMPOSÉE ET CONVEXITÉ
% ─────────────────────────────────────────────────────────────────────────────

\chapter{Dérivation : composée et convexité}

\section*{Introduction}

\textit{En Première, nous avons appris à dériver des fonctions simples (polynômes, quotients, racines). Mais comment dériver $e^{x^2}$, $\ln(\cos x)$ ou $\sqrt{2x+1}$ ? Et comment savoir si une courbe « s'arrondit vers le haut » ou « vers le bas » ?}

\medskip

Ce chapitre introduit deux outils puissants : la \textbf{dérivée d'une composition}, qui permet de dériver presque toute fonction, et la notion de \textbf{convexité}, qui affine l'étude des fonctions grâce à la dérivée seconde.

\medskip

\textbf{Objectifs du chapitre :}
\begin{itemize}[leftmargin=*]
  \item Dériver une fonction composée $f = v \circ u$
  \item Connaître les formules de dérivation des composées usuelles
  \item Définir et interpréter la dérivée seconde
  \item Caractériser la convexité et la concavité par $f''$
  \item Identifier les points d'inflexion
  \item Utiliser la convexité pour établir des inégalités
\end{itemize}

% ── I. Dérivée d'une fonction composée ──
\section{Dérivée d'une fonction composée}

\subsection{Composition de fonctions}

\begin{definition}[Fonction composée]
Soient $u$ et $v$ deux fonctions. La \textbf{composée} de $v$ par $u$, notée $v \circ u$, est la fonction définie par :
\[ (v \circ u)(x) = v\bigl(u(x)\bigr) \]
On lit « $v$ rond $u$ » ou « $v$ composée avec $u$ ». On applique d'abord $u$, puis $v$ au résultat.
\end{definition}

\begin{exemple}
Si $u(x) = x^2 + 1$ et $v(X) = \sqrt{X}$, alors $(v \circ u)(x) = \sqrt{x^2 + 1}$.

Si $u(x) = 2x + 3$ et $v(X) = e^X$, alors $(v \circ u)(x) = e^{2x+3}$.
\end{exemple}

\subsection{Formule de dérivation de la composée}

\begin{theoreme}[Dérivée d'une composée]
Si $u$ est dérivable en $x$ et $v$ est dérivable en $u(x)$, alors $v \circ u$ est dérivable en $x$ et :
\[ (v \circ u)'(x) = u'(x) \times v'\bigl(u(x)\bigr) \]
\end{theoreme}

\textit{Formulation mnémotechnique :} « la dérivée de la composée est la dérivée de l'intérieur multipliée par la dérivée de l'extérieur évaluée à l'intérieur ».

\begin{attention}[L'ordre compte]
La formule $(v \circ u)' = u' \times (v' \circ u)$ n'est pas la même chose que $(u \circ v)' = v' \times (u' \circ v)$. L'ordre de la composition est crucial.
\end{attention}

\subsection{Table des dérivées composées usuelles}

\begin{propriete}[Formules de dérivation des composées]
Dans le tableau ci-dessous, $u$ désigne une fonction dérivable.

\begin{center}
\renewcommand{\arraystretch}{2}
\begin{tabular}{ccc}
\toprule
\textbf{Fonction} & \textbf{Dérivée} & \textbf{Conditions} \\
\midrule
$e^{u}$ & $u' \, e^{u}$ & \\
$\ln u$ & $\ds\frac{u'}{u}$ & $u > 0$ \\
$u^n$ \; ($n \in \Z$) & $n \, u' \, u^{n-1}$ & $u \neq 0$ si $n < 0$ \\
$\sqrt{u}$ & $\ds\frac{u'}{2\sqrt{u}}$ & $u > 0$ \\
$\cos u$ & $-u' \sin u$ & \\
$\sin u$ & $u' \cos u$ & \\
\bottomrule
\end{tabular}
\end{center}
\end{propriete}

\begin{methode}[Dériver une composée]
\begin{enumerate}
  \item \textbf{Identifier} la fonction « extérieure » $v$ et la fonction « intérieure » $u$.
  \item \textbf{Calculer} $u'(x)$.
  \item \textbf{Appliquer} la formule : dérivée de l'extérieur en $u(x)$, multipliée par $u'(x)$.
\end{enumerate}
\end{methode}

\begin{exemple}[Dérivation de $e^{u}$]
Soit $f(x) = e^{3x^2 - x}$. Ici $u(x) = 3x^2 - x$ et $v(X) = e^X$.

$u'(x) = 6x - 1$, donc $f'(x) = (6x-1)\,e^{3x^2 - x}$.
\end{exemple}

\begin{exemple}[Dérivation de $\ln u$]
Soit $f(x) = \ln(x^2 + 1)$. Ici $u(x) = x^2 + 1$ (toujours $> 0$).

$u'(x) = 2x$, donc $f'(x) = \ds\frac{2x}{x^2 + 1}$.
\end{exemple}

\begin{exemple}[Dérivation de $\sqrt{u}$]
Soit $f(x) = \sqrt{2x + 5}$ pour $x > -\ds\frac{5}{2}$. Ici $u(x) = 2x + 5$.

$u'(x) = 2$, donc $f'(x) = \ds\frac{2}{2\sqrt{2x+5}} = \frac{1}{\sqrt{2x+5}}$.
\end{exemple}

\begin{exemple}[Double composition]
Soit $f(x) = e^{\sqrt{x}}$ pour $x > 0$. On pose $g(x) = \sqrt{x}$, donc $f(x) = e^{g(x)}$.

$g'(x) = \ds\frac{1}{2\sqrt{x}}$, donc $f'(x) = g'(x) \cdot e^{g(x)} = \ds\frac{e^{\sqrt{x}}}{2\sqrt{x}}$.
\end{exemple}

% ── II. Dérivée seconde ──
\section{Dérivée seconde}

\begin{definition}[Dérivée seconde]
Si $f$ est dérivable et si $f'$ est elle-même dérivable, on appelle \textbf{dérivée seconde} de $f$ la fonction $f'' = (f')'$.

On note aussi $f''(x)$ ou $\ds\frac{\mathrm{d}^2 f}{\mathrm{d}x^2}$.
\end{definition}

\begin{exemple}
Soit $f(x) = x^3 - 6x^2 + 9x + 1$.
\begin{itemize}
  \item $f'(x) = 3x^2 - 12x + 9$
  \item $f''(x) = 6x - 12$
\end{itemize}
\end{exemple}

\begin{exemple}
Soit $f(x) = e^{2x}$.
\begin{itemize}
  \item $f'(x) = 2e^{2x}$
  \item $f''(x) = 4e^{2x}$
\end{itemize}
\end{exemple}

% ── III. Convexité ──
\section{Convexité et concavité}

\subsection{Définitions}

\begin{definition}[Fonction convexe, fonction concave]
Soit $f$ une fonction définie sur un intervalle $I$.
\begin{itemize}
  \item $f$ est \textbf{convexe} sur $I$ si, pour tous $a, b \in I$ et tout $t \in [0,1]$ :
  \[ f\bigl(ta + (1-t)b\bigr) \leqslant t\,f(a) + (1-t)\,f(b) \]
  \item $f$ est \textbf{concave} sur $I$ si l'inégalité est dans l'autre sens ($\geqslant$).
\end{itemize}
\end{definition}

\textit{Interprétation géométrique :}
\begin{itemize}
  \item $f$ convexe $\Leftrightarrow$ la courbe est \textbf{en dessous} de chacune de ses cordes (segments joignant deux points de la courbe). La courbe « s'arrondit vers le haut ».
  \item $f$ concave $\Leftrightarrow$ la courbe est \textbf{au-dessus} de chacune de ses cordes. La courbe « s'arrondit vers le bas ».
\end{itemize}

\begin{center}
\begin{tikzpicture}
% Convexe
\begin{scope}
  \draw[->] (-0.3,0) -- (4.5,0);
  \draw[->] (0,-0.3) -- (0,4);
  \draw[indigo, very thick, domain=0.3:4, samples=100] plot (\x, {0.3*(\x-2)^2 + 0.5});
  \draw[rougeerr, thick, dashed] (0.5,{0.3*(0.5-2)^2+0.5}) -- (3.8,{0.3*(3.8-2)^2+0.5});
  \fill[indigo] (0.5,{0.3*(0.5-2)^2+0.5}) circle (2pt);
  \fill[indigo] (3.8,{0.3*(3.8-2)^2+0.5}) circle (2pt);
  \node[below] at (2.2,-0.3) {\textbf{Convexe}};
  \node[rougeerr, font=\small] at (2.8,2.5) {corde};
\end{scope}
% Concave
\begin{scope}[xshift=6cm]
  \draw[->] (-0.3,0) -- (4.5,0);
  \draw[->] (0,-0.3) -- (0,4);
  \draw[indigo, very thick, domain=0.3:4, samples=100] plot (\x, {-0.3*(\x-2)^2 + 3.5});
  \draw[rougeerr, thick, dashed] (0.5,{-0.3*(0.5-2)^2+3.5}) -- (3.8,{-0.3*(3.8-2)^2+3.5});
  \fill[indigo] (0.5,{-0.3*(0.5-2)^2+3.5}) circle (2pt);
  \fill[indigo] (3.8,{-0.3*(3.8-2)^2+3.5}) circle (2pt);
  \node[below] at (2.2,-0.3) {\textbf{Concave}};
  \node[rougeerr, font=\small] at (2.8,1.8) {corde};
\end{scope}
\end{tikzpicture}
\captionof{figure}{Courbe convexe (sous la corde) et concave (au-dessus de la corde).}
\end{center}

\subsection{Caractérisation par $f''$}

\begin{theoreme}[Caractérisation de la convexité par la dérivée seconde]
Soit $f$ deux fois dérivable sur un intervalle $I$.
\begin{itemize}
  \item $f$ est \textbf{convexe} sur $I$ $\Leftrightarrow$ $f''(x) \geqslant 0$ pour tout $x \in I$.
  \item $f$ est \textbf{concave} sur $I$ $\Leftrightarrow$ $f''(x) \leqslant 0$ pour tout $x \in I$.
\end{itemize}
\end{theoreme}

\textit{Autrement dit :} convexe $\Leftrightarrow$ $f'$ croissante ; concave $\Leftrightarrow$ $f'$ décroissante.

\begin{exemple}
\begin{itemize}
  \item $f(x) = x^2$ : $f''(x) = 2 > 0$ partout, donc $f$ est convexe sur $\R$.
  \item $f(x) = e^x$ : $f''(x) = e^x > 0$ partout, donc $f$ est convexe sur $\R$.
  \item $f(x) = \ln x$ : $f''(x) = -\ds\frac{1}{x^2} < 0$ pour $x > 0$, donc $f$ est concave sur $]0, +\infty[$.
  \item $f(x) = \sqrt{x}$ : $f''(x) = -\ds\frac{1}{4x^{3/2}} < 0$ pour $x > 0$, donc $f$ est concave sur $]0, +\infty[$.
\end{itemize}
\end{exemple}

\subsection{Position de la courbe par rapport aux tangentes}

\begin{propriete}[Tangente et convexité]
\begin{itemize}
  \item Si $f$ est \textbf{convexe}, sa courbe est \textbf{au-dessus} de chacune de ses tangentes.
  \item Si $f$ est \textbf{concave}, sa courbe est \textbf{en dessous} de chacune de ses tangentes.
\end{itemize}
\end{propriete}

\begin{demonstration}
Soit $f$ convexe et deux fois dérivable. La tangente en $a$ a pour équation $T_a(x) = f(a) + f'(a)(x - a)$. Posons $g(x) = f(x) - T_a(x)$. Alors $g(a) = 0$ et $g'(a) = 0$. De plus, $g''(x) = f''(x) \geqslant 0$, donc $g'$ est croissante. Comme $g'(a) = 0$, on a $g'(x) \leqslant 0$ pour $x \leqslant a$ et $g'(x) \geqslant 0$ pour $x \geqslant a$. Donc $g$ admet un minimum en $a$ et $g(x) \geqslant g(a) = 0$ pour tout $x$. Ainsi $f(x) \geqslant T_a(x)$. \qed
\end{demonstration}

% ── IV. Point d'inflexion ──
\section{Point d'inflexion}

\begin{definition}[Point d'inflexion]
Un point $A\bigl(a, f(a)\bigr)$ est un \textbf{point d'inflexion} de la courbe de $f$ si la convexité change de sens en $a$ : la courbe passe de convexe à concave, ou de concave à convexe.
\end{definition}

\begin{propriete}[Caractérisation du point d'inflexion]
Si $f$ est deux fois dérivable, alors $A\bigl(a, f(a)\bigr)$ est un point d'inflexion si et seulement si $f''$ \textbf{change de signe} en $a$.

Condition nécessaire : $f''(a) = 0$ (mais ce n'est pas suffisant).
\end{propriete}

\begin{attention}[$f''(a) = 0$ ne suffit pas]
$f(x) = x^4$ vérifie $f''(0) = 0$ mais $f''(x) = 12x^2 \geqslant 0$ partout : pas de changement de signe, donc pas de point d'inflexion en $0$.
\end{attention}

\begin{exemple}
Soit $f(x) = x^3 - 3x^2 + 2x$.
\begin{itemize}
  \item $f'(x) = 3x^2 - 6x + 2$
  \item $f''(x) = 6x - 6 = 6(x - 1)$
\end{itemize}
$f''(x) = 0 \Leftrightarrow x = 1$. De plus, $f''(x) < 0$ pour $x < 1$ et $f''(x) > 0$ pour $x > 1$ : changement de signe en $x = 1$.

Le point $\bigl(1, f(1)\bigr) = (1, 0)$ est un point d'inflexion. En ce point, la tangente \textbf{traverse} la courbe.
\end{exemple}

\begin{propriete}[Tangente au point d'inflexion]
En un point d'inflexion, la tangente \textbf{traverse} la courbe : la courbe passe d'un côté à l'autre de la tangente.
\end{propriete}

% ── V. Convexité et inégalités ──
\section{Convexité et inégalités}

La convexité fournit des inégalités remarquables fondamentales en analyse et en probabilités.

\begin{propriete}[Inégalité $e^x \geqslant 1 + x$]
Pour tout $x \in \R$ :
\[ e^x \geqslant 1 + x \]
avec égalité si et seulement si $x = 0$.
\end{propriete}

\begin{demonstration}
La fonction $f(x) = e^x$ est convexe sur $\R$ (car $f''(x) = e^x > 0$). Sa tangente en $0$ est $T_0(x) = 1 + x$. Par la propriété de convexité, la courbe est au-dessus de sa tangente, donc $e^x \geqslant 1 + x$ pour tout $x$, avec égalité uniquement au point de tangence $x = 0$. \qed
\end{demonstration}

\begin{propriete}[Inégalité $\ln x \leqslant x - 1$]
Pour tout $x > 0$ :
\[ \ln x \leqslant x - 1 \]
avec égalité si et seulement si $x = 1$.
\end{propriete}

\begin{demonstration}
La fonction $\ln$ est concave sur $]0, +\infty[$ (car $(\ln x)'' = -\ds\frac{1}{x^2} < 0$). Sa tangente en $1$ est $T_1(x) = x - 1$. Par concavité, la courbe est en dessous de sa tangente, donc $\ln x \leqslant x - 1$. \qed
\end{demonstration}

\begin{propriete}[Inégalité arithmético-géométrique]
Pour tous réels $a, b \geqslant 0$ :
\[ \sqrt{ab} \leqslant \frac{a + b}{2} \]
La moyenne géométrique est inférieure ou égale à la moyenne arithmétique.
\end{propriete}

\begin{demonstration}[par la concavité de $\sqrt{x}$]
La fonction $g(x) = \sqrt{x}$ est concave sur $[0, +\infty[$ (car $g''(x) = -\ds\frac{1}{4x^{3/2}} < 0$ pour $x > 0$). Par concavité :
\[ g\!\left(\frac{a+b}{2}\right) \geqslant \frac{g(a) + g(b)}{2} \]
soit $\sqrt{\ds\frac{a+b}{2}} \geqslant \ds\frac{\sqrt{a} + \sqrt{b}}{2}$.

On peut aussi démontrer directement : $\ds\frac{a+b}{2} - \sqrt{ab} = \frac{(\sqrt{a} - \sqrt{b})^2}{2} \geqslant 0$. \qed
\end{demonstration}

% ── VI. Étude complète ──
\section{Étude complète d'une fonction}

\begin{methode}[Mener une étude complète]
\begin{enumerate}
  \item \textbf{Ensemble de définition} et éventuelles symétries.
  \item \textbf{Limites} aux bornes du domaine (asymptotes).
  \item \textbf{Dérivée $f'$} : signe, tableau de variations.
  \item \textbf{Dérivée seconde $f''$} : signe, convexité, points d'inflexion.
  \item \textbf{Tableau complet} et \textbf{courbe}.
\end{enumerate}
\end{methode}

\begin{exemple}[Étude de $f(x) = x\,e^{-x}$]
\textbf{Domaine :} $\R$.

\textbf{Limites :}
\begin{itemize}
  \item En $+\infty$ : $f(x) = \ds\frac{x}{e^x} \to 0$ par croissances comparées. AH : $y = 0$.
  \item En $-\infty$ : $x \to -\infty$ et $e^{-x} \to +\infty$, donc $f(x) \to -\infty$.
\end{itemize}

\textbf{Dérivée :} $f'(x) = e^{-x} + x \cdot (-e^{-x}) = (1-x)\,e^{-x}$.

$f'(x) = 0 \Leftrightarrow x = 1$ (car $e^{-x} > 0$ toujours). $f'(x) > 0$ pour $x < 1$, $f'(x) < 0$ pour $x > 1$.

Maximum en $x = 1$ : $f(1) = e^{-1} = \ds\frac{1}{e}$.

\textbf{Dérivée seconde :} $f'(x) = (1-x)e^{-x}$. Par la formule du produit :
\[ f''(x) = (-1) \cdot e^{-x} + (1-x) \cdot (-e^{-x}) = e^{-x}\bigl[-1 -(1-x)\bigr] = (x-2)\,e^{-x} \]

$f''(x) = 0 \Leftrightarrow x = 2$. $f''(x) < 0$ pour $x < 2$ (concave), $f''(x) > 0$ pour $x > 2$ (convexe).

Point d'inflexion en $\bigl(2, 2e^{-2}\bigr) \approx (2 ;\, 0{,}27)$.
\end{exemple}

% ── Bilan ──
\section*{Bilan du chapitre}

\begin{retenir}
\begin{itemize}[leftmargin=*]
  \item \textbf{Composée :} $(v \circ u)'(x) = u'(x) \times v'\bigl(u(x)\bigr)$
  \item \textbf{Formules clés :} $(e^u)' = u'e^u$, \; $(\ln u)' = \ds\frac{u'}{u}$, \; $(\sqrt{u}\,)' = \ds\frac{u'}{2\sqrt{u}}$, \; $(u^n)' = nu'u^{n-1}$
  \item \textbf{Dérivée seconde :} $f'' = (f')'$
  \item \textbf{Convexité :} $f'' \geqslant 0 \Leftrightarrow f$ convexe $\Leftrightarrow$ courbe sous les cordes, au-dessus des tangentes
  \item \textbf{Concavité :} $f'' \leqslant 0 \Leftrightarrow f$ concave $\Leftrightarrow$ courbe au-dessus des cordes, sous les tangentes
  \item \textbf{Point d'inflexion :} changement de signe de $f''$ ; la tangente traverse la courbe
  \item \textbf{Inégalités :} $e^x \geqslant 1 + x$, \; $\ln x \leqslant x - 1$, \; $\sqrt{ab} \leqslant \ds\frac{a+b}{2}$
\end{itemize}
\end{retenir}


% ─────────────────────────────────────────────────────────────────────────────
% CHAPITRE 8 — CONTINUITÉ DES FONCTIONS
% ─────────────────────────────────────────────────────────────────────────────

\chapter{Continuité des fonctions}

\section*{Introduction}

\textit{Peut-on tracer le graphe d'une fonction « sans lever le crayon » ? Si une température passe de $15\,°$C à $25\,°$C au cours d'une journée, est-elle forcément passée par $20\,°$C ? Si $f(1) < 0$ et $f(3) > 0$, existe-t-il un $c$ entre $1$ et $3$ tel que $f(c) = 0$ ?}

\medskip

La notion de \textbf{continuité} formalise l'idée intuitive d'une fonction « sans saut ». C'est un concept fondamental de l'analyse, qui conduit au \textbf{théorème des valeurs intermédiaires} (TVI), l'un des résultats les plus utilisés pour prouver l'existence de solutions d'équations.

\medskip

\textbf{Objectifs du chapitre :}
\begin{itemize}[leftmargin=*]
  \item Définir la continuité en un point et sur un intervalle
  \item Connaître le lien entre dérivabilité et continuité
  \item Énoncer et appliquer le théorème des valeurs intermédiaires (TVI)
  \item Utiliser le corollaire du TVI pour prouver l'existence d'une solution unique
  \item Mettre en œuvre la méthode de dichotomie
  \item Étudier des suites définies par $u_{n+1} = f(u_n)$
\end{itemize}

\begin{histoire}[La continuité : de l'intuition à la rigueur]
Le concept de continuité a une histoire tumultueuse. Au XVIII\ieme{} siècle, Euler et les analystes travaillent avec des fonctions « continues » de manière informelle. C'est \textbf{Bernard Bolzano} (1781--1848), mathématicien et philosophe tchèque, qui démontre en 1817 le premier résultat rigoureux de type TVI, dans un article resté longtemps méconnu. Indépendamment, \textbf{Cauchy} formalise la continuité dans son \textit{Cours d'analyse} (1821). Enfin, \textbf{Karl Weierstrass} (1815--1897) parachève la rigueur en donnant la définition $\varepsilon$-$\delta$ que l'on utilise encore aujourd'hui. Il construit aussi la première fonction continue partout mais dérivable nulle part, brisant l'intuition que continuité implique dérivabilité.
\end{histoire}

% ── I. Continuité en un point ──
\section{Continuité en un point}

\begin{definition}[Continuité en un point]
Soit $f$ une fonction définie sur un intervalle $I$ et $a \in I$. On dit que $f$ est \textbf{continue en $a$} si :
\[ \lim_{x \to a} f(x) = f(a) \]
Autrement dit, la limite de $f$ en $a$ existe et vaut $f(a)$.
\end{definition}

\textit{Intuition :} $f$ est continue en $a$ si, lorsque $x$ se rapproche de $a$, $f(x)$ se rapproche de $f(a)$. Pas de « saut », pas de « trou » dans le graphe.

\begin{exemple}
\begin{itemize}
  \item Les fonctions polynomiales sont continues en tout point de $\R$.
  \item $f(x) = \ds\frac{1}{x}$ est continue en tout point de $]0, +\infty[$ (et de $]-\infty, 0[$), mais n'est pas définie en $0$.
  \item $f(x) = |x|$ est continue en tout point de $\R$, y compris en $0$.
\end{itemize}
\end{exemple}

\begin{propriete}[Fonctions continues usuelles]
Les fonctions suivantes sont continues sur leur domaine de définition :
\begin{itemize}
  \item les fonctions polynomiales (sur $\R$)
  \item les fonctions rationnelles (sur leur domaine)
  \item $\sqrt{x}$ (sur $[0, +\infty[$)
  \item $e^x$ (sur $\R$)
  \item $\ln x$ (sur $]0, +\infty[$)
  \item $\cos x$ et $\sin x$ (sur $\R$)
\end{itemize}
\end{propriete}

\subsection{Dérivable implique continue}

\begin{theoreme}[Dérivabilité et continuité]
Si $f$ est \textbf{dérivable} en $a$, alors $f$ est \textbf{continue} en $a$.
\end{theoreme}

\begin{demonstration}
Si $f$ est dérivable en $a$, alors $\ds\lim_{x \to a} \frac{f(x) - f(a)}{x - a} = f'(a)$ (limite finie). Or :
\[ f(x) - f(a) = \frac{f(x) - f(a)}{x-a} \cdot (x - a) \]
Quand $x \to a$, le premier facteur tend vers $f'(a)$ et le second vers $0$. Donc $f(x) - f(a) \to 0$, c'est-à-dire $f(x) \to f(a)$. Ainsi $f$ est continue en $a$. \qed
\end{demonstration}

\begin{attention}[La réciproque est fausse]
$f(x) = |x|$ est continue en $0$ mais pas dérivable en $0$ (la courbe présente un « angle »). La continuité est une condition nécessaire mais pas suffisante pour la dérivabilité.
\end{attention}

% ── II. Continuité sur un intervalle ──
\section{Continuité sur un intervalle}

\begin{definition}[Continuité sur un intervalle]
$f$ est \textbf{continue sur un intervalle $I$} si elle est continue en chaque point de $I$.
\end{definition}

\begin{propriete}[Opérations et continuité]
Somme, produit, quotient (si dénominateur non nul) et composée de fonctions continues sont des fonctions continues.
\end{propriete}

\begin{propriete}[Image d'un intervalle]
Si $f$ est continue sur un intervalle $[a, b]$, alors l'image $f([a, b])$ est aussi un intervalle.

En particulier, $f$ atteint toutes les valeurs comprises entre $f(a)$ et $f(b)$.
\end{propriete}

\textit{C'est l'idée clé : une fonction continue ne peut pas « sauter » de valeurs.}

% ── III. Théorème des valeurs intermédiaires ──
\section{Théorème des valeurs intermédiaires}

\begin{theoreme}[Théorème des valeurs intermédiaires (TVI)]
Soit $f$ une fonction \textbf{continue} sur un intervalle $[a, b]$.

Pour tout réel $k$ compris entre $f(a)$ et $f(b)$, il existe au moins un réel $c \in [a, b]$ tel que $f(c) = k$.
\end{theoreme}

\textit{En clair :} si $f$ est continue et passe de la valeur $f(a)$ à la valeur $f(b)$, elle prend toutes les valeurs intermédiaires. Elle ne peut pas « sauter » par-dessus une valeur.

\begin{center}
\begin{tikzpicture}
\begin{axis}[
  width=10cm, height=7cm,
  xlabel={$x$}, ylabel={$y$},
  xmin=-0.5, xmax=5, ymin=-1, ymax=5,
  xtick=\empty, ytick=\empty,
  axis lines=middle,
  clip=false
]
\addplot[indigo, very thick, domain=0.3:4.5, samples=100] {0.3*(x-1)*(x-3)*(x-4) + 2.5};
\addplot[rougeerr, dashed, thick, domain=-0.5:5] {2};
\draw[gray, dashed] (axis cs:0.5,0) -- (axis cs:0.5,{0.3*(0.5-1)*(0.5-3)*(0.5-4)+2.5});
\draw[gray, dashed] (axis cs:4.3,0) -- (axis cs:4.3,{0.3*(4.3-1)*(4.3-3)*(4.3-4)+2.5});
\node[below, font=\small] at (axis cs:0.5,0) {$a$};
\node[below, font=\small] at (axis cs:4.3,0) {$b$};
\node[left, rougeerr, font=\small] at (axis cs:0,2) {$k$};
\node[left, font=\small, gray] at (axis cs:0,{0.3*(0.5-1)*(0.5-3)*(0.5-4)+2.5}) {$f(a)$};
\node[left, font=\small, gray] at (axis cs:0,{0.3*(4.3-1)*(4.3-3)*(4.3-4)+2.5}) {$f(b)$};
% Mark intersection points
\fill[vertdef] (axis cs:0.72,2) circle (3pt);
\fill[vertdef] (axis cs:2.3,2) circle (3pt);
\fill[vertdef] (axis cs:3.98,2) circle (3pt);
\node[below, vertdef, font=\small] at (axis cs:0.72,0) {$c_1$};
\node[below, vertdef, font=\small] at (axis cs:2.3,0) {$c_2$};
\node[below, vertdef, font=\small] at (axis cs:3.98,0) {$c_3$};
\draw[vertdef, dashed] (axis cs:0.72,0) -- (axis cs:0.72,2);
\draw[vertdef, dashed] (axis cs:2.3,0) -- (axis cs:2.3,2);
\draw[vertdef, dashed] (axis cs:3.98,0) -- (axis cs:3.98,2);
\end{axis}
\end{tikzpicture}
\captionof{figure}{TVI : la droite $y = k$ coupe la courbe en au moins un point entre $a$ et $b$.}
\end{center}

\subsection{Corollaire : existence d'un zéro}

\begin{propriete}[Corollaire du TVI]
Si $f$ est continue sur $[a, b]$ et si $f(a)$ et $f(b)$ sont de \textbf{signes contraires} (c'est-à-dire $f(a) \cdot f(b) < 0$), alors il existe au moins un $c \in \mathopen{]}a, b\mathclose{[}$ tel que $f(c) = 0$.
\end{propriete}

\begin{exemple}
Montrons que l'équation $x^3 + x - 1 = 0$ a au moins une solution dans $[0, 1]$.

Posons $f(x) = x^3 + x - 1$. $f$ est continue sur $\R$ (polynôme).

$f(0) = -1 < 0$ et $f(1) = 1 > 0$. Comme $f(0) \cdot f(1) < 0$, par le corollaire du TVI, il existe $c \in \mathopen{]}0, 1\mathclose{[}$ tel que $f(c) = 0$.
\end{exemple}

\begin{attention}
Le TVI garantit l'\textbf{existence} d'un $c$ mais pas son \textbf{unicité}. Il peut y avoir plusieurs solutions (voir la figure ci-dessus avec $c_1$, $c_2$, $c_3$). Pour l'unicité, il faut une hypothèse supplémentaire (monotonie).
\end{attention}

% ── IV. TVI et monotonie : bijection ──
\section{TVI et monotonie : unicité de la solution}

\begin{theoreme}[Bijection]
Si $f$ est \textbf{continue et strictement monotone} sur un intervalle $I$, alors pour tout $k$ dans l'image de $f$, l'équation $f(x) = k$ a \textbf{une unique solution} dans $I$.
\end{theoreme}

\textit{Interprétation :} une fonction continue strictement monotone réalise une \textbf{bijection} de $I$ sur $f(I)$. Elle est « inversible ».

\begin{methode}[Prouver l'existence et l'unicité d'une solution]
Pour montrer que $f(x) = k$ a une unique solution sur $[a, b]$ :
\begin{enumerate}
  \item Vérifier que $f$ est \textbf{continue} sur $[a, b]$.
  \item Vérifier que $f$ est \textbf{strictement monotone} sur $[a, b]$ (étudier le signe de $f'$).
  \item Vérifier que $k$ est compris entre $f(a)$ et $f(b)$.
  \item \textbf{Conclure :} par le TVI appliqué à une fonction continue strictement monotone, $f(x) = k$ admet une unique solution $c \in [a, b]$.
\end{enumerate}
\end{methode}

\begin{exemple}
Montrons que l'équation $x^3 + x - 1 = 0$ a une \textbf{unique} solution sur $\R$.

Posons $f(x) = x^3 + x - 1$. $f'(x) = 3x^2 + 1 > 0$ pour tout $x \in \R$ : $f$ est strictement croissante sur $\R$.

De plus, $\ds\lim_{x \to -\infty} f(x) = -\infty$ et $\ds\lim_{x \to +\infty} f(x) = +\infty$, donc $f$ prend toutes les valeurs réelles.

Par le théorème de bijection, $f(x) = 0$ a une unique solution $\alpha \in \R$. On sait que $\alpha \in \mathopen{]}0, 1\mathclose{[}$ car $f(0) < 0 < f(1)$.
\end{exemple}

\begin{exemple}
Montrons que $e^x = 3x$ a une unique solution sur $[1, 2]$.

Posons $g(x) = e^x - 3x$. $g$ est continue et dérivable sur $\R$. $g'(x) = e^x - 3$.

$g'(x) = 0 \Leftrightarrow x = \ln 3 \approx 1{,}10$. Sur $[\ln 3, 2]$, $g' > 0$ : $g$ est strictement croissante.

$g(\ln 3) = 3 - 3\ln 3 \approx -0{,}30 < 0$ et $g(2) = e^2 - 6 \approx 1{,}39 > 0$.

Comme $g$ est strictement croissante sur $[\ln 3, 2]$ avec $g(\ln 3) < 0 < g(2)$, par le théorème de bijection, $g(x) = 0$ admet une unique solution dans $[\ln 3, 2] \subset [1, 2]$.
\end{exemple}

% ── V. Méthode de dichotomie ──
\section{Méthode de dichotomie}

\begin{definition}[Dichotomie]
La \textbf{dichotomie} est une méthode algorithmique pour approcher une solution de $f(x) = 0$ sur $[a, b]$ (quand $f(a) \cdot f(b) < 0$). On coupe l'intervalle en deux à chaque étape.
\end{definition}

\begin{methode}[Algorithme de dichotomie]
\begin{enumerate}
  \item Poser $a_0 = a$ et $b_0 = b$ avec $f(a_0) \cdot f(b_0) < 0$.
  \item Calculer le milieu $m = \ds\frac{a_0 + b_0}{2}$.
  \item Si $f(a_0) \cdot f(m) \leqslant 0$, la solution est dans $[a_0, m]$ : poser $b_1 = m$.
  \item Sinon, la solution est dans $[m, b_0]$ : poser $a_1 = m$.
  \item Recommencer avec $[a_1, b_1]$, et ainsi de suite.
\end{enumerate}

À chaque étape, la longueur de l'intervalle est divisée par $2$. Après $n$ étapes, la solution est connue à $\ds\frac{b-a}{2^n}$ près.
\end{methode}

\begin{exemple}[Dichotomie pour $x^3 + x - 1 = 0$ sur $[0, 1]$]
$f(x) = x^3 + x - 1$. On sait que $f(0) = -1 < 0$ et $f(1) = 1 > 0$.

\begin{center}
\renewcommand{\arraystretch}{1.4}
\begin{tabular}{ccccl}
\toprule
Étape & $a$ & $b$ & $m$ & Signe de $f(m)$ \\
\midrule
1 & $0$ & $1$ & $0{,}5$ & $f(0{,}5) = -0{,}375 < 0$ \\
2 & $0{,}5$ & $1$ & $0{,}75$ & $f(0{,}75) \approx 0{,}172 > 0$ \\
3 & $0{,}5$ & $0{,}75$ & $0{,}625$ & $f(0{,}625) \approx -0{,}131 < 0$ \\
4 & $0{,}625$ & $0{,}75$ & $0{,}6875$ & $f(0{,}6875) \approx 0{,}013 > 0$ \\
5 & $0{,}625$ & $0{,}6875$ & $0{,}65625$ & $f(0{,}65625) \approx -0{,}061 < 0$ \\
\bottomrule
\end{tabular}
\end{center}

Après 5 étapes : $c \in [0{,}65625 ;\, 0{,}6875]$, soit $c \approx 0{,}68$ à $0{,}02$ près.

Après 10 étapes : précision $\ds\frac{1}{2^{10}} \approx 10^{-3}$.
\end{exemple}

\begin{methode}[Implémentation en Python]
\begin{verbatim}
def dichotomie(f, a, b, epsilon):
    """Approche une solution de f(x)=0 sur [a,b]
       à epsilon près par dichotomie."""
    while b - a > epsilon:
        m = (a + b) / 2
        if f(a) * f(m) <= 0:
            b = m
        else:
            a = m
    return (a + b) / 2

# Exemple : x^3 + x - 1 = 0
f = lambda x: x**3 + x - 1
print(dichotomie(f, 0, 1, 1e-6))  # 0.6823278...
\end{verbatim}
\end{methode}

\begin{attention}[Nombre d'étapes nécessaires]
Pour obtenir une précision $\varepsilon$, il faut $n$ étapes avec $\ds\frac{b-a}{2^n} \leqslant \varepsilon$, soit $n \geqslant \ds\frac{\ln\!\left(\frac{b-a}{\varepsilon}\right)}{\ln 2}$.

Par exemple, pour $\varepsilon = 10^{-6}$ sur $[0, 1]$ : $n \geqslant \ds\frac{\ln(10^6)}{\ln 2} \approx 20$ étapes.
\end{attention}

% ── VI. Suites et point fixe ──
\section{Suites $u_{n+1} = f(u_n)$ et point fixe}

\subsection{Définition}

\begin{definition}[Suite récurrente $u_{n+1} = f(u_n)$]
Soit $f$ une fonction et $u_0$ un réel. La suite définie par :
\[ u_0 \text{ donné}, \qquad u_{n+1} = f(u_n) \]
est une \textbf{suite récurrente} d'ordre 1. Chaque terme est l'image du précédent par $f$.
\end{definition}

\begin{definition}[Point fixe]
Un \textbf{point fixe} de $f$ est un réel $\ell$ tel que $f(\ell) = \ell$.

Graphiquement, c'est un point d'intersection de la courbe $y = f(x)$ et de la droite $y = x$.
\end{definition}

\subsection{Convergence et point fixe}

\begin{propriete}[Limite et point fixe]
Si la suite $(u_n)$ définie par $u_{n+1} = f(u_n)$ converge vers $\ell$ et si $f$ est continue en $\ell$, alors $\ell$ est un point fixe de $f$ : $f(\ell) = \ell$.
\end{propriete}

\begin{demonstration}
Si $\ds\lim_{n \to +\infty} u_n = \ell$, alors $\ds\lim_{n \to +\infty} u_{n+1} = \ell$ aussi. En passant à la limite dans $u_{n+1} = f(u_n)$, et comme $f$ est continue en $\ell$, on obtient $\ell = f(\ell)$. \qed
\end{demonstration}

\begin{methode}[Étudier une suite $u_{n+1} = f(u_n)$]
\begin{enumerate}
  \item \textbf{Identifier $f$} et son domaine.
  \item \textbf{Chercher les points fixes} en résolvant $f(\ell) = \ell$.
  \item \textbf{Étudier la monotonie de $(u_n)$} (souvent par récurrence).
  \item \textbf{Montrer la convergence} : utiliser le théorème de convergence monotone (suite croissante majorée ou décroissante minorée).
  \item \textbf{Identifier la limite} : c'est un point fixe de $f$.
\end{enumerate}
\end{methode}

\begin{exemple}
Soit $u_0 = 4$ et $u_{n+1} = \ds\frac{1}{2}u_n + 1$.

\textbf{Fonction :} $f(x) = \ds\frac{x}{2} + 1$.

\textbf{Point fixe :} $f(\ell) = \ell \Leftrightarrow \ds\frac{\ell}{2} + 1 = \ell \Leftrightarrow \ell = 2$.

\textbf{Monotonie :} $u_{n+1} - u_n = f(u_n) - u_n = \ds\frac{u_n}{2} + 1 - u_n = 1 - \frac{u_n}{2}$.

Si $u_n > 2$, alors $u_{n+1} - u_n < 0$ : la suite est décroissante.

\textbf{Minoration :} montrons que $u_n > 2$ pour tout $n$. $u_0 = 4 > 2$. Si $u_n > 2$, alors $u_{n+1} = \ds\frac{u_n}{2} + 1 > \frac{2}{2} + 1 = 2$. OK par récurrence.

\textbf{Conclusion :} $(u_n)$ est décroissante et minorée par $2$, donc elle converge. Sa limite est le point fixe $\ell = 2$.

On peut vérifier : $u_0 = 4$, $u_1 = 3$, $u_2 = 2{,}5$, $u_3 = 2{,}25$, $u_4 = 2{,}125$\ldots{} la suite converge bien vers $2$.
\end{exemple}

\begin{exemple}[Suite $u_{n+1} = \cos(u_n)$]
Soit $u_0 = 0$ et $u_{n+1} = \cos(u_n)$.

Les premiers termes : $u_0 = 0$, $u_1 = 1$, $u_2 \approx 0{,}540$, $u_3 \approx 0{,}858$, $u_4 \approx 0{,}654$\ldots

La suite semble converger vers un point fixe $\ell$ vérifiant $\cos(\ell) = \ell$. Ce nombre, $\ell \approx 0{,}7391$, est appelé le \textbf{point fixe de Dottie} du cosinus. On peut montrer la convergence en utilisant le fait que $|\cos'(x)| = |\sin x| < 1$ sur un voisinage de $\ell$.
\end{exemple}

% ── Bilan ──
\section*{Bilan du chapitre}

\begin{retenir}
\begin{itemize}[leftmargin=*]
  \item \textbf{Continuité en $a$ :} $\ds\lim_{x \to a} f(x) = f(a)$ (pas de saut)
  \item \textbf{Dérivable $\Rightarrow$ continue} (mais la réciproque est fausse : $|x|$)
  \item \textbf{TVI :} si $f$ continue sur $[a,b]$ et $k$ entre $f(a)$ et $f(b)$, alors $\exists\, c \in [a,b]$, $f(c) = k$
  \item \textbf{Corollaire :} si $f(a) \cdot f(b) < 0$, alors $\exists\, c$, $f(c) = 0$
  \item \textbf{Bijection :} continue + strictement monotone $\Rightarrow$ solution unique
  \item \textbf{Dichotomie :} couper l'intervalle en $2$ à chaque étape ; précision $\ds\frac{b-a}{2^n}$ après $n$ étapes
  \item \textbf{Point fixe :} si $u_{n+1} = f(u_n)$ converge vers $\ell$ et $f$ continue, alors $f(\ell) = \ell$
\end{itemize}
\end{retenir}

% ─────────────────────────────────────────────────────────────────────────────
% CHAPITRE 9 — FONCTION LOGARITHME NÉPÉRIEN
% ─────────────────────────────────────────────────────────────────────────────

\chapter{Fonction logarithme népérien}

\section*{Introduction}

\textit{Qu'est-ce que le pH d'une solution chimique ? Comment mesure-t-on l'intensité d'un séisme sur l'échelle de Richter ? Pourquoi les décibels mesurent-ils le son sur une échelle « compressée » ?}

\medskip

Toutes ces grandeurs font intervenir le \textbf{logarithme}, une fonction qui transforme les multiplications en additions et permet de « comprimer » des échelles de grandeur très étendues. Le logarithme népérien, noté $\ln$, est défini comme la \textbf{fonction réciproque de l'exponentielle}. C'est l'une des fonctions les plus importantes de l'analyse mathématique.

\medskip

\textbf{Objectifs du chapitre :}
\begin{itemize}[leftmargin=*]
  \item Connaître la définition du logarithme népérien comme réciproque de l'exponentielle
  \item Maîtriser les propriétés algébriques du logarithme
  \item Calculer la dérivée de $\ln u$
  \item Connaître les limites et les croissances comparées
  \item Résoudre des équations et inéquations avec le logarithme
  \item Étudier des fonctions faisant intervenir $\ln$
\end{itemize}

\begin{histoire}[L'invention des logarithmes]
Le concept de logarithme est né au début du XVII\ieme{} siècle pour répondre à un besoin pratique : simplifier les calculs astronomiques. En 1614, le mathématicien écossais John \textbf{Napier} (ou Neper) publie la \textit{Mirifici Logarithmorum Canonis Descriptio}, introduisant les premières tables de logarithmes. L'idée géniale est de transformer les multiplications — pénibles à la main — en additions, beaucoup plus simples.

\medskip

En 1624, Henry \textbf{Briggs} adapte l'idée en base 10 (logarithme décimal). Le logarithme \textbf{népérien}, lié au nombre $e \approx 2{,}718$, émerge plus tard des travaux d'\textbf{Euler} (XVIII\ieme{} siècle) et s'impose comme le logarithme « naturel » de l'analyse mathématique, grâce à ses propriétés de dérivation remarquables.
\end{histoire}

% ── I. Définition ──
\section{Définition du logarithme népérien}

\subsection{La fonction exponentielle est une bijection}

Rappelons que la fonction exponentielle $\exp : \R \to \,]0\,;+\infty[$ est \textbf{continue} et \textbf{strictement croissante} sur $\R$. De plus, $\ds\lim_{x \to -\infty} e^x = 0$ et $\ds\lim_{x \to +\infty} e^x = +\infty$.

\medskip

D'après le théorème des valeurs intermédiaires, pour tout réel $b > 0$, l'équation $e^x = b$ admet \textbf{une unique solution}. On peut donc définir une fonction réciproque.

\begin{definition}[Logarithme népérien]
La \textbf{fonction logarithme népérien}, notée $\ln$, est définie sur $]0\,;+\infty[$ par :
\[ \forall\, x > 0,\quad y = \ln x \iff x = e^y \]
Autrement dit, $\ln x$ est l'unique réel $y$ tel que $e^y = x$.
\end{definition}

\begin{attention}[Ensemble de définition]
La fonction $\ln$ n'est définie que pour $x > 0$. On ne peut \textbf{jamais} écrire $\ln 0$, $\ln(-1)$ ou $\ln(\text{expression négative})$. Avant tout calcul, vérifier que l'argument est \textbf{strictement positif}.
\end{attention}

\subsection{Propriétés immédiates}

\begin{propriete}[Valeurs fondamentales et relations exp-ln]
Pour tout réel $x > 0$ et tout réel $y$ :
\begin{enumerate}
  \item $\ln 1 = 0$ \quad (car $e^0 = 1$)
  \item $\ln e = 1$ \quad (car $e^1 = e$)
  \item $\ln(e^y) = y$ pour tout $y \in \R$
  \item $e^{\ln x} = x$ pour tout $x > 0$
\end{enumerate}
\end{propriete}

\begin{demonstration}
Les propriétés 3 et 4 découlent directement de la définition : $\ln$ et $\exp$ sont réciproques l'une de l'autre.
\begin{itemize}
  \item $\ln(e^y) = y$ car $e^y$ est le nombre dont le logarithme vaut $y$.
  \item $e^{\ln x} = x$ car $\ln x$ est le réel dont l'exponentielle vaut $x$. \qed
\end{itemize}
\end{demonstration}

\begin{propriete}[Signe du logarithme]
\begin{itemize}
  \item Si $0 < x < 1$, alors $\ln x < 0$
  \item Si $x = 1$, alors $\ln x = 0$
  \item Si $x > 1$, alors $\ln x > 0$
\end{itemize}
\end{propriete}

\begin{exemple}
Quelques valeurs : $\ln 1 = 0$, \quad $\ln e = 1$, \quad $\ln e^2 = 2$, \quad $\ln e^{-3} = -3$, \quad $\ln \sqrt{e} = \frac{1}{2}$.
\end{exemple}

% ── II. Propriétés algébriques ──
\section{Propriétés algébriques}

\begin{theoreme}[Relation fonctionnelle du logarithme]
Pour tous réels $a > 0$ et $b > 0$ :
\[ \ln(ab) = \ln a + \ln b \]
Le logarithme \textbf{transforme un produit en somme}.
\end{theoreme}

\begin{demonstration}
Posons $\alpha = \ln a$ et $\beta = \ln b$. Alors $a = e^\alpha$ et $b = e^\beta$. On a :
\[ ab = e^\alpha \cdot e^\beta = e^{\alpha + \beta} \]
Donc $\ln(ab) = \alpha + \beta = \ln a + \ln b$. \qed
\end{demonstration}

\begin{propriete}[Conséquences algébriques]
Pour tous réels $a > 0$, $b > 0$ et tout entier relatif $n$ :
\begin{enumerate}
  \item $\ds\ln\left(\frac{a}{b}\right) = \ln a - \ln b$ \quad (quotient $\to$ différence)
  \item $\ds\ln\left(\frac{1}{b}\right) = -\ln b$ \quad (inverse $\to$ opposé)
  \item $\ds\ln(a^n) = n\ln a$ \quad (puissance $\to$ produit)
  \item $\ds\ln(\sqrt{a}) = \frac{1}{2}\ln a$ \quad (racine $\to$ demi-logarithme)
\end{enumerate}
\end{propriete}

\begin{demonstration}
\begin{enumerate}
  \item $\ds\ln\left(\frac{a}{b}\right) = \ln\left(a \cdot \frac{1}{b}\right) = \ln a + \ln\frac{1}{b}$. Or $\frac{1}{b} = e^{-\ln b}$, donc $\ln\frac{1}{b} = -\ln b$.
  \item Cas particulier du point 1 avec $a = 1$.
  \item Par récurrence : $\ln(a^n) = \ln(\underbrace{a \cdot a \cdots a}_{n \text{ fois}}) = n\ln a$.
  \item $\ln(\sqrt{a}) = \ln(a^{1/2}) = \frac{1}{2}\ln a$. \qed
\end{enumerate}
\end{demonstration}

\begin{attention}[Erreurs fréquentes avec le logarithme]
\begin{itemize}
  \item $\ln(a + b) \neq \ln a + \ln b$ \quad (le log d'une somme ne se simplifie \textbf{pas})
  \item $\ln(a - b) \neq \ln a - \ln b$ \quad (même erreur)
  \item $(\ln a)^2 \neq \ln(a^2)$ \quad ($(\ln a)^2 = (\ln a)(\ln a)$ tandis que $\ln(a^2) = 2\ln a$)
  \item $\ds\frac{\ln a}{\ln b} \neq \ln\frac{a}{b}$
\end{itemize}
\end{attention}

\begin{exemple}[Simplifications]
\begin{itemize}
  \item $\ln 8 = \ln 2^3 = 3\ln 2$
  \item $\ln 12 = \ln(4 \times 3) = \ln 4 + \ln 3 = 2\ln 2 + \ln 3$
  \item $\ds\ln\frac{5}{3} = \ln 5 - \ln 3$
  \item $\ln\sqrt{50} = \frac{1}{2}\ln 50 = \frac{1}{2}(\ln 25 + \ln 2) = \frac{1}{2}(2\ln 5 + \ln 2) = \ln 5 + \frac{1}{2}\ln 2$
\end{itemize}
\end{exemple}

\begin{methode}[Écrire un nombre en fonction de $\ln 2$ et $\ln 3$]
Pour exprimer $\ln N$ en fonction de $\ln 2$ et $\ln 3$, décomposer $N$ en produit de puissances de 2 et de 3.

\medskip

\textit{Exemple :} $\ln 72 = \ln(8 \times 9) = \ln 2^3 + \ln 3^2 = 3\ln 2 + 2\ln 3$.
\end{methode}

% ── III. Étude de la fonction logarithme népérien ──
\section{Étude de la fonction logarithme népérien}

\subsection{Dérivée}

\begin{theoreme}[Dérivée du logarithme népérien]
La fonction $\ln$ est dérivable sur $]0\,;+\infty[$ et :
\[ (\ln x)' = \frac{1}{x} \]
\end{theoreme}

\begin{demonstration}
Soit $f(x) = e^x$ et $g(x) = \ln x$. On a $f(g(x)) = x$ pour tout $x > 0$. En dérivant par la formule de la dérivée d'une composée :
\[ f'(g(x)) \cdot g'(x) = 1 \implies e^{\ln x} \cdot g'(x) = 1 \implies x \cdot g'(x) = 1 \implies g'(x) = \frac{1}{x} \qedhere \]
\end{demonstration}

\begin{propriete}[Dérivée de $\ln u$]
Si $u$ est une fonction dérivable et \textbf{strictement positive} sur un intervalle $I$, alors $\ln u$ est dérivable sur $I$ et :
\[ (\ln u)' = \frac{u'}{u} \]
\end{propriete}

\begin{exemple}[Dérivées]
\begin{enumerate}
  \item $f(x) = \ln(2x+1)$ pour $x > -\frac{1}{2}$ : $f'(x) = \ds\frac{2}{2x+1}$
  \item $g(x) = \ln(x^2+1)$ : $g'(x) = \ds\frac{2x}{x^2+1}$
  \item $h(x) = x\ln x$ pour $x > 0$ : $h'(x) = \ln x + x \cdot \frac{1}{x} = \ln x + 1$
  \item $k(x) = (\ln x)^2$ pour $x > 0$ : $k'(x) = 2\ln x \cdot \frac{1}{x} = \frac{2\ln x}{x}$
\end{enumerate}
\end{exemple}

\subsection{Sens de variation}

\begin{propriete}[Variations de $\ln$]
La fonction $\ln$ est \textbf{strictement croissante} sur $]0\,;+\infty[$.

\medskip

Conséquences : pour $a > 0$ et $b > 0$ :
\begin{itemize}
  \item $\ln a = \ln b \iff a = b$
  \item $\ln a < \ln b \iff a < b$
  \item $\ln a \leqslant \ln b \iff a \leqslant b$
\end{itemize}
\end{propriete}

\begin{demonstration}
Pour tout $x > 0$, $(\ln x)' = \frac{1}{x} > 0$. Donc $\ln$ est strictement croissante. \qed
\end{demonstration}

\subsection{Limites}

\begin{propriete}[Limites de $\ln$]
\[ \lim_{x \to 0^+} \ln x = -\infty \qquad \text{et} \qquad \lim_{x \to +\infty} \ln x = +\infty \]
\end{propriete}

\begin{demonstration}
\begin{itemize}
  \item Posons $x = e^t$. Quand $x \to +\infty$, $t \to +\infty$, donc $\ln x = t \to +\infty$.
  \item Posons $x = e^t$. Quand $x \to 0^+$, $t \to -\infty$, donc $\ln x = t \to -\infty$. \qed
\end{itemize}
\end{demonstration}

\begin{center}
\begin{tikzpicture}
\begin{axis}[
  width=12cm, height=8cm,
  axis lines=center,
  xlabel={$x$}, ylabel={$y$},
  xmin=-1, xmax=8, ymin=-4, ymax=4,
  xtick={1,2,...,7}, ytick={-3,-2,...,3},
  grid=both, grid style={gray!20},
  legend pos=south east
]
\addplot[domain=0.04:8, samples=200, thick, indigo] {ln(x)};
\addlegendentry{$y = \ln x$}
\addplot[domain=-3:2.1, samples=100, thick, rougeerr, dashed] {exp(x)};
\addlegendentry{$y = e^x$}
\addplot[domain=-1:4.5, samples=50, thin, gray, dashed] {x};
\addlegendentry{$y = x$}
\end{axis}
\end{tikzpicture}
\captionof{figure}{Courbes de $\ln$ et $\exp$, symétriques par rapport à la droite $y = x$}
\end{center}

\subsection{Courbes réciproques}

\begin{propriete}[Symétrie des courbes $\exp$ et $\ln$]
Les courbes de $\exp$ et $\ln$ sont \textbf{symétriques par rapport à la droite $y = x$}.

\medskip

En effet, le point $(a, b)$ est sur la courbe de $\exp$ si et seulement si $b = e^a$, c'est-à-dire $a = \ln b$, ce qui signifie que $(b, a)$ est sur la courbe de $\ln$.
\end{propriete}

% ── IV. Croissances comparées ──
\section{Croissances comparées}

\begin{theoreme}[Croissances comparées]
\begin{enumerate}
  \item $\ds\lim_{x \to +\infty} \frac{\ln x}{x} = 0$ \qquad (la puissance « l'emporte » sur le logarithme)
  \item $\ds\lim_{x \to 0^+} x\ln x = 0$ \qquad (la puissance « l'emporte » sur le logarithme)
\end{enumerate}
Plus généralement, pour tout entier $n \geqslant 1$ :
\[ \lim_{x \to +\infty} \frac{\ln x}{x^n} = 0 \qquad \text{et} \qquad \lim_{x \to 0^+} x^n \ln x = 0 \]
\end{theoreme}

\begin{demonstration}[de $\ds\lim_{x \to +\infty} \frac{\ln x}{x} = 0$]
Posons $x = e^t$. Quand $x \to +\infty$, $t \to +\infty$. On a :
\[ \frac{\ln x}{x} = \frac{t}{e^t} \]
Or on sait que $\ds\lim_{t \to +\infty} \frac{t}{e^t} = 0$ (croissance comparée exp/polynôme). Donc $\ds\lim_{x \to +\infty} \frac{\ln x}{x} = 0$. \qed
\end{demonstration}

\begin{methode}[Utiliser les croissances comparées]
Face à une forme indéterminée impliquant $\ln$ :
\begin{itemize}
  \item $\ds\frac{\ln x}{x^n}$ en $+\infty$ : le résultat est $0$ (la puissance gagne)
  \item $x^n \ln x$ en $0^+$ : le résultat est $0$ (la puissance gagne)
\end{itemize}
Penser éventuellement au changement de variable $x = e^t$ pour se ramener à des limites connues.
\end{methode}

\begin{exemple}[Limites avec croissances comparées]
\begin{enumerate}
  \item $\ds\lim_{x \to +\infty} \frac{\ln x}{\sqrt{x}} = \lim_{x \to +\infty} \frac{\ln x}{x^{1/2}} = 0$
  \item $\ds\lim_{x \to +\infty} \frac{(\ln x)^2}{x}$ : on pose $t = \ln x$, $x = e^t$, d'où $\frac{t^2}{e^t} \to 0$
  \item $\ds\lim_{x \to 0^+} x^2 \ln x = 0$
  \item $\ds\lim_{x \to 0^+} \sqrt{x}\,\ln x = 0$
\end{enumerate}
\end{exemple}

% ── V. Équations et inéquations ──
\section{Équations et inéquations logarithmiques}

\begin{methode}[Résoudre une équation avec $\ln$]
\begin{enumerate}
  \item \textbf{Déterminer le domaine de validité} : tous les arguments de $\ln$ doivent être $> 0$.
  \item \textbf{Utiliser les propriétés algébriques} pour simplifier (regrouper les $\ln$).
  \item \textbf{Exploiter l'injectivité} : $\ln A = \ln B \iff A = B$ (avec $A, B > 0$).
  \item \textbf{Ou passer à l'exponentielle} : $\ln A = c \iff A = e^c$.
  \item \textbf{Vérifier} que la solution appartient au domaine de validité.
\end{enumerate}
\end{methode}

\begin{exemple}[Équations]
\begin{enumerate}
  \item $\ln x = 3$ \quad $\Longleftrightarrow$ \quad $x = e^3$
  \item $\ln(2x-1) = \ln(x+4)$ \quad $\Longleftrightarrow$ \quad $2x - 1 = x + 4$ (si $2x-1 > 0$ et $x+4 > 0$) \quad $\Longleftrightarrow$ \quad $x = 5$. Vérification : $2(5)-1=9>0$ et $5+4=9>0$. \checkmark
  \item $\ln x + \ln(x-2) = \ln 3$ \quad $\Longleftrightarrow$ \quad $\ln[x(x-2)] = \ln 3$ (pour $x > 2$) \quad $\Longleftrightarrow$ \quad $x^2 - 2x = 3$ \quad $\Longleftrightarrow$ \quad $x^2 - 2x - 3 = 0$ \quad $\Longleftrightarrow$ \quad $(x-3)(x+1)=0$. Comme $x > 2$, la solution est $x = 3$.
\end{enumerate}
\end{exemple}

\begin{methode}[Résoudre une inéquation avec $\ln$]
On utilise la \textbf{croissance} de $\ln$ :
\begin{itemize}
  \item $\ln A \leqslant \ln B \iff A \leqslant B$ \quad (avec $A > 0$ et $B > 0$)
  \item $\ln A \leqslant c \iff A \leqslant e^c$ \quad (avec $A > 0$)
\end{itemize}
Ne pas oublier la condition de définition ($A > 0$).
\end{methode}

\begin{exemple}
$\ln(3x+1) \leqslant 2$ \quad $\Longleftrightarrow$ \quad $3x + 1 \leqslant e^2$ et $3x + 1 > 0$ \quad $\Longleftrightarrow$ \quad $x \leqslant \frac{e^2 - 1}{3}$ et $x > -\frac{1}{3}$.

Solution : $x \in \left]-\frac{1}{3}\,;\,\frac{e^2-1}{3}\right]$.
\end{exemple}

% ── VI. Études de fonctions ──
\section{Études de fonctions avec $\ln$}

\begin{methode}[Étude complète d'une fonction]
Pour étudier une fonction $f$ faisant intervenir $\ln$ :
\begin{enumerate}
  \item Déterminer le \textbf{domaine de définition} ($\ln$ exige un argument $> 0$)
  \item Calculer la \textbf{dérivée} et étudier son signe
  \item Dresser le \textbf{tableau de variation}
  \item Calculer les \textbf{limites} aux bornes (croissances comparées si besoin)
  \item Rechercher d'éventuelles \textbf{asymptotes}
  \item Tracer la \textbf{courbe}
\end{enumerate}
\end{methode}

\begin{exemple}[Étude de $f(x) = x - \ln x$]
\textbf{Domaine :} $D_f = \,]0\,;+\infty[$

\textbf{Dérivée :} $f'(x) = 1 - \frac{1}{x} = \frac{x-1}{x}$

\textbf{Signe :} $f'(x) < 0$ si $0 < x < 1$, \quad $f'(x) = 0$ si $x = 1$, \quad $f'(x) > 0$ si $x > 1$.

\textbf{Minimum :} $f(1) = 1 - \ln 1 = 1$.

\textbf{Limites :}
\begin{itemize}
  \item $\ds\lim_{x \to 0^+} f(x) = 0 - (-\infty) = +\infty$
  \item $\ds\lim_{x \to +\infty} f(x) = +\infty$ \quad car $f(x) = x\left(1 - \frac{\ln x}{x}\right)$ et $\frac{\ln x}{x} \to 0$
\end{itemize}
\end{exemple}

\begin{exemple}[Étude de $g(x) = x\ln x$]
\textbf{Domaine :} $D_g = \,]0\,;+\infty[$

\textbf{Dérivée :} $g'(x) = \ln x + 1$

\textbf{Signe :} $g'(x) = 0 \iff \ln x = -1 \iff x = e^{-1} = \frac{1}{e}$. On a $g'(x) < 0$ sur $]0\,;\frac{1}{e}[$ et $g'(x) > 0$ sur $]\frac{1}{e}\,;+\infty[$.

\textbf{Minimum :} $g\left(\frac{1}{e}\right) = \frac{1}{e} \cdot (-1) = -\frac{1}{e}$.

\textbf{Limites :} $\ds\lim_{x \to 0^+} x\ln x = 0$ (croissance comparée) et $\ds\lim_{x \to +\infty} x\ln x = +\infty$.
\end{exemple}

\begin{exemple}[Étude de $h(x) = \ln(x^2 + 1)$]
\textbf{Domaine :} $D_h = \R$ \quad (car $x^2 + 1 > 0$ pour tout $x$)

\textbf{Dérivée :} $h'(x) = \ds\frac{2x}{x^2+1}$

\textbf{Signe :} $h'(x) < 0$ si $x < 0$, \quad $h'(x) = 0$ si $x = 0$, \quad $h'(x) > 0$ si $x > 0$.

\textbf{Minimum :} $h(0) = \ln 1 = 0$.

\textbf{Parité :} $h(-x) = \ln((-x)^2 + 1) = \ln(x^2+1) = h(x)$. La fonction est \textbf{paire}, sa courbe est symétrique par rapport à l'axe des ordonnées.
\end{exemple}

% ── Bilan ──
\section*{Bilan du chapitre}

\begin{retenir}
\begin{itemize}[leftmargin=*]
  \item \textbf{Définition :} $y = \ln x \iff x = e^y$ \quad ($x > 0$)
  \item \textbf{Valeurs :} $\ln 1 = 0$, \quad $\ln e = 1$, \quad $\ln(e^y) = y$, \quad $e^{\ln x} = x$
  \item \textbf{Propriétés algébriques :}
    $\ln(ab) = \ln a + \ln b$, \quad $\ln\frac{a}{b} = \ln a - \ln b$, \quad $\ln(a^n) = n\ln a$
  \item \textbf{Dérivée :} $(\ln x)' = \frac{1}{x}$, \quad $(\ln u)' = \frac{u'}{u}$
  \item \textbf{Variations :} $\ln$ est strictement croissante sur $]0\,;+\infty[$
  \item \textbf{Limites :} $\ds\lim_{x \to 0^+} \ln x = -\infty$, \quad $\ds\lim_{x \to +\infty} \ln x = +\infty$
  \item \textbf{Croissances comparées :} $\ds\lim_{x \to +\infty} \frac{\ln x}{x} = 0$, \quad $\ds\lim_{x \to 0^+} x\ln x = 0$
  \item \textbf{Équations :} $\ln A = \ln B \iff A = B$ \quad ($A,B > 0$)
\end{itemize}
\end{retenir}


% ─────────────────────────────────────────────────────────────────────────────
% CHAPITRE 10 — FONCTIONS TRIGONOMÉTRIQUES
% ─────────────────────────────────────────────────────────────────────────────

\chapter{Fonctions trigonométriques}

\section*{Introduction}

\textit{Pourquoi le son d'une guitare, la lumière d'un laser et la marée haute à Saint-Malo obéissent-ils aux mêmes lois mathématiques ?}

\medskip

Tous ces phénomènes sont \textbf{périodiques} : ils se reproduisent identiquement à intervalles réguliers. Les fonctions sinus et cosinus, issues du cercle trigonométrique, sont les outils mathématiques naturels pour modéliser ces phénomènes. Ce chapitre étudie ces fonctions du point de vue de l'\textbf{analyse} : dérivabilité, variations, équations et inéquations.

\medskip

\textbf{Objectifs du chapitre :}
\begin{itemize}[leftmargin=*]
  \item Maîtriser le cercle trigonométrique et les valeurs remarquables
  \item Connaître les dérivées de $\sin$ et $\cos$
  \item Étudier les variations des fonctions trigonométriques
  \item Résoudre des équations et inéquations trigonométriques
  \item Étudier des fonctions faisant intervenir $\sin$ et $\cos$
\end{itemize}

\begin{histoire}[De l'astronomie à l'analyse]
La trigonométrie est née de l'observation du ciel. Les astronomes babyloniens et grecs (Hipparque, II\ieme{} siècle av. J.-C.) calculaient déjà des tables de cordes pour prédire les positions des astres. Les mathématiciens indiens (Aryabhata, V\ieme{} siècle) introduisent les fonctions sinus (\textit{jya}) et cosinus. Les Arabes (Al-Battani, IX\ieme{} siècle) perfectionnent ces tables et développent la tangente.

\medskip

Au XVIII\ieme{} siècle, \textbf{Euler} systématise les fonctions trigonométriques comme fonctions de la variable réelle et établit la formule $e^{ix} = \cos x + i\sin x$. En 1807, \textbf{Joseph Fourier} montre que toute fonction périodique peut s'écrire comme somme de sinus et cosinus (séries de Fourier), révolutionnant la physique et l'ingénierie — du traitement du son à la compression d'images (JPEG, MP3).
\end{histoire}

% ── I. Rappels sur le cercle trigonométrique ──
\section{Rappels sur le cercle trigonométrique}

\subsection{Définitions}

\begin{definition}[Radian et cercle trigonométrique]
Le \textbf{cercle trigonométrique} est le cercle de centre $O$ et de rayon 1, muni du sens direct (sens inverse des aiguilles d'une montre). À tout réel $x$, on associe un unique point $M$ sur ce cercle, obtenu en « enroulant » la droite des réels autour du cercle à partir du point $(1, 0)$.

\medskip

Les coordonnées de $M$ sont $(\cos x,\, \sin x)$.
\end{definition}

\begin{center}
\begin{tikzpicture}[scale=2.5]
  % Cercle
  \draw[gray!40] (-1.3,0) -- (1.3,0);
  \draw[gray!40] (0,-1.3) -- (0,1.3);
  \draw[thick] (0,0) circle (1);
  % Point M
  \coordinate (M) at ({cos(50)},{sin(50)});
  \draw[thick,indigo] (0,0) -- (M);
  \fill[indigo] (M) circle (1pt) node[above right] {$M(\cos x,\sin x)$};
  % Projections
  \draw[dashed,vertdef] (M) -- ({cos(50)},0) node[below] {$\cos x$};
  \draw[dashed,rougeerr] (M) -- (0,{sin(50)}) node[left] {$\sin x$};
  % Arc
  \draw[thick,orangemeth,-{Stealth}] (0.35,0) arc (0:50:0.35) node[midway,right] {$x$};
  % Axes
  \node[below right] at (1.3,0) {$x$};
  \node[above left] at (0,1.3) {$y$};
  % Points remarquables
  \fill (1,0) circle (0.8pt) node[below right] {$1$};
  \fill (0,1) circle (0.8pt) node[above left] {$1$};
  \fill (-1,0) circle (0.8pt) node[below left] {$-1$};
  \fill (0,-1) circle (0.8pt) node[below left] {$-1$};
\end{tikzpicture}
\captionof{figure}{Cercle trigonométrique et coordonnées du point $M$}
\end{center}

\subsection{Valeurs remarquables}

\begin{propriete}[Tableau des valeurs remarquables]
\begin{center}
\renewcommand{\arraystretch}{1.8}
\begin{tabular}{|c|c|c|c|c|c|}
\hline
$x$ & $0$ & $\ds\frac{\pi}{6}$ & $\ds\frac{\pi}{4}$ & $\ds\frac{\pi}{3}$ & $\ds\frac{\pi}{2}$ \\
\hline
$\cos x$ & $1$ & $\ds\frac{\sqrt{3}}{2}$ & $\ds\frac{\sqrt{2}}{2}$ & $\ds\frac{1}{2}$ & $0$ \\
\hline
$\sin x$ & $0$ & $\ds\frac{1}{2}$ & $\ds\frac{\sqrt{2}}{2}$ & $\ds\frac{\sqrt{3}}{2}$ & $1$ \\
\hline
\end{tabular}
\end{center}
\end{propriete}

\subsection{Propriétés fondamentales}

\begin{propriete}[Propriétés de $\cos$ et $\sin$]
Pour tout réel $x$ :
\begin{enumerate}
  \item $\cos^2 x + \sin^2 x = 1$ \quad (relation fondamentale)
  \item $-1 \leqslant \cos x \leqslant 1$ et $-1 \leqslant \sin x \leqslant 1$
  \item $\cos$ et $\sin$ sont \textbf{périodiques} de période $2\pi$ : $\cos(x + 2\pi) = \cos x$, $\sin(x + 2\pi) = \sin x$
  \item $\cos$ est \textbf{paire} : $\cos(-x) = \cos x$
  \item $\sin$ est \textbf{impaire} : $\sin(-x) = -\sin x$
  \item $\cos(\pi - x) = -\cos x$ et $\sin(\pi - x) = \sin x$
  \item $\cos\left(\frac{\pi}{2} - x\right) = \sin x$ et $\sin\left(\frac{\pi}{2} - x\right) = \cos x$
\end{enumerate}
\end{propriete}

\begin{attention}[Degrés et radians]
En mathématiques, les angles sont \textbf{toujours en radians}. Ne pas confondre $\sin(30)$ (30 radians, valeur sans intérêt) et $\sin\frac{\pi}{6} = 0{,}5$ (30 degrés). La formule de conversion est : $\alpha_{\text{rad}} = \frac{\pi}{180}\,\alpha_{\text{deg}}$.
\end{attention}

% ── II. Dérivées des fonctions trigonométriques ──
\section{Dérivées des fonctions trigonométriques}

\begin{theoreme}[Dérivées de $\sin$ et $\cos$]
Les fonctions $\sin$ et $\cos$ sont dérivables sur $\R$ et :
\[ (\sin x)' = \cos x \qquad \text{et} \qquad (\cos x)' = -\sin x \]
\end{theoreme}

\begin{demonstration}[de $(\sin x)' = \cos x$]
On utilise la définition du nombre dérivé :
\[ \frac{\sin(x+h) - \sin x}{h} = \frac{\sin x \cos h + \cos x \sin h - \sin x}{h} = \sin x \cdot \frac{\cos h - 1}{h} + \cos x \cdot \frac{\sin h}{h} \]
Or, quand $h \to 0$ : $\ds\frac{\sin h}{h} \to 1$ et $\ds\frac{\cos h - 1}{h} \to 0$.

Donc $(\sin x)' = \sin x \cdot 0 + \cos x \cdot 1 = \cos x$. \qed
\end{demonstration}

\begin{propriete}[Dérivées des composées]
Si $u$ est une fonction dérivable :
\[ (\sin u)' = u' \cos u \qquad \text{et} \qquad (\cos u)' = -u' \sin u \]
\end{propriete}

\begin{exemple}[Dérivées]
\begin{enumerate}
  \item $f(x) = \sin(3x)$ : $f'(x) = 3\cos(3x)$
  \item $g(x) = \cos(2x + 1)$ : $g'(x) = -2\sin(2x+1)$
  \item $h(x) = \sin^2 x = (\sin x)^2$ : $h'(x) = 2\sin x \cdot \cos x = \sin(2x)$
  \item $k(x) = x + \cos x$ : $k'(x) = 1 - \sin x$
\end{enumerate}
\end{exemple}

\begin{attention}[Signe moins pour le cosinus]
La dérivée de $\cos x$ est $-\sin x$ (avec un signe $-$). C'est une source d'erreur fréquente. Moyen mnémotechnique : « cosinus perd un signe en dérivant ».
\end{attention}

% ── III. Variations des fonctions sin et cos ──
\section{Variations des fonctions sinus et cosinus}

\subsection{Étude de la fonction sinus sur $[0\,;2\pi]$}

\begin{propriete}[Variations de $\sin$ sur $[0\,;2\pi]$]
La dérivée de $\sin$ est $\cos$. Son signe sur $[0\,;2\pi]$ donne :
\begin{itemize}
  \item $\cos x > 0$ sur $\left[0\,;\frac{\pi}{2}\right[$ : $\sin$ est strictement croissante
  \item $\cos x < 0$ sur $\left]\frac{\pi}{2}\,;\frac{3\pi}{2}\right[$ : $\sin$ est strictement décroissante
  \item $\cos x > 0$ sur $\left]\frac{3\pi}{2}\,;2\pi\right]$ : $\sin$ est strictement croissante
\end{itemize}
Maximum : $\sin\frac{\pi}{2} = 1$. \quad Minimum : $\sin\frac{3\pi}{2} = -1$.
\end{propriete}

\subsection{Étude de la fonction cosinus sur $[0\,;2\pi]$}

\begin{propriete}[Variations de $\cos$ sur $[0\,;2\pi]$]
La dérivée de $\cos$ est $-\sin$. Son signe sur $[0\,;2\pi]$ donne :
\begin{itemize}
  \item $-\sin x < 0$ sur $]0\,;\pi[$ : $\cos$ est strictement décroissante
  \item $-\sin x > 0$ sur $]\pi\,;2\pi[$ : $\cos$ est strictement croissante
\end{itemize}
Maximum : $\cos 0 = \cos 2\pi = 1$. \quad Minimum : $\cos\pi = -1$.
\end{propriete}

\begin{center}
\begin{tikzpicture}
\begin{axis}[
  width=14cm, height=6cm,
  axis lines=center,
  xlabel={$x$}, ylabel={$y$},
  xmin=-0.5, xmax=7, ymin=-1.5, ymax=1.5,
  xtick={0, 1.5708, 3.1416, 4.7124, 6.2832},
  xticklabels={$0$, $\frac{\pi}{2}$, $\pi$, $\frac{3\pi}{2}$, $2\pi$},
  ytick={-1,0,1},
  grid=both, grid style={gray!15},
  legend pos=north east
]
\addplot[domain=0:6.5, samples=200, thick, indigo] {sin(deg(x))};
\addlegendentry{$y = \sin x$}
\addplot[domain=0:6.5, samples=200, thick, rougeerr, dashed] {cos(deg(x))};
\addlegendentry{$y = \cos x$}
\end{axis}
\end{tikzpicture}
\captionof{figure}{Courbes des fonctions sinus et cosinus sur $[0\,;2\pi]$}
\end{center}

% ── IV. Équations trigonométriques ──
\section{Résolution d'équations trigonométriques}

\begin{methode}[Équation $\cos x = a$]
Pour $a \in [-1\,;1]$, l'équation $\cos x = a$ a exactement deux familles de solutions :
\[ x = \arccos a + 2k\pi \quad \text{ou} \quad x = -\arccos a + 2k\pi \qquad (k \in \Z) \]
Que l'on écrit aussi : $x = \pm\, \arccos a + 2k\pi$, \quad $k \in \Z$.

\medskip

\textit{Cas particuliers :} $\cos x = 1 \iff x = 2k\pi$ ; \quad $\cos x = 0 \iff x = \frac{\pi}{2} + k\pi$ ; \quad $\cos x = -1 \iff x = \pi + 2k\pi$.
\end{methode}

\begin{methode}[Équation $\sin x = b$]
Pour $b \in [-1\,;1]$, l'équation $\sin x = b$ a exactement deux familles de solutions :
\[ x = \arcsin b + 2k\pi \quad \text{ou} \quad x = \pi - \arcsin b + 2k\pi \qquad (k \in \Z) \]

\textit{Cas particuliers :} $\sin x = 0 \iff x = k\pi$ ; \quad $\sin x = 1 \iff x = \frac{\pi}{2} + 2k\pi$ ; \quad $\sin x = -1 \iff x = -\frac{\pi}{2} + 2k\pi$.
\end{methode}

\begin{exemple}[Résolution d'équations]
\begin{enumerate}
  \item $\cos x = \frac{1}{2}$ : on sait que $\cos\frac{\pi}{3} = \frac{1}{2}$, donc $x = \frac{\pi}{3} + 2k\pi$ ou $x = -\frac{\pi}{3} + 2k\pi$, $k \in \Z$.
  \item $\sin x = \frac{\sqrt{2}}{2}$ : on sait que $\sin\frac{\pi}{4} = \frac{\sqrt{2}}{2}$, donc $x = \frac{\pi}{4} + 2k\pi$ ou $x = \frac{3\pi}{4} + 2k\pi$, $k \in \Z$.
  \item $2\cos x - 1 = 0$ sur $[0\,;2\pi]$ : $\cos x = \frac{1}{2}$, d'où $x = \frac{\pi}{3}$ ou $x = \frac{5\pi}{3}$.
\end{enumerate}
\end{exemple}

\begin{attention}[Ne pas oublier les familles de solutions]
Contrairement aux équations avec $\ln$ ou $\exp$, une équation trigonométrique a \textbf{une infinité de solutions} (à cause de la périodicité). Il faut toujours écrire le « $+ 2k\pi$ ». Sur un intervalle borné, on sélectionne ensuite les valeurs de $k$ qui conviennent.
\end{attention}

% ── V. Inéquations trigonométriques ──
\section{Inéquations trigonométriques}

\begin{methode}[Résoudre une inéquation sur un intervalle]
Pour résoudre $\cos x \geqslant a$ (ou $\sin x \leqslant b$) sur un intervalle $[0\,;2\pi]$ :
\begin{enumerate}
  \item Résoudre l'équation $\cos x = a$ (ou $\sin x = b$) sur cet intervalle.
  \item Utiliser le tableau de variation ou la courbe pour déterminer les intervalles où l'inégalité est satisfaite.
  \item S'aider du cercle trigonométrique pour visualiser.
\end{enumerate}
\end{methode}

\begin{exemple}
Résoudre $\sin x \geqslant \frac{1}{2}$ sur $[0\,;2\pi]$.

On résout d'abord $\sin x = \frac{1}{2}$ : $x = \frac{\pi}{6}$ ou $x = \frac{5\pi}{6}$.

D'après la courbe de $\sin$, $\sin x \geqslant \frac{1}{2}$ entre ces deux valeurs. Solution : $x \in \left[\frac{\pi}{6}\,;\,\frac{5\pi}{6}\right]$.
\end{exemple}

% ── VI. Études de fonctions trigonométriques ──
\section{Études de fonctions trigonométriques}

\begin{methode}[Exploiter la périodicité et la parité]
Avant de dériver, vérifier si la fonction est :
\begin{itemize}
  \item \textbf{Périodique} de période $T$ : étudier sur un intervalle de longueur $T$ suffit.
  \item \textbf{Paire} : étudier sur $[0\,;+\infty[$ puis compléter par symétrie.
  \item \textbf{Impaire} : même principe, avec symétrie par rapport à $O$.
\end{itemize}
\end{methode}

\begin{exemple}[Étude de $f(x) = \cos x + \sin x$]
\textbf{Périodicité :} $f(x + 2\pi) = \cos(x+2\pi) + \sin(x+2\pi) = \cos x + \sin x = f(x)$. Période $2\pi$.

On étudie $f$ sur $[0\,;2\pi]$.

\textbf{Dérivée :} $f'(x) = -\sin x + \cos x$

\textbf{Annulation :} $f'(x) = 0 \iff \cos x = \sin x \iff \tan x = 1$ (quand $\cos x \neq 0$), d'où $x = \frac{\pi}{4}$ ou $x = \frac{5\pi}{4}$ sur $[0\,;2\pi]$.

\textbf{Maximum :} $f\left(\frac{\pi}{4}\right) = \frac{\sqrt{2}}{2} + \frac{\sqrt{2}}{2} = \sqrt{2}$

\textbf{Minimum :} $f\left(\frac{5\pi}{4}\right) = -\frac{\sqrt{2}}{2} - \frac{\sqrt{2}}{2} = -\sqrt{2}$

\medskip

\textit{Remarque :} on peut aussi écrire $f(x) = \sqrt{2}\cos\left(x - \frac{\pi}{4}\right)$, ce qui montre directement l'amplitude $\sqrt{2}$ et le déphasage $\frac{\pi}{4}$.
\end{exemple}

\begin{exemple}[Étude de $g(x) = \cos(2x)$]
\textbf{Périodicité :} $g(x + \pi) = \cos(2(x+\pi)) = \cos(2x + 2\pi) = \cos(2x) = g(x)$. Période $\pi$.

\textbf{Parité :} $g(-x) = \cos(-2x) = \cos(2x) = g(x)$. Fonction paire.

On étudie $g$ sur $[0\,;\pi]$.

\textbf{Dérivée :} $g'(x) = -2\sin(2x)$

$g'(x) = 0 \iff \sin(2x) = 0 \iff 2x = k\pi \iff x = \frac{k\pi}{2}$. Sur $[0\,;\pi]$ : $x = 0$, $x = \frac{\pi}{2}$, $x = \pi$.

$g$ est décroissante sur $\left[0\,;\frac{\pi}{2}\right]$ et croissante sur $\left[\frac{\pi}{2}\,;\pi\right]$.
\end{exemple}

% ── Bilan ──
\section*{Bilan du chapitre}

\begin{retenir}
\begin{itemize}[leftmargin=*]
  \item \textbf{Relation fondamentale :} $\cos^2 x + \sin^2 x = 1$
  \item \textbf{Dérivées :} $(\sin x)' = \cos x$, \quad $(\cos x)' = -\sin x$
  \item \textbf{Composées :} $(\sin u)' = u'\cos u$, \quad $(\cos u)' = -u'\sin u$
  \item \textbf{Périodicité :} $\sin$ et $\cos$ sont de période $2\pi$
  \item \textbf{Parité :} $\cos$ est paire, $\sin$ est impaire
  \item \textbf{Équation $\cos x = a$ :} $x = \pm\,\arccos a + 2k\pi$
  \item \textbf{Équation $\sin x = b$ :} $x = \arcsin b + 2k\pi$ ou $x = \pi - \arcsin b + 2k\pi$
  \item \textbf{Variations :} $\sin$ croissante sur $\left[-\frac{\pi}{2}\,;\frac{\pi}{2}\right]$, $\cos$ décroissante sur $[0\,;\pi]$
  \item \textbf{Stratégie :} exploiter périodicité et parité pour réduire l'intervalle d'étude
\end{itemize}
\end{retenir}


% ─────────────────────────────────────────────────────────────────────────────
% CHAPITRE 11 — PRIMITIVES ET ÉQUATIONS DIFFÉRENTIELLES
% ─────────────────────────────────────────────────────────────────────────────

\chapter{Primitives et équations différentielles}

\section*{Introduction}

\textit{Si l'on connaît la vitesse d'un objet à chaque instant, comment retrouver sa position ? Si la dérivée d'une fonction est connue, peut-on « remonter » à la fonction elle-même ?}

\medskip

Ce problème inverse de la dérivation conduit à la notion de \textbf{primitive}. Dériver, c'est calculer un taux de variation instantané ; « primitiver », c'est reconstruire une fonction à partir de son taux de variation. Les primitives servent aussi à résoudre des \textbf{équations différentielles}, c'est-à-dire des équations dont l'inconnue est une fonction et qui relient la fonction à sa dérivée. Ces équations modélisent des phénomènes fondamentaux : décroissance radioactive, refroidissement, circuits électriques, dynamique des populations.

\medskip

\textbf{Objectifs du chapitre :}
\begin{itemize}[leftmargin=*]
  \item Connaître la définition de primitive et savoir en calculer
  \item Utiliser le tableau des primitives usuelles, y compris les composées
  \item Déterminer une primitive vérifiant une condition initiale
  \item Résoudre l'équation différentielle $y' = ay$
  \item Résoudre l'équation différentielle $y' = ay + b$
  \item Modéliser des situations concrètes par des équations différentielles
\end{itemize}

\begin{histoire}[La naissance du calcul infinitésimal]
Le concept de primitive est indissociable de l'invention du calcul infinitésimal au XVII\ieme{} siècle. \textbf{Isaac Newton} (1643--1727) développe sa « méthode des fluxions » pour résoudre des problèmes de mécanique : connaissant la vitesse, retrouver la position. Indépendamment, \textbf{Gottfried Wilhelm Leibniz} (1646--1716) crée le calcul différentiel et intégral, introduisant les notations $\frac{dy}{dx}$ et $\int$ encore utilisées aujourd'hui. Leur dispute de priorité est l'une des plus célèbres de l'histoire des sciences. Le théorème fondamental de l'analyse, qui relie primitive et intégrale, unifie leurs deux approches.
\end{histoire}

% ── I. Notion de primitive ──
\section{Notion de primitive}

\subsection{Définition}

\begin{definition}[Primitive]
Soit $f$ une fonction définie sur un intervalle $I$. On dit que $F$ est une \textbf{primitive} de $f$ sur $I$ si $F$ est dérivable sur $I$ et :
\[ \forall\, x \in I,\quad F'(x) = f(x) \]
\end{definition}

\textit{Autrement dit, une primitive de $f$ est une fonction dont la dérivée est $f$. Primitiver, c'est le problème inverse de dériver.}

\begin{exemple}
\begin{itemize}
  \item $F(x) = x^3$ est une primitive de $f(x) = 3x^2$ car $F'(x) = 3x^2$.
  \item $G(x) = \sin x$ est une primitive de $g(x) = \cos x$ car $G'(x) = \cos x$.
  \item $H(x) = e^x$ est une primitive de $h(x) = e^x$ car $H'(x) = e^x$.
\end{itemize}
\end{exemple}

\subsection{Unicité à une constante près}

\begin{theoreme}[Primitives et constante]
Si $F$ est une primitive de $f$ sur un intervalle $I$, alors :
\begin{enumerate}
  \item Pour toute constante $C \in \R$, la fonction $G = F + C$ est aussi une primitive de $f$.
  \item \textbf{Toute} primitive de $f$ sur $I$ est de la forme $F + C$ avec $C \in \R$.
\end{enumerate}
Autrement dit, \textbf{deux primitives d'une même fonction diffèrent d'une constante}.
\end{theoreme}

\begin{demonstration}
\begin{enumerate}
  \item $(F + C)' = F' + 0 = f$. Donc $F + C$ est bien une primitive de $f$.
  \item Si $G$ est aussi une primitive de $f$, alors $(G - F)' = G' - F' = f - f = 0$. Or, une fonction de dérivée nulle sur un intervalle est constante. Donc $G - F = C$ pour un certain $C \in \R$, d'où $G = F + C$. \qed
\end{enumerate}
\end{demonstration}

\begin{attention}[Intervalle obligatoire]
Le résultat « deux primitives diffèrent d'une constante » n'est vrai que sur un \textbf{intervalle}. Sur une réunion de deux intervalles disjoints, la constante peut être différente sur chacun.
\end{attention}

\subsection{Existence}

\begin{theoreme}[Existence des primitives]
Toute fonction \textbf{continue} sur un intervalle $I$ admet des primitives sur $I$.
\end{theoreme}

\textit{Ce théorème, admis, garantit que le problème « trouver $F$ telle que $F' = f$ » a toujours une solution dès que $f$ est continue.}

% ── II. Tableau des primitives ──
\section{Tableau des primitives usuelles}

\subsection{Primitives des fonctions de référence}

\begin{propriete}[Primitives usuelles]
\begin{center}
\renewcommand{\arraystretch}{2}
\begin{tabular}{|c|c|c|}
\hline
\textbf{Fonction $f(x)$} & \textbf{Primitive $F(x)$} & \textbf{Intervalle} \\
\hline
$k$ (constante) & $kx$ & $\R$ \\
\hline
$x^n$ ($n \in \N$) & $\ds\frac{x^{n+1}}{n+1}$ & $\R$ \\
\hline
$\ds\frac{1}{x^n}$ ($n \geqslant 2$) & $\ds\frac{-1}{(n-1)x^{n-1}}$ & $]0\,;+\infty[$ \\
\hline
$\ds\frac{1}{\sqrt{x}}$ & $2\sqrt{x}$ & $]0\,;+\infty[$ \\
\hline
$\ds\frac{1}{x}$ & $\ln|x|$ & $]0\,;+\infty[$ ou $]-\infty\,;0[$ \\
\hline
$e^x$ & $e^x$ & $\R$ \\
\hline
$\cos x$ & $\sin x$ & $\R$ \\
\hline
$\sin x$ & $-\cos x$ & $\R$ \\
\hline
\end{tabular}
\end{center}
\end{propriete}

\begin{attention}[Signe moins pour $\sin$]
La primitive de $\sin x$ est $-\cos x$ (avec un signe $-$). Vérification : $(-\cos x)' = \sin x$. \checkmark
\end{attention}

\subsection{Primitives des composées}

\begin{propriete}[Primitives de fonctions composées]
Soit $u$ une fonction dérivable. Les primitives des fonctions composées suivantes sont :
\begin{center}
\renewcommand{\arraystretch}{2}
\begin{tabular}{|c|c|}
\hline
\textbf{Fonction $f(x)$} & \textbf{Primitive $F(x)$} \\
\hline
$u'(x) \cdot [u(x)]^n$ \quad ($n \neq -1$) & $\ds\frac{[u(x)]^{n+1}}{n+1}$ \\
\hline
$\ds\frac{u'(x)}{u(x)}$ & $\ln|u(x)|$ \\
\hline
$u'(x)\,e^{u(x)}$ & $e^{u(x)}$ \\
\hline
$u'(x)\cos u(x)$ & $\sin u(x)$ \\
\hline
$u'(x)\sin u(x)$ & $-\cos u(x)$ \\
\hline
\end{tabular}
\end{center}
\end{propriete}

\begin{methode}[Reconnaître une forme composée]
Pour primitiver $f(x)$, chercher si $f$ est de la forme $u' \times g(u)$ où $g$ est une fonction dont on connaît la primitive. Si $f(x) = \alpha\, u'(x)\, g(u(x))$, alors une primitive est $\alpha\, G(u(x))$ où $G$ est une primitive de $g$.
\end{methode}

\begin{exemple}[Calcul de primitives]
\begin{enumerate}
  \item Primitive de $f(x) = (2x+1)^3$ : on pose $u(x) = 2x+1$, $u'(x) = 2$. On a $f(x) = \frac{1}{2} \cdot 2 \cdot (2x+1)^3 = \frac{1}{2} u'u^3$. Primitive : $F(x) = \frac{1}{2} \cdot \frac{(2x+1)^4}{4} = \frac{(2x+1)^4}{8}$.
  \item Primitive de $g(x) = \frac{2x}{x^2+1}$ : on a $u = x^2+1$, $u' = 2x$, donc $g = \frac{u'}{u}$. Primitive : $G(x) = \ln(x^2+1)$. (Pas de valeur absolue car $x^2+1 > 0$.)
  \item Primitive de $h(x) = 3e^{3x}$ : on a $u = 3x$, $u' = 3$, donc $h = u'e^u$. Primitive : $H(x) = e^{3x}$.
  \item Primitive de $k(x) = \cos(5x)$ : on a $k(x) = \frac{1}{5} \cdot 5\cos(5x)$. Primitive : $K(x) = \frac{1}{5}\sin(5x)$.
\end{enumerate}
\end{exemple}

% ── III. Primitive avec condition initiale ──
\section{Primitive vérifiant une condition initiale}

\begin{propriete}[Unicité avec condition initiale]
Soit $f$ une fonction continue sur un intervalle $I$, et $x_0 \in I$, $y_0 \in \R$. Il existe une \textbf{unique} primitive $F$ de $f$ sur $I$ telle que $F(x_0) = y_0$.
\end{propriete}

\begin{methode}[Déterminer la primitive avec condition]
\begin{enumerate}
  \item Trouver la forme générale des primitives : $F(x) = G(x) + C$
  \item Utiliser la condition $F(x_0) = y_0$ pour calculer $C$
  \item En déduire l'expression de $F$
\end{enumerate}
\end{methode}

\begin{exemple}
Trouver la primitive $F$ de $f(x) = 2x + 3$ telle que $F(1) = 5$.

Forme générale : $F(x) = x^2 + 3x + C$.

Condition : $F(1) = 1 + 3 + C = 4 + C = 5$, donc $C = 1$.

Conclusion : $F(x) = x^2 + 3x + 1$.
\end{exemple}

% ── IV. Équation différentielle y' = ay ──
\section{Équation différentielle $y' = ay$}

\begin{definition}[Équation différentielle]
Une \textbf{équation différentielle} est une équation dont l'inconnue est une \textbf{fonction} $y$ et qui fait intervenir $y$ et sa dérivée $y'$ (voire $y''$, etc.). \textbf{Résoudre} une équation différentielle, c'est trouver toutes les fonctions $y$ qui la vérifient.
\end{definition}

\begin{theoreme}[Solutions de $y' = ay$]
Soit $a \in \R^*$. Les solutions de l'équation différentielle $y' = ay$ sont les fonctions :
\[ y(x) = Ce^{ax} \qquad \text{où } C \in \R \]
\end{theoreme}

\begin{demonstration}
\textbf{Vérification :} si $y(x) = Ce^{ax}$, alors $y'(x) = Cae^{ax} = a \cdot Ce^{ax} = ay(x)$. \checkmark

\textbf{Unicité :} Soit $y$ une solution et posons $g(x) = y(x)\,e^{-ax}$. Alors :
\[ g'(x) = y'(x)\,e^{-ax} + y(x)(-a)e^{-ax} = e^{-ax}[y'(x) - ay(x)] = 0 \]
car $y' = ay$. Donc $g$ est constante : $g(x) = C$, d'où $y(x) = Ce^{ax}$. \qed
\end{demonstration}

\begin{propriete}[Condition initiale]
L'équation $y' = ay$ avec la condition $y(x_0) = y_0$ admet une unique solution :
\[ y(x) = y_0\,e^{a(x-x_0)} \]
\end{propriete}

\begin{exemple}
Résoudre $y' = -3y$ avec $y(0) = 5$.

Solution générale : $y(x) = Ce^{-3x}$. Condition : $y(0) = C = 5$.

Solution : $y(x) = 5e^{-3x}$.
\end{exemple}

% ── V. Équation différentielle y' = ay + b ──
\section{Équation différentielle $y' = ay + b$}

\begin{methode}[Stratégie de résolution]
Pour résoudre $y' = ay + b$ (avec $a \neq 0$) :
\begin{enumerate}
  \item Trouver la \textbf{solution particulière constante} $y_p$ : si $y = k$ est constante, $y' = 0$, donc $0 = ak + b$, d'où $k = -\frac{b}{a}$.
  \item Résoudre l'\textbf{équation homogène} $y' = ay$ : solutions $Ce^{ax}$.
  \item La \textbf{solution générale} est $y(x) = Ce^{ax} - \frac{b}{a}$.
\end{enumerate}
\end{methode}

\begin{theoreme}[Solutions de $y' = ay + b$]
Soit $a \in \R^*$ et $b \in \R$. Les solutions de l'équation différentielle $y' = ay + b$ sont :
\[ y(x) = Ce^{ax} - \frac{b}{a} \qquad \text{où } C \in \R \]
\end{theoreme}

\begin{demonstration}
Posons $y_p = -\frac{b}{a}$ (solution constante) et $z(x) = y(x) - y_p = y(x) + \frac{b}{a}$.

Alors $z'(x) = y'(x)$ et $y' = ay + b$ s'écrit :
\[ z' = a\left(z - \frac{b}{a}\right) + b = az - b + b = az \]
Donc $z$ est solution de $z' = az$, ce qui donne $z(x) = Ce^{ax}$.

Finalement, $y(x) = z(x) + y_p = Ce^{ax} - \frac{b}{a}$. \qed
\end{demonstration}

\begin{exemple}[Résolution complète]
Résoudre $y' = -2y + 6$ avec $y(0) = 1$.

\textbf{Solution particulière constante :} $y_p = -\frac{6}{-2} = 3$.

\textbf{Solution générale :} $y(x) = Ce^{-2x} + 3$.

\textbf{Condition initiale :} $y(0) = C + 3 = 1$, donc $C = -2$.

\textbf{Solution :} $y(x) = -2e^{-2x} + 3$.

\medskip

\textit{Interprétation :} $y(x) \to 3$ quand $x \to +\infty$ (la solution tend vers la solution constante). C'est typique d'un phénomène qui se stabilise.
\end{exemple}

% ── VI. Applications ──
\section{Applications des équations différentielles}

\begin{exemple}[Décroissance radioactive]
Un isotope radioactif se désintègre selon la loi $N'(t) = -\lambda N(t)$, où $N(t)$ est le nombre de noyaux à l'instant $t$ et $\lambda > 0$ la constante de désintégration.

C'est une équation $y' = ay$ avec $a = -\lambda < 0$. Solution : $N(t) = N_0\,e^{-\lambda t}$.

La \textbf{demi-vie} $t_{1/2}$ est le temps au bout duquel la moitié des noyaux se sont désintégrés :
\[ N_0\,e^{-\lambda t_{1/2}} = \frac{N_0}{2} \implies e^{-\lambda t_{1/2}} = \frac{1}{2} \implies t_{1/2} = \frac{\ln 2}{\lambda} \]
\end{exemple}

\begin{exemple}[Loi de refroidissement de Newton]
La température $T(t)$ d'un corps plongé dans un milieu à température constante $T_m$ vérifie :
\[ T'(t) = -k(T(t) - T_m) \]
En posant $\theta(t) = T(t) - T_m$, on obtient $\theta' = -k\theta$, d'où $\theta(t) = (T_0 - T_m)e^{-kt}$.

Ainsi : $T(t) = T_m + (T_0 - T_m)\,e^{-kt}$.

La température tend exponentiellement vers $T_m$ : le corps se refroidit (ou se réchauffe) vers la température ambiante.
\end{exemple}

\begin{exemple}[Circuit RC]
Dans un circuit RC (résistance + condensateur), la tension $u(t)$ aux bornes du condensateur vérifie :
\[ u'(t) = -\frac{1}{RC}\,u(t) + \frac{E}{RC} \]
C'est une équation $y' = ay + b$ avec $a = -\frac{1}{RC}$ et $b = \frac{E}{RC}$.

Solution : $u(t) = Ce^{-t/RC} + E$. Avec $u(0) = 0$ : $u(t) = E\left(1 - e^{-t/RC}\right)$.

La constante de temps $\tau = RC$ caractérise la vitesse de charge : à $t = 5\tau$, le condensateur est chargé à $99{,}3\,\%$.
\end{exemple}

% ── Bilan ──
\section*{Bilan du chapitre}

\begin{retenir}
\begin{itemize}[leftmargin=*]
  \item \textbf{Primitive :} $F$ est une primitive de $f$ si $F' = f$
  \item \textbf{Unicité :} deux primitives de $f$ diffèrent d'une constante
  \item \textbf{Primitives usuelles :} $x^n \to \frac{x^{n+1}}{n+1}$, \quad $\frac{1}{x} \to \ln|x|$, \quad $e^x \to e^x$, \quad $\cos x \to \sin x$, \quad $\sin x \to -\cos x$
  \item \textbf{Composées :} $\frac{u'}{u} \to \ln|u|$, \quad $u'e^u \to e^u$, \quad $u'u^n \to \frac{u^{n+1}}{n+1}$
  \item \textbf{Condition initiale :} fixe la constante, donne une unique primitive
  \item \textbf{$y' = ay$ :} solutions $y = Ce^{ax}$
  \item \textbf{$y' = ay + b$ :} solutions $y = Ce^{ax} - \frac{b}{a}$ (solution particulière constante $-\frac{b}{a}$)
\end{itemize}
\end{retenir}


% ─────────────────────────────────────────────────────────────────────────────
% CHAPITRE 12 — CALCUL INTÉGRAL
% ─────────────────────────────────────────────────────────────────────────────

\chapter{Calcul intégral}

\section*{Introduction}

\textit{Comment calculer l'aire d'une région délimitée par une courbe ? Comment relier cette aire au calcul de primitives ? Comment calculer la valeur moyenne d'une grandeur qui varie ?}

\medskip

Le calcul intégral répond à ces questions en établissant un lien fondamental entre deux problèmes apparemment distincts : le calcul d'aires et la recherche de primitives. Ce lien, exprimé par le \textbf{théorème fondamental de l'analyse}, est l'une des plus belles réalisations des mathématiques. L'intégrale a des applications dans toutes les sciences : calcul de travail en physique, de doses en médecine, de probabilités, de moyennes, d'aires et de volumes.

\medskip

\textbf{Objectifs du chapitre :}
\begin{itemize}[leftmargin=*]
  \item Comprendre la notion d'intégrale comme aire algébrique
  \item Connaître et appliquer le théorème fondamental de l'analyse
  \item Maîtriser les propriétés de l'intégrale (linéarité, Chasles, positivité)
  \item Calculer des intégrales par recherche de primitives
  \item Calculer une valeur moyenne
  \item Utiliser l'intégration par parties
  \item Calculer des aires entre deux courbes
\end{itemize}

\begin{histoire}[Du calcul d'aires à l'intégrale]
Le problème du calcul d'aires remonte à l'Antiquité. \textbf{Archimède} (III\ieme{} siècle av. J.-C.) calcule l'aire sous une parabole par une méthode d'exhaustion, anticipant l'intégrale de vingt siècles. Au XVII\ieme{} siècle, \textbf{Cavalieri} développe la méthode des indivisibles. Mais c'est le \textbf{théorème fondamental}, établi indépendamment par \textbf{Newton} et \textbf{Leibniz} vers 1670, qui révolutionne le sujet : il montre que calculer une aire revient à trouver une primitive. Leibniz introduit le symbole $\int$ (un « S » allongé, pour « somme ») et la notation $\int_a^b f(x)\,\mathrm{d}x$, toujours utilisée aujourd'hui.

\medskip

Au XIX\ieme{} siècle, \textbf{Riemann} donne une définition rigoureuse de l'intégrale comme limite de sommes de rectangles, puis \textbf{Lebesgue} (1902) en propose une généralisation puissante qui fonde la théorie moderne de la mesure.
\end{histoire}

% ── I. Intégrale d'une fonction continue positive ──
\section{Intégrale d'une fonction continue positive}

\subsection{Aire sous une courbe}

\begin{definition}[Intégrale d'une fonction continue positive]
Soit $f$ une fonction continue et \textbf{positive} sur un intervalle $[a\,;b]$. L'\textbf{intégrale de $f$ de $a$ à $b$}, notée $\ds\int_a^b f(x)\,\mathrm{d}x$, est l'\textbf{aire} (en unités d'aire) de la surface délimitée par :
\begin{itemize}
  \item la courbe $\mathcal{C}_f$
  \item l'axe des abscisses
  \item les droites verticales $x = a$ et $x = b$
\end{itemize}
\end{definition}

\begin{center}
\begin{tikzpicture}
\begin{axis}[
  width=10cm, height=6cm,
  axis lines=center,
  xlabel={$x$}, ylabel={$y$},
  xmin=-0.5, xmax=5, ymin=-0.5, ymax=4,
  xtick={1,3.5},
  xticklabels={$a$,$b$},
  ytick=\empty,
  grid=none
]
\addplot[domain=0.3:4.5, samples=100, thick, indigo, name path=f] {0.3*x^2 - 0.8*x + 2.5};
\addplot[domain=0.3:4.5, samples=100, draw=none, name path=zero] {0};
\addplot[fill=indigobg, fill opacity=0.6] fill between[of=f and zero, soft clip={domain=1:3.5}];
\draw[dashed, gray] (axis cs:1,0) -- (axis cs:1,{0.3*1-0.8+2.5});
\draw[dashed, gray] (axis cs:3.5,0) -- (axis cs:3.5,{0.3*12.25-2.8+2.5});
\node[indigo] at (axis cs:2.25,1.2) {\large $\ds\int_a^b f(x)\,\mathrm{d}x$};
\node[indigo, above] at (axis cs:4, 3.5) {$\mathcal{C}_f$};
\end{axis}
\end{tikzpicture}
\captionof{figure}{L'intégrale de $f$ positive est l'aire sous la courbe}
\end{center}

\textit{L'unité d'aire (u.a.) est l'aire du rectangle formé par une unité sur chaque axe.}

\subsection{Approche par les sommes de Riemann}

\begin{methode}[Intuition des sommes de Riemann]
On approche l'aire sous la courbe en la découpant en $n$ rectangles de largeur $\frac{b-a}{n}$. La somme des aires de ces rectangles :
\[ S_n = \frac{b-a}{n}\sum_{k=0}^{n-1} f\left(a + k\cdot\frac{b-a}{n}\right) \]
converge vers $\ds\int_a^b f(x)\,\mathrm{d}x$ quand $n \to +\infty$.
\end{methode}

% ── II. Théorème fondamental ──
\section{Théorème fondamental de l'analyse}

\begin{theoreme}[Théorème fondamental]
Soit $f$ une fonction continue sur un intervalle $[a\,;b]$ et $F$ une primitive de $f$ sur $[a\,;b]$. Alors :
\[ \int_a^b f(x)\,\mathrm{d}x = F(b) - F(a) = \Big[F(x)\Big]_a^b \]
\end{theoreme}

\textit{Ce théorème est le pilier du calcul intégral : il ramène le calcul d'une intégrale (un problème d'aire, a priori difficile) au calcul d'une primitive (un problème algébrique).}

\begin{demonstration}[Idée de la démonstration]
Posons $\Phi(x) = \int_a^x f(t)\,\mathrm{d}t$ pour $x \in [a\,;b]$. On montre que $\Phi$ est dérivable et $\Phi'(x) = f(x)$. Donc $\Phi$ est une primitive de $f$. Il existe une constante $C$ telle que $\Phi(x) = F(x) + C$.

Or $\Phi(a) = \int_a^a f(t)\,\mathrm{d}t = 0$, donc $0 = F(a) + C$, d'où $C = -F(a)$.

Alors $\Phi(b) = \int_a^b f(t)\,\mathrm{d}t = F(b) + C = F(b) - F(a)$. \qed
\end{demonstration}

\begin{methode}[Calculer une intégrale]
\begin{enumerate}
  \item Trouver une primitive $F$ de $f$
  \item Calculer $F(b) - F(a)$
  \item Notation : $\Big[F(x)\Big]_a^b = F(b) - F(a)$
\end{enumerate}
\end{methode}

\begin{exemple}[Calculs d'intégrales]
\begin{enumerate}
  \item $\ds\int_0^3 2x\,\mathrm{d}x = \Big[x^2\Big]_0^3 = 9 - 0 = 9$
  \item $\ds\int_1^e \frac{1}{x}\,\mathrm{d}x = \Big[\ln x\Big]_1^e = \ln e - \ln 1 = 1 - 0 = 1$
  \item $\ds\int_0^1 e^x\,\mathrm{d}x = \Big[e^x\Big]_0^1 = e^1 - e^0 = e - 1$
  \item $\ds\int_0^{\pi} \sin x\,\mathrm{d}x = \Big[-\cos x\Big]_0^{\pi} = -\cos\pi - (-\cos 0) = -(-1) + 1 = 2$
  \item $\ds\int_1^4 (3x^2 - 2x + 1)\,\mathrm{d}x = \Big[x^3 - x^2 + x\Big]_1^4 = (64 - 16 + 4) - (1 - 1 + 1) = 52 - 1 = 51$
\end{enumerate}
\end{exemple}

% ── III. Propriétés de l'intégrale ──
\section{Propriétés de l'intégrale}

\subsection{Conventions et relation de Chasles}

\begin{propriete}[Conventions]
Pour toute fonction $f$ continue et tout réel $a$ :
\begin{enumerate}
  \item $\ds\int_a^a f(x)\,\mathrm{d}x = 0$
  \item $\ds\int_b^a f(x)\,\mathrm{d}x = -\int_a^b f(x)\,\mathrm{d}x$
\end{enumerate}
\end{propriete}

\begin{theoreme}[Relation de Chasles]
Pour toute fonction $f$ continue et tous réels $a$, $b$, $c$ :
\[ \int_a^b f(x)\,\mathrm{d}x + \int_b^c f(x)\,\mathrm{d}x = \int_a^c f(x)\,\mathrm{d}x \]
\end{theoreme}

\begin{demonstration}
$\ds\int_a^b f + \int_b^c f = [F(b)-F(a)] + [F(c)-F(b)] = F(c) - F(a) = \int_a^c f$. \qed
\end{demonstration}

\subsection{Linéarité}

\begin{propriete}[Linéarité de l'intégrale]
Pour toutes fonctions continues $f$, $g$ sur $[a\,;b]$ et tout réel $\lambda$ :
\begin{enumerate}
  \item $\ds\int_a^b [f(x) + g(x)]\,\mathrm{d}x = \int_a^b f(x)\,\mathrm{d}x + \int_a^b g(x)\,\mathrm{d}x$
  \item $\ds\int_a^b \lambda\, f(x)\,\mathrm{d}x = \lambda\int_a^b f(x)\,\mathrm{d}x$
\end{enumerate}
\end{propriete}

\subsection{Positivité et ordre}

\begin{propriete}[Positivité]
Si $f$ est continue et \textbf{positive} sur $[a\,;b]$ (avec $a \leqslant b$), alors :
\[ \int_a^b f(x)\,\mathrm{d}x \geqslant 0 \]
\end{propriete}

\begin{propriete}[Ordre]
Si $f$ et $g$ sont continues sur $[a\,;b]$ avec $a \leqslant b$ et $f(x) \leqslant g(x)$ pour tout $x \in [a\,;b]$, alors :
\[ \int_a^b f(x)\,\mathrm{d}x \leqslant \int_a^b g(x)\,\mathrm{d}x \]
\end{propriete}

\begin{attention}[L'ordre n'est vrai que si $a \leqslant b$]
Les propriétés de positivité et d'ordre supposent $a \leqslant b$. Si $a > b$, l'inégalité s'inverse (car $\int_b^a f = -\int_a^b f$).
\end{attention}

% ── IV. Fonction de signe variable ──
\section{Intégrale d'une fonction de signe quelconque}

\subsection{Aire algébrique}

\begin{definition}[Aire algébrique]
Si $f$ prend des valeurs positives et négatives sur $[a\,;b]$, l'intégrale $\ds\int_a^b f(x)\,\mathrm{d}x$ est une \textbf{aire algébrique} :
\begin{itemize}
  \item Les aires au-dessus de l'axe des abscisses comptent \textbf{positivement}
  \item Les aires en dessous de l'axe des abscisses comptent \textbf{négativement}
\end{itemize}
\end{definition}

\begin{attention}[Intégrale et aire totale]
L'intégrale $\ds\int_a^b f(x)\,\mathrm{d}x$ n'est \textbf{pas} en général l'aire totale de la surface entre la courbe et l'axe. Si $f$ change de signe, des parties se compensent.
\end{attention}

\subsection{Calcul de l'aire totale}

\begin{methode}[Aire totale entre une courbe et l'axe des abscisses]
Pour calculer l'aire totale de la surface entre $\mathcal{C}_f$ et l'axe $Ox$ sur $[a\,;b]$, on intègre la \textbf{valeur absolue} de $f$ :
\[ \mathcal{A} = \int_a^b |f(x)|\,\mathrm{d}x \]
En pratique :
\begin{enumerate}
  \item Chercher les zéros de $f$ sur $[a\,;b]$ : $x_1, x_2, \ldots$
  \item Découper $[a\,;b]$ en sous-intervalles où $f$ garde un signe constant
  \item Sur chaque sous-intervalle, calculer $\int |f|$ (= $\int f$ si $f \geqslant 0$, $= -\int f$ si $f \leqslant 0$)
  \item Sommer les résultats
\end{enumerate}
\end{methode}

\begin{exemple}
Calculer l'aire entre la courbe de $f(x) = x^2 - 1$ et l'axe $Ox$ sur $[-2\,;2]$.

$f$ s'annule en $x = -1$ et $x = 1$. Sur $[-2\,;-1]$ et $[1\,;2]$, $f \geqslant 0$. Sur $[-1\,;1]$, $f \leqslant 0$.

\begin{align*}
\mathcal{A} &= \int_{-2}^{-1} (x^2-1)\,\mathrm{d}x - \int_{-1}^{1} (x^2-1)\,\mathrm{d}x + \int_{1}^{2} (x^2-1)\,\mathrm{d}x \\
&= \left[\frac{x^3}{3}-x\right]_{-2}^{-1} - \left[\frac{x^3}{3}-x\right]_{-1}^{1} + \left[\frac{x^3}{3}-x\right]_{1}^{2} \\
&= \frac{4}{3} + \frac{4}{3} + \frac{4}{3} = 4
\end{align*}
\end{exemple}

% ── V. Valeur moyenne ──
\section{Valeur moyenne d'une fonction}

\begin{definition}[Valeur moyenne]
Soit $f$ une fonction continue sur $[a\,;b]$ avec $a < b$. La \textbf{valeur moyenne} de $f$ sur $[a\,;b]$ est :
\[ \mu = \frac{1}{b-a}\int_a^b f(x)\,\mathrm{d}x \]
\end{definition}

\textit{Interprétation géométrique :} $\mu$ est la hauteur du rectangle de base $[a\,;b]$ qui a la même aire que la surface sous la courbe de $f$. Autrement dit, $\mu \times (b-a) = \int_a^b f(x)\,\mathrm{d}x$.

\begin{exemple}
Valeur moyenne de $f(x) = x^2$ sur $[0\,;3]$ :
\[ \mu = \frac{1}{3-0}\int_0^3 x^2\,\mathrm{d}x = \frac{1}{3}\left[\frac{x^3}{3}\right]_0^3 = \frac{1}{3} \cdot \frac{27}{3} = \frac{1}{3} \cdot 9 = 3 \]
\end{exemple}

\begin{exemple}[Application physique]
La température (en °C) au cours d'une journée de 12 h est modélisée par $T(t) = -t^2 + 12t + 5$ pour $t \in [0\,;12]$. La température moyenne est :
\[ T_{\text{moy}} = \frac{1}{12}\int_0^{12} (-t^2 + 12t + 5)\,\mathrm{d}t = \frac{1}{12}\left[-\frac{t^3}{3}+6t^2+5t\right]_0^{12} = \frac{1}{12}\left(-576+864+60\right) = \frac{348}{12} = 29\text{ °C} \]
\end{exemple}

% ── VI. Intégration par parties ──
\section{Intégration par parties}

\begin{theoreme}[Intégration par parties (IPP)]
Si $u$ et $v$ sont deux fonctions dérivables sur $[a\,;b]$ avec $u'$ et $v'$ continues, alors :
\[ \int_a^b u(x)\,v'(x)\,\mathrm{d}x = \Big[u(x)\,v(x)\Big]_a^b - \int_a^b u'(x)\,v(x)\,\mathrm{d}x \]
\end{theoreme}

\begin{demonstration}
La formule de dérivation d'un produit donne $(uv)' = u'v + uv'$, d'où $uv' = (uv)' - u'v$.

En intégrant de $a$ à $b$ :
\[ \int_a^b uv' = \int_a^b (uv)' - \int_a^b u'v = \Big[uv\Big]_a^b - \int_a^b u'v \qedhere \]
\end{demonstration}

\begin{methode}[Choix de $u$ et $v'$ pour l'IPP]
Le but est de \textbf{simplifier} l'intégrale. Règle générale :
\begin{itemize}
  \item Choisir $u$ = la fonction qui se \textbf{simplifie en dérivant} (polynôme, $\ln x$)
  \item Choisir $v'$ = la fonction qu'on sait \textbf{primitiver facilement} ($e^x$, $\cos x$, $\sin x$)
\end{itemize}
Moyen mnémotechnique : \textbf{LATE} — Logarithme, Algébrique (polynôme), Trigonométrique, Exponentielle. On choisit $u$ en priorité dans cet ordre.
\end{methode}

\begin{exemple}[IPP classiques]
\begin{enumerate}
  \item $\ds\int_0^1 xe^x\,\mathrm{d}x$. On pose $u = x$, $v' = e^x$, donc $u' = 1$, $v = e^x$.
  \[ \int_0^1 xe^x\,\mathrm{d}x = \Big[xe^x\Big]_0^1 - \int_0^1 e^x\,\mathrm{d}x = e - \Big[e^x\Big]_0^1 = e - (e - 1) = 1 \]

  \item $\ds\int_1^e \ln x\,\mathrm{d}x$. On pose $u = \ln x$, $v' = 1$, donc $u' = \frac{1}{x}$, $v = x$.
  \[ \int_1^e \ln x\,\mathrm{d}x = \Big[x\ln x\Big]_1^e - \int_1^e 1\,\mathrm{d}x = (e \cdot 1 - 1 \cdot 0) - (e-1) = e - e + 1 = 1 \]

  \item $\ds\int_0^{\pi} x\sin x\,\mathrm{d}x$. On pose $u = x$, $v' = \sin x$, donc $u' = 1$, $v = -\cos x$.
  \[ \int_0^{\pi} x\sin x\,\mathrm{d}x = \Big[-x\cos x\Big]_0^{\pi} + \int_0^{\pi} \cos x\,\mathrm{d}x = \pi + \Big[\sin x\Big]_0^{\pi} = \pi + 0 = \pi \]
\end{enumerate}
\end{exemple}

% ── VII. Aire entre deux courbes ──
\section{Aire entre deux courbes}

\begin{propriete}[Aire entre deux courbes]
Soient $f$ et $g$ deux fonctions continues sur $[a\,;b]$. L'aire de la surface comprise entre les courbes $\mathcal{C}_f$ et $\mathcal{C}_g$ et les droites $x = a$, $x = b$ est :
\[ \mathcal{A} = \int_a^b |f(x) - g(x)|\,\mathrm{d}x \]
\end{propriete}

\begin{methode}[Calculer l'aire entre deux courbes]
\begin{enumerate}
  \item Résoudre $f(x) = g(x)$ pour trouver les points d'intersection.
  \item Déterminer le signe de $f(x) - g(x)$ sur chaque sous-intervalle.
  \item Sur chaque sous-intervalle $[x_i, x_{i+1}]$ :
  \begin{itemize}
    \item Si $f \geqslant g$ : ajouter $\ds\int_{x_i}^{x_{i+1}} [f(x) - g(x)]\,\mathrm{d}x$
    \item Si $f \leqslant g$ : ajouter $\ds\int_{x_i}^{x_{i+1}} [g(x) - f(x)]\,\mathrm{d}x$
  \end{itemize}
\end{enumerate}
\end{methode}

\begin{exemple}
Calculer l'aire entre les courbes de $f(x) = x^2$ et $g(x) = x$ sur $[0\,;1]$.

Intersection : $x^2 = x \iff x(x-1) = 0 \iff x = 0$ ou $x = 1$.

Sur $]0\,;1[$, $x > x^2$ (car $0 < x < 1$), donc $g(x) > f(x)$.

\[ \mathcal{A} = \int_0^1 (x - x^2)\,\mathrm{d}x = \left[\frac{x^2}{2} - \frac{x^3}{3}\right]_0^1 = \frac{1}{2} - \frac{1}{3} = \frac{1}{6} \text{ u.a.} \]
\end{exemple}

\begin{exemple}[Aire entre une parabole et une droite]
Calculer l'aire entre $f(x) = x^2 - 2x$ et $g(x) = x$ sur leur domaine d'intersection.

$f(x) = g(x) \iff x^2 - 2x = x \iff x^2 - 3x = 0 \iff x(x-3) = 0$. Points d'intersection : $x = 0$ et $x = 3$.

Sur $]0\,;3[$, $g(x) - f(x) = x - x^2 + 2x = 3x - x^2 = x(3-x) > 0$.

\[ \mathcal{A} = \int_0^3 (3x - x^2)\,\mathrm{d}x = \left[\frac{3x^2}{2} - \frac{x^3}{3}\right]_0^3 = \frac{27}{2} - 9 = \frac{27}{2} - \frac{18}{2} = \frac{9}{2} \text{ u.a.} \]
\end{exemple}

% ── VIII. Compléments ──
\section{Compléments}

\subsection{Suites et intégrales}

\begin{exemple}[Suite définie par une intégrale]
Soit $I_n = \ds\int_0^1 x^n\,\mathrm{d}x$ pour $n \in \N$.

\textbf{Calcul :} $I_n = \left[\frac{x^{n+1}}{n+1}\right]_0^1 = \frac{1}{n+1}$.

\textbf{Limite :} $\ds\lim_{n \to +\infty} I_n = 0$. Interprétation : la courbe de $x^n$ sur $[0\,;1]$ « s'écrase » vers 0 quand $n$ grandit (sauf en $x = 1$).
\end{exemple}

\begin{exemple}[Relation de récurrence par IPP]
Soit $J_n = \ds\int_0^1 x^n e^x\,\mathrm{d}x$. Par IPP avec $u = x^n$ et $v' = e^x$ :
\[ J_n = \Big[x^n e^x\Big]_0^1 - \int_0^1 nx^{n-1}e^x\,\mathrm{d}x = e - nJ_{n-1} \]
Avec $J_0 = \int_0^1 e^x\,\mathrm{d}x = e - 1$, on peut calculer $J_n$ de proche en proche.
\end{exemple}

\subsection{Application en probabilités}

\begin{exemple}[Loi uniforme]
Si une variable aléatoire $X$ suit une loi uniforme sur $[a\,;b]$, sa densité est $f(x) = \frac{1}{b-a}$ et son espérance est :
\[ E(X) = \int_a^b x \cdot \frac{1}{b-a}\,\mathrm{d}x = \frac{1}{b-a}\left[\frac{x^2}{2}\right]_a^b = \frac{b^2-a^2}{2(b-a)} = \frac{a+b}{2} \]
L'espérance est le milieu de l'intervalle.
\end{exemple}

% ── Bilan ──
\section*{Bilan du chapitre}

\begin{retenir}
\begin{itemize}[leftmargin=*]
  \item \textbf{Théorème fondamental :} $\ds\int_a^b f(x)\,\mathrm{d}x = F(b) - F(a) = \Big[F(x)\Big]_a^b$
  \item \textbf{Relation de Chasles :} $\ds\int_a^b f + \int_b^c f = \int_a^c f$
  \item \textbf{Linéarité :} $\ds\int_a^b (\alpha f + \beta g) = \alpha\int_a^b f + \beta\int_a^b g$
  \item \textbf{Positivité :} si $f \geqslant 0$ et $a \leqslant b$, alors $\ds\int_a^b f \geqslant 0$
  \item \textbf{Valeur moyenne :} $\mu = \ds\frac{1}{b-a}\int_a^b f(x)\,\mathrm{d}x$
  \item \textbf{IPP :} $\ds\int_a^b uv' = \Big[uv\Big]_a^b - \int_a^b u'v$
  \item \textbf{Aire entre deux courbes :} $\mathcal{A} = \ds\int_a^b |f(x) - g(x)|\,\mathrm{d}x$
  \item \textbf{Conventions :} $\ds\int_a^a f = 0$ et $\ds\int_b^a f = -\int_a^b f$
\end{itemize}
\end{retenir}

\part{Probabilités}

% ─────────────────────────────────────────────────────────────────────────────
% CHAPITRE 13 — ÉPREUVES INDÉPENDANTES ET LOI BINOMIALE
% ─────────────────────────────────────────────────────────────────────────────

\chapter{Épreuves indépendantes et loi binomiale}

\section*{Introduction}

\textit{Un fabricant produit 10\,000 pièces par jour. On sait que 2\,\% d'entre elles sont défectueuses. On prélève un échantillon de 50 pièces au hasard : combien peut-on s'attendre à en trouver de défectueuses ? Quelle est la probabilité de n'en trouver aucune ? Au moins 3 ?}

\medskip

Ces questions apparaissent dans le \textbf{contrôle qualité industriel}, les \textbf{tests médicaux} (probabilité de faux positifs), les \textbf{sondages} et même les jeux de hasard. Elles ont toutes un point commun : on \textbf{répète} une même expérience aléatoire, de manière \textbf{indépendante}, et on compte le nombre de « succès ». Le modèle mathématique qui répond à ces questions est la \textbf{loi binomiale}, l'une des lois de probabilité les plus importantes.

\medskip

\textbf{Objectifs du chapitre :}
\begin{itemize}[leftmargin=*]
  \item Reconnaître des épreuves indépendantes et un schéma de Bernoulli
  \item Connaître et utiliser la loi binomiale $\mathcal{B}(n,p)$
  \item Calculer des probabilités, l'espérance et l'écart-type
  \item Résoudre des problèmes de seuil
\end{itemize}

\begin{histoire}[Jacques Bernoulli et l'\textit{Ars Conjectandi}]
La loi binomiale porte le nom du mathématicien suisse \textbf{Jacques Bernoulli} (1654--1705). Son ouvrage \textit{Ars Conjectandi} (« L'Art de conjecturer »), publié en 1713 après sa mort, est considéré comme l'acte de naissance de la \textbf{théorie des probabilités} en tant que discipline mathématique. Bernoulli y développe systématiquement le calcul des combinaisons et formule la loi des grands nombres sous sa forme originelle. La famille Bernoulli, installée à Bâle, a produit au moins huit mathématiciens remarquables sur trois générations --- un cas unique dans l'histoire des sciences.
\end{histoire}

% ── I. Indépendance ──
\section{Indépendance de deux événements}

\subsection{Définition}

\begin{definition}[Événements indépendants]
Deux événements $A$ et $B$ de probabilités non nulles sont \textbf{indépendants} si la réalisation de l'un ne modifie pas la probabilité de l'autre. Formellement :
\[ P(A \cap B) = P(A) \times P(B) \]
\end{definition}

\textit{Intuition :} connaître le résultat de $A$ n'apporte aucune information sur $B$, et réciproquement. Par exemple, les résultats de deux lancers de dé successifs sont indépendants.

\begin{attention}[Indépendance $\neq$ incompatibilité]
Deux événements \textbf{incompatibles} ($A \cap B = \emptyset$) ne sont \textbf{jamais} indépendants (sauf cas trivial de probabilité nulle). En effet, si $A$ est réalisé, on sait que $B$ ne l'est pas : la réalisation de l'un influence bien l'autre.
\end{attention}

\begin{exemple}[Lancer de deux dés]
On lance deux dés équilibrés. Soit $A$ : « le premier dé donne 6 » et $B$ : « le second dé donne un nombre pair ».

$P(A) = \dfrac{1}{6}$, \quad $P(B) = \dfrac{1}{2}$, \quad $P(A \cap B) = \dfrac{1}{6} \times \dfrac{1}{2} = \dfrac{1}{12}$.

On vérifie : $P(A) \times P(B) = \dfrac{1}{12} = P(A \cap B)$, donc $A$ et $B$ sont indépendants.
\end{exemple}

\subsection{Succession d'épreuves indépendantes}

\begin{propriete}[Généralisation]
Si $A_1, A_2, \ldots, A_n$ sont des événements \textbf{mutuellement indépendants}, alors :
\[ P(A_1 \cap A_2 \cap \cdots \cap A_n) = P(A_1) \times P(A_2) \times \cdots \times P(A_n) \]
\end{propriete}

\begin{methode}[Vérifier l'indépendance en pratique]
En pratique, l'indépendance est souvent \textbf{postulée} par le modèle (tirage avec remise, expériences physiquement séparées). On ne la « démontre » pas à partir des données ; c'est une hypothèse de modélisation.
\end{methode}

% ── II. Épreuve de Bernoulli ──
\section{Épreuve de Bernoulli}

\begin{definition}[Épreuve de Bernoulli]
Une \textbf{épreuve de Bernoulli} est une expérience aléatoire n'admettant que \textbf{deux issues} :
\begin{itemize}
  \item le \textbf{succès} $S$, de probabilité $p$ ;
  \item l'\textbf{échec} $\bar{S}$, de probabilité $q = 1 - p$.
\end{itemize}
Le paramètre $p \in [0, 1]$ est appelé \textbf{paramètre} de l'épreuve de Bernoulli.
\end{definition}

\begin{exemple}
\begin{itemize}
  \item Lancer une pièce : $S$ = « obtenir pile », $p = 0{,}5$.
  \item Contrôle qualité : $S$ = « pièce défectueuse », $p = 0{,}02$.
  \item Test médical : $S$ = « test positif », $p$ dépend de la sensibilité du test.
\end{itemize}
\end{exemple}

\begin{definition}[Variable de Bernoulli]
La \textbf{variable aléatoire de Bernoulli} $X$ associée vaut :
\[ X = \begin{cases} 1 & \text{en cas de succès (probabilité } p\text{)} \\ 0 & \text{en cas d'échec (probabilité } 1-p\text{)} \end{cases} \]
On note $X \sim \mathcal{B}(p)$. Son espérance est $E(X) = p$ et sa variance $V(X) = p(1-p)$.
\end{definition}

% ── III. Schéma de Bernoulli ──
\section{Schéma de Bernoulli}

\begin{definition}[Schéma de Bernoulli]
Un \textbf{schéma de Bernoulli} de paramètres $n$ et $p$ consiste en la \textbf{répétition de $n$ épreuves de Bernoulli identiques et indépendantes}, chacune de paramètre $p$.

Les trois conditions sont essentielles :
\begin{enumerate}
  \item Chaque épreuve n'a que \textbf{deux issues} (succès ou échec).
  \item La probabilité de succès $p$ est la \textbf{même} à chaque épreuve.
  \item Les épreuves sont \textbf{indépendantes} les unes des autres.
\end{enumerate}
\end{definition}

\begin{attention}[Tirage sans remise]
Un tirage \textbf{sans remise} dans une urne ne constitue \textbf{pas} un schéma de Bernoulli, car la probabilité change à chaque tirage et les tirages ne sont pas indépendants. Toutefois, si la population est très grande par rapport à l'échantillon, on peut \textbf{assimiler} un tirage sans remise à un schéma de Bernoulli.
\end{attention}

\begin{exemple}[Arbre de probabilité pour $n = 3$]
On lance 3 fois une pièce truquée avec $p = P(S) = 0{,}6$ et $q = P(\bar{S}) = 0{,}4$.

\begin{center}
\begin{tikzpicture}[
  grow=right,
  level distance=2.8cm,
  sibling distance=1.2cm,
  every node/.style={font=\small},
  edge from parent/.style={draw,indigo,-{Stealth[length=3pt]}},
  level 1/.style={sibling distance=3.2cm},
  level 2/.style={sibling distance=1.6cm},
  level 3/.style={sibling distance=0.8cm}
]
\node[indigo,font=\bfseries] {Départ}
  child {node {$\bar{S}$}
    child {node {$\bar{S}$}
      child {node {$\bar{S}$} edge from parent node[below,font=\scriptsize] {$0{,}4$}}
      child {node {$S$} edge from parent node[above,font=\scriptsize] {$0{,}6$}}
    edge from parent node[below,font=\scriptsize] {$0{,}4$}}
    child {node {$S$}
      child {node {$\bar{S}$} edge from parent node[below,font=\scriptsize] {$0{,}4$}}
      child {node {$S$} edge from parent node[above,font=\scriptsize] {$0{,}6$}}
    edge from parent node[above,font=\scriptsize] {$0{,}6$}}
  edge from parent node[below,font=\scriptsize] {$0{,}4$}}
  child {node {$S$}
    child {node {$\bar{S}$}
      child {node {$\bar{S}$} edge from parent node[below,font=\scriptsize] {$0{,}4$}}
      child {node {$S$} edge from parent node[above,font=\scriptsize] {$0{,}6$}}
    edge from parent node[below,font=\scriptsize] {$0{,}4$}}
    child {node {$S$}
      child {node {$\bar{S}$} edge from parent node[below,font=\scriptsize] {$0{,}4$}}
      child {node {$S$} edge from parent node[above,font=\scriptsize] {$0{,}6$}}
    edge from parent node[above,font=\scriptsize] {$0{,}6$}}
  edge from parent node[above,font=\scriptsize] {$0{,}6$}};
\end{tikzpicture}
\captionof{figure}{Arbre de probabilité d'un schéma de Bernoulli avec $n=3$ et $p=0{,}6$}
\end{center}

Sur cet arbre, chaque chemin de longueur 3 correspond à une issue. Par exemple, le chemin $S$-$S$-$\bar{S}$ (2 succès puis 1 échec) a pour probabilité $0{,}6 \times 0{,}6 \times 0{,}4 = 0{,}144$.

Il y a $\binom{3}{2} = 3$ chemins contenant exactement 2 succès ($SS\bar{S}$, $S\bar{S}S$, $\bar{S}SS$), donc $P(X = 2) = 3 \times 0{,}6^2 \times 0{,}4 = 0{,}432$.
\end{exemple}

% ── IV. Loi binomiale ──
\section{Loi binomiale}

\begin{definition}[Loi binomiale]
Soit $X$ le nombre de succès dans un schéma de Bernoulli de paramètres $n$ et $p$. La variable $X$ suit la \textbf{loi binomiale} de paramètres $n$ et $p$, notée :
\[ X \sim \mathcal{B}(n, p) \]
$X$ prend les valeurs $0, 1, 2, \ldots, n$.
\end{definition}

\begin{propriete}[Formule fondamentale]
Si $X \sim \mathcal{B}(n, p)$, alors pour tout $k \in \{0, 1, \ldots, n\}$ :
\[ \boxed{P(X = k) = \binom{n}{k}\, p^k\, (1-p)^{n-k}} \]
\end{propriete}

\begin{methode}[Comprendre chaque facteur]
La formule se décompose ainsi :
\begin{itemize}
  \item $\binom{n}{k}$ : nombre de façons de \textbf{placer} les $k$ succès parmi les $n$ épreuves (combinaisons) ;
  \item $p^k$ : probabilité d'obtenir $k$ \textbf{succès} (chacun de probabilité $p$) ;
  \item $(1-p)^{n-k}$ : probabilité d'obtenir $n-k$ \textbf{échecs} (chacun de probabilité $1-p$).
\end{itemize}
Le produit $p^k (1-p)^{n-k}$ est la probabilité d'\textbf{un} chemin particulier contenant $k$ succès et $n-k$ échecs. On multiplie par $\binom{n}{k}$ car il y a autant de chemins de ce type.
\end{methode}

\begin{exemple}[Contrôle qualité]
Une usine produit des pièces dont 3\,\% sont défectueuses. On prélève $n = 10$ pièces (avec remise). Soit $X$ le nombre de pièces défectueuses : $X \sim \mathcal{B}(10 ;\, 0{,}03)$.

\medskip

$P(X = 0) = \binom{10}{0}\, (0{,}03)^0\, (0{,}97)^{10} = (0{,}97)^{10} \approx 0{,}7374$

\medskip

$P(X = 1) = \binom{10}{1}\, (0{,}03)^1\, (0{,}97)^{9} = 10 \times 0{,}03 \times (0{,}97)^9 \approx 0{,}2281$

\medskip

$P(X = 2) = \binom{10}{2}\, (0{,}03)^2\, (0{,}97)^{8} = 45 \times 0{,}0009 \times (0{,}97)^8 \approx 0{,}0317$
\end{exemple}

% ── V. Espérance, variance, écart-type ──
\section{Espérance, variance et écart-type}

\begin{propriete}[Paramètres de la loi binomiale]
Si $X \sim \mathcal{B}(n, p)$, alors :
\begin{align*}
E(X) &= np \\
V(X) &= np(1-p) \\
\sigma(X) &= \sqrt{np(1-p)}
\end{align*}
\end{propriete}

\textit{Interprétation :} si l'on répète $n$ fois une expérience de probabilité de succès $p$, on s'attend \textbf{en moyenne} à observer $np$ succès.

\begin{demonstration}[de $E(X) = np$ --- admise ici, démontrée au chapitre~14]
La démonstration repose sur l'écriture $X = X_1 + X_2 + \cdots + X_n$ où chaque $X_i \sim \mathcal{B}(p)$, et sur la linéarité de l'espérance. Voir chapitre~14.
\end{demonstration}

\begin{exemple}
On lance 100 fois un dé équilibré. Soit $X$ le nombre de 6 obtenus : $X \sim \mathcal{B}(100 ;\, \tfrac{1}{6})$.

$E(X) = 100 \times \dfrac{1}{6} \approx 16{,}7$, \quad $\sigma(X) = \sqrt{100 \times \dfrac{1}{6} \times \dfrac{5}{6}} \approx 3{,}73$.

En moyenne, on obtient environ 17 fois le 6, avec un écart-type d'environ 3,7.
\end{exemple}

% ── VI. Calculs pratiques ──
\section{Calculs pratiques de probabilités cumulées}

\begin{methode}[Probabilités cumulées]
Pour $X \sim \mathcal{B}(n, p)$ :
\begin{itemize}
  \item $P(X \leqslant k) = \ds\sum_{i=0}^{k} \binom{n}{i}\, p^i\, (1-p)^{n-i}$
  \item $P(X \geqslant k) = 1 - P(X \leqslant k-1)$
  \item $P(X \geqslant 1) = 1 - P(X = 0) = 1 - (1-p)^n$
\end{itemize}
En pratique, on utilise la calculatrice ou un tableur pour ces calculs (fonction \texttt{binomCdf} ou \texttt{LOI.BINOMIALE}).
\end{methode}

\begin{exemple}[Au moins un succès]
Un contrôleur inspecte $n = 20$ pièces, la probabilité de défaut étant $p = 0{,}05$. Quelle est la probabilité de trouver au moins une pièce défectueuse ?

$P(X \geqslant 1) = 1 - P(X = 0) = 1 - (0{,}95)^{20} \approx 1 - 0{,}3585 = 0{,}6415$.

Il y a environ 64\,\% de chances de trouver au moins une pièce défectueuse.
\end{exemple}

% ── VII. Problèmes de seuil ──
\section{Problèmes de seuil}

\begin{methode}[Trouver $n$ pour atteindre un seuil de probabilité]
On cherche le plus petit entier $n$ tel que $P(X \geqslant 1) \geqslant \alpha$ (par exemple $\alpha = 0{,}95$).

On résout :
\begin{align*}
P(X \geqslant 1) \geqslant \alpha &\iff 1 - (1-p)^n \geqslant \alpha \\
&\iff (1-p)^n \leqslant 1 - \alpha \\
&\iff n \ln(1-p) \leqslant \ln(1-\alpha) \quad \text{(car } \ln(1-p) < 0\text{)} \\
&\iff n \geqslant \frac{\ln(1-\alpha)}{\ln(1-p)}
\end{align*}
\end{methode}

\begin{exemple}
Un test de dépistage a une probabilité $p = 0{,}01$ de détecter une anomalie chez une personne saine (faux positif). Combien de personnes saines faut-il tester pour que la probabilité d'observer au moins un faux positif dépasse 95\,\% ?

$n \geqslant \dfrac{\ln(0{,}05)}{\ln(0{,}99)} = \dfrac{-2{,}996}{-0{,}01005} \approx 298{,}1$

Il faut tester au moins $n = 299$ personnes saines.
\end{exemple}

% ── Bilan ──
\section*{Bilan du chapitre}

\begin{retenir}
\begin{itemize}[leftmargin=*]
  \item Deux événements $A$ et $B$ sont \textbf{indépendants} si $P(A \cap B) = P(A) \times P(B)$
  \item Une \textbf{épreuve de Bernoulli} a deux issues : succès ($p$) et échec ($1-p$)
  \item Un \textbf{schéma de Bernoulli} : $n$ épreuves de Bernoulli identiques et indépendantes
  \item \textbf{Loi binomiale} $X \sim \mathcal{B}(n,p)$ : $P(X = k) = \ds\binom{n}{k}\, p^k (1-p)^{n-k}$
  \item $E(X) = np$, \quad $V(X) = np(1-p)$, \quad $\sigma(X) = \sqrt{np(1-p)}$
  \item $P(X \geqslant 1) = 1 - (1-p)^n$
  \item \textbf{Seuil :} $n \geqslant \dfrac{\ln(1-\alpha)}{\ln(1-p)}$ pour $P(X \geqslant 1) \geqslant \alpha$
\end{itemize}
\end{retenir}


% ─────────────────────────────────────────────────────────────────────────────
% CHAPITRE 14 — SOMMES DE VARIABLES ALÉATOIRES
% ─────────────────────────────────────────────────────────────────────────────

\chapter{Sommes de variables aléatoires}

\section*{Introduction}

\textit{Un assureur couvre 500 contrats. Chaque contrat génère un coût aléatoire. Comment prévoir le coût total ? Un joueur lance 10 dés : comment calculer la moyenne des résultats, et avec quelle précision ?}

\medskip

Dès que l'on additionne des grandeurs aléatoires --- coûts, mesures, scores --- on a besoin de connaître les propriétés de la \textbf{somme de variables aléatoires}. Ce chapitre établit les règles fondamentales du calcul de l'espérance et de la variance d'une somme, et en tire des conséquences capitales : la justification des formules $E(X) = np$ et $V(X) = np(1-p)$ pour la loi binomiale, et les propriétés de la \textbf{moyenne d'échantillon}.

\medskip

\textbf{Objectifs du chapitre :}
\begin{itemize}[leftmargin=*]
  \item Utiliser la linéarité de l'espérance
  \item Calculer la variance d'une somme de variables indépendantes
  \item Retrouver l'espérance et la variance de la loi binomiale
  \item Connaître les propriétés de la moyenne d'un échantillon aléatoire
\end{itemize}

\begin{histoire}[Gauss, Legendre et la théorie des erreurs]
Au début du XIX\ieme{} siècle, les astronomes cherchent à déterminer les orbites des planètes à partir de mesures entachées d'erreurs. \textbf{Adrien-Marie Legendre} (1805) et \textbf{Carl Friedrich Gauss} (1809) développent indépendamment la \textbf{méthode des moindres carrés}, qui repose sur la somme de variables aléatoires. Plus tard, le statisticien belge \textbf{Adolphe Quetelet} (1796--1874) applique ces idées à la société : il introduit le concept d'« homme moyen » et montre que les caractéristiques humaines (taille, poids) suivent des lois statistiques régulières lorsqu'on moyenne un grand nombre d'observations.
\end{histoire}

% ── I. Linéarité de l'espérance ──
\section{Linéarité de l'espérance}

\begin{propriete}[Espérance d'une somme]
Soient $X$ et $Y$ deux variables aléatoires (pas nécessairement indépendantes). Alors :
\[ \boxed{E(X + Y) = E(X) + E(Y)} \]
Plus généralement, pour $X_1, X_2, \ldots, X_n$ :
\[ E(X_1 + X_2 + \cdots + X_n) = E(X_1) + E(X_2) + \cdots + E(X_n) \]
\end{propriete}

\begin{attention}[Aucune condition d'indépendance]
La linéarité de l'espérance est \textbf{toujours vraie}, que les variables soient indépendantes ou non. C'est une propriété remarquable et très puissante.
\end{attention}

\begin{propriete}[Espérance d'un multiple]
Pour toute variable aléatoire $X$ et tout réel $a$ :
\[ E(aX) = a \cdot E(X) \]
Plus généralement, pour $a, b \in \R$ :
\[ E(aX + b) = a \cdot E(X) + b \]
\end{propriete}

\begin{exemple}[Somme de deux dés]
On lance deux dés équilibrés. Soit $X_1$ et $X_2$ les résultats, et $S = X_1 + X_2$.

$E(X_1) = E(X_2) = \dfrac{1+2+3+4+5+6}{6} = 3{,}5$

$E(S) = E(X_1) + E(X_2) = 3{,}5 + 3{,}5 = 7$

En moyenne, la somme de deux dés vaut 7. Pas besoin de calculer les 36 cas !
\end{exemple}

\begin{exemple}[Gain net d'un jeu]
Un joueur paye 2\,\texteuro{} pour lancer un dé. Il gagne 5\,\texteuro{} s'il obtient 6, et rien sinon. Son gain net est $G = 5X - 2$ où $X \sim \mathcal{B}(\tfrac{1}{6})$.

$E(G) = 5 \times E(X) - 2 = 5 \times \dfrac{1}{6} - 2 = \dfrac{5}{6} - 2 = -\dfrac{7}{6} \approx -1{,}17\,\text{\texteuro}$.

Le jeu est défavorable au joueur : il perd en moyenne 1,17\,\texteuro{} par partie.
\end{exemple}

% ── II. Variance d'une somme ──
\section{Variance d'une somme de variables indépendantes}

\begin{propriete}[Variance d'une somme --- cas indépendant]
Si $X$ et $Y$ sont deux variables aléatoires \textbf{indépendantes}, alors :
\[ \boxed{V(X + Y) = V(X) + V(Y)} \]
Plus généralement, si $X_1, \ldots, X_n$ sont \textbf{mutuellement indépendantes} :
\[ V(X_1 + X_2 + \cdots + X_n) = V(X_1) + V(X_2) + \cdots + V(X_n) \]
\end{propriete}

\begin{attention}[Condition d'indépendance indispensable]
Contrairement à l'espérance, l'additivité de la variance \textbf{nécessite l'indépendance}. Si les variables ne sont pas indépendantes, la formule peut être fausse. Par exemple, $V(X + X) = V(2X) = 4\,V(X) \neq 2\,V(X)$ en général.
\end{attention}

\begin{propriete}[Variance d'un multiple]
Pour toute variable aléatoire $X$ et tout réel $a$ :
\[ V(aX) = a^2 \cdot V(X) \qquad \text{et} \qquad \sigma(aX) = |a| \cdot \sigma(X) \]
De plus, l'ajout d'une constante ne change pas la variance : $V(X + b) = V(X)$.
\end{propriete}

\begin{exemple}[Somme de deux dés (suite)]
$V(X_1) = V(X_2) = E(X^2) - (E(X))^2 = \dfrac{1+4+9+16+25+36}{6} - 3{,}5^2 = \dfrac{91}{6} - \dfrac{49}{4} = \dfrac{35}{12}$

Les deux dés étant indépendants : $V(S) = V(X_1) + V(X_2) = \dfrac{35}{12} + \dfrac{35}{12} = \dfrac{35}{6} \approx 5{,}83$.

$\sigma(S) = \sqrt{\dfrac{35}{6}} \approx 2{,}42$.
\end{exemple}

% ── III. Application à la loi binomiale ──
\section{Application à la loi binomiale}

\begin{propriete}[Décomposition de la loi binomiale]
Si $X \sim \mathcal{B}(n, p)$, alors $X$ se décompose comme une somme de $n$ variables de Bernoulli indépendantes :
\[ X = X_1 + X_2 + \cdots + X_n \quad \text{où chaque } X_i \sim \mathcal{B}(p) \]
\end{propriete}

\begin{demonstration}[de $E(X) = np$ et $V(X) = np(1-p)$]
Chaque $X_i \sim \mathcal{B}(p)$ vérifie $E(X_i) = p$ et $V(X_i) = p(1-p)$.

\medskip

\textbf{Espérance :} par linéarité,
\[ E(X) = E(X_1) + E(X_2) + \cdots + E(X_n) = \underbrace{p + p + \cdots + p}_{n \text{ termes}} = np \]

\textbf{Variance :} les $X_i$ étant indépendantes,
\[ V(X) = V(X_1) + V(X_2) + \cdots + V(X_n) = \underbrace{p(1-p) + \cdots + p(1-p)}_{n \text{ termes}} = np(1-p) \]

D'où $\sigma(X) = \sqrt{np(1-p)}$. \qed
\end{demonstration}

\textit{Remarque :} cette démonstration est bien plus élégante que le calcul direct par la formule $E(X) = \sum_{k=0}^{n} k\,\binom{n}{k}\,p^k(1-p)^{n-k}$, qui nécessite des manipulations combinatoires techniques.

% ── IV. Échantillon et moyenne ──
\section{Échantillon et moyenne d'échantillon}

\subsection{Échantillon aléatoire}

\begin{definition}[Échantillon aléatoire]
Soit $X$ une variable aléatoire d'espérance $\mu = E(X)$ et de variance $\sigma^2 = V(X)$. Un \textbf{échantillon aléatoire} de taille $n$ est un $n$-uplet $(X_1, X_2, \ldots, X_n)$ de variables aléatoires \textbf{indépendantes et de même loi} que $X$.
\end{definition}

\subsection{Somme et moyenne}

\begin{definition}
On définit :
\begin{itemize}
  \item La \textbf{somme} : $S_n = X_1 + X_2 + \cdots + X_n$
  \item La \textbf{moyenne d'échantillon} : $M_n = \dfrac{S_n}{n} = \dfrac{X_1 + X_2 + \cdots + X_n}{n}$
\end{itemize}
\end{definition}

\begin{propriete}[Paramètres de la somme]
\[ E(S_n) = n\mu \qquad V(S_n) = n\sigma^2 \qquad \sigma(S_n) = \sigma\sqrt{n} \]
\end{propriete}

\begin{propriete}[Paramètres de la moyenne]
\[ \boxed{E(M_n) = \mu \qquad V(M_n) = \frac{\sigma^2}{n} \qquad \sigma(M_n) = \frac{\sigma}{\sqrt{n}}} \]
\end{propriete}

\begin{demonstration}
$M_n = \frac{1}{n} S_n$, donc :
\[ E(M_n) = \frac{1}{n} E(S_n) = \frac{1}{n} \times n\mu = \mu \]
\[ V(M_n) = \frac{1}{n^2} V(S_n) = \frac{1}{n^2} \times n\sigma^2 = \frac{\sigma^2}{n} \]
\[ \sigma(M_n) = \frac{1}{n} \sigma(S_n) = \frac{1}{n} \times \sigma\sqrt{n} = \frac{\sigma}{\sqrt{n}} \qed \]
\end{demonstration}

\begin{methode}[Interprétation de $\sigma(M_n) = \sigma / \sqrt{n}$]
La moyenne d'échantillon a la même espérance que $X$ (elle est \textbf{sans biais}), mais son écart-type est divisé par $\sqrt{n}$. Cela signifie que :
\begin{itemize}
  \item La moyenne se \textbf{concentre} autour de $\mu$ quand $n$ augmente.
  \item Pour diviser l'incertitude par 2, il faut \textbf{multiplier} la taille de l'échantillon par 4.
  \item Pour diviser l'incertitude par 10, il faut multiplier $n$ par 100.
\end{itemize}
La précision ne s'améliore qu'en $\sqrt{n}$ : c'est la \textbf{loi en racine de $n$}.
\end{methode}

\begin{exemple}[Assurance]
Un assureur couvre $n = 10\,000$ contrats. Le coût annuel par contrat a pour espérance $\mu = 400$\,\texteuro{} et écart-type $\sigma = 1\,200$\,\texteuro{} (forte variabilité individuelle).

\medskip

\textbf{Coût total :} $E(S_n) = 10\,000 \times 400 = 4\,000\,000$\,\texteuro{}, \quad $\sigma(S_n) = 1\,200 \times \sqrt{10\,000} = 120\,000$\,\texteuro{}.

\medskip

\textbf{Coût moyen par contrat :} $E(M_n) = 400$\,\texteuro{}, \quad $\sigma(M_n) = \dfrac{1\,200}{\sqrt{10\,000}} = 12$\,\texteuro{}.

\medskip

Bien que chaque contrat soit très aléatoire ($\sigma = 1\,200$\,\texteuro{}), le coût moyen par contrat est prévisible à 12\,\texteuro{} près. C'est le principe fondamental de l'assurance : la \textbf{mutualisation des risques}.
\end{exemple}

% ── Bilan ──
\section*{Bilan du chapitre}

\begin{retenir}
\begin{itemize}[leftmargin=*]
  \item \textbf{Linéarité de l'espérance} (toujours vraie) : $E(X+Y) = E(X) + E(Y)$, \quad $E(aX+b) = aE(X) + b$
  \item \textbf{Variance d'une somme} (si indépendantes) : $V(X+Y) = V(X) + V(Y)$
  \item \textbf{Variance d'un multiple :} $V(aX) = a^2 V(X)$
  \item \textbf{Loi binomiale :} $X = X_1 + \cdots + X_n$ avec $X_i \sim \mathcal{B}(p)$ $\Rightarrow$ $E(X) = np$, $V(X) = np(1-p)$
  \item \textbf{Moyenne d'échantillon :} $E(M_n) = \mu$, \quad $V(M_n) = \dfrac{\sigma^2}{n}$, \quad $\sigma(M_n) = \dfrac{\sigma}{\sqrt{n}}$
\end{itemize}
\end{retenir}


% ─────────────────────────────────────────────────────────────────────────────
% CHAPITRE 15 — CONCENTRATION ET LOI DES GRANDS NOMBRES
% ─────────────────────────────────────────────────────────────────────────────

\chapter{Concentration et loi des grands nombres}

\section*{Introduction}

\textit{Lorsqu'on lance un dé un grand nombre de fois, la fréquence d'apparition du 6 se rapproche de $1/6$. Lorsqu'on mesure la même grandeur physique de nombreuses fois, la moyenne des mesures se stabilise. Pourquoi ?}

\medskip

Ce phénomène de \textbf{stabilisation des moyennes}, observé empiriquement depuis l'Antiquité, est l'un des résultats les plus profonds des probabilités : la \textbf{loi des grands nombres}. Ce chapitre en donne une démonstration rigoureuse, fondée sur l'\textbf{inégalité de Bienaymé-Tchebychev}. Il s'agit d'un résultat théorique fondamental qui justifie l'utilisation des statistiques : c'est parce que les moyennes convergent que les sondages, les essais cliniques et le contrôle qualité ont un sens.

\medskip

\textbf{Objectifs du chapitre :}
\begin{itemize}[leftmargin=*]
  \item Connaître et utiliser l'inégalité de Bienaymé-Tchebychev
  \item Énoncer et comprendre la loi des grands nombres
  \item Appliquer l'inégalité de concentration à la moyenne d'échantillon
  \item Déterminer la taille d'échantillon nécessaire pour estimer une proportion
\end{itemize}

\begin{histoire}[De Bernoulli à Kolmogorov : trois siècles pour une preuve]
La loi des grands nombres a une histoire longue et riche.

\textbf{Jacques Bernoulli} (1713) en donne la première version dans l'\textit{Ars Conjectandi} : il montre que la fréquence observée d'un événement converge vers sa probabilité théorique, mais sa démonstration est très technique et limitée au cas des épreuves de Bernoulli.

En 1853, \textbf{Irénée-Jules Bienaymé} (1796--1878), statisticien français méconnu, établit une inégalité générale reliant la probabilité qu'une variable s'écarte de sa moyenne à sa variance. En 1867, le mathématicien russe \textbf{Pafnouti Tchebychev} (1821--1894) redécouvre et popularise cette inégalité, qui porte aujourd'hui leurs deux noms.

C'est finalement \textbf{Andreï Kolmogorov} (1903--1987) qui, en 1933, pose les fondements axiomatiques des probabilités et démontre la version la plus générale de la loi des grands nombres (convergence presque sûre).
\end{histoire}

% ── I. Inégalité de Bienaymé-Tchebychev ──
\section{Inégalité de Bienaymé-Tchebychev}

\subsection{Inégalité de Markov}

\begin{propriete}[Inégalité de Markov]
Soit $X$ une variable aléatoire \textbf{positive} d'espérance $E(X)$. Pour tout réel $a > 0$ :
\[ P(X \geqslant a) \leqslant \frac{E(X)}{a} \]
\end{propriete}

\textit{Intuition :} si la moyenne de $X$ est petite, alors $X$ ne peut pas être souvent grande. C'est une borne grossière mais universelle.

\subsection{Inégalité de Bienaymé-Tchebychev}

\begin{theoreme}[Inégalité de Bienaymé-Tchebychev]
Soit $X$ une variable aléatoire d'espérance $\mu = E(X)$ et de variance $V(X)$. Pour tout réel $\delta > 0$ :
\[ \boxed{P\!\left(\left|X - \mu\right| \geqslant \delta\right) \leqslant \frac{V(X)}{\delta^2}} \]
\end{theoreme}

\begin{demonstration}
La variable aléatoire $Y = (X - \mu)^2$ est positive, et $E(Y) = E\!\left[(X - \mu)^2\right] = V(X)$.

On applique l'inégalité de Markov à $Y$ avec le seuil $a = \delta^2 > 0$ :
\[ P(Y \geqslant \delta^2) \leqslant \frac{E(Y)}{\delta^2} = \frac{V(X)}{\delta^2} \]

Or $Y \geqslant \delta^2 \iff (X - \mu)^2 \geqslant \delta^2 \iff |X - \mu| \geqslant \delta$.

Donc $P(|X - \mu| \geqslant \delta) \leqslant \dfrac{V(X)}{\delta^2}$. \qed
\end{demonstration}

\begin{methode}[Lecture de l'inégalité]
L'inégalité de Bienaymé-Tchebychev dit : \textbf{la probabilité que $X$ s'écarte de son espérance d'au moins $\delta$ est majorée par $V(X)/\delta^2$}.

\medskip

De manière équivalente, en posant $\delta = k\sigma$ avec $\sigma = \sqrt{V(X)}$ :
\[ P(|X - \mu| \geqslant k\sigma) \leqslant \frac{1}{k^2} \]

\begin{center}
\begin{tabular}{c|ccc}
\toprule
$k$ (en nombre d'écarts-types) & 2 & 3 & 5 \\
\midrule
$P(|X - \mu| \geqslant k\sigma) \leqslant$ & $\dfrac{1}{4} = 25\,\%$ & $\dfrac{1}{9} \approx 11\,\%$ & $\dfrac{1}{25} = 4\,\%$ \\
\bottomrule
\end{tabular}
\end{center}
\end{methode}

\begin{attention}[Borne non optimale]
L'inégalité de Bienaymé-Tchebychev est une borne \textbf{universelle} : elle est valable pour \textbf{toute} variable aléatoire, sans hypothèse sur la forme de sa loi. En contrepartie, elle est souvent \textbf{très pessimiste}. Par exemple, pour une loi normale, $P(|X - \mu| \geqslant 2\sigma) \approx 4{,}6\,\%$, bien mieux que la borne de 25\,\%.

Sa force réside dans sa \textbf{généralité} : elle permet de démontrer des résultats théoriques sans connaître la loi exacte de $X$.
\end{attention}

\begin{exemple}
Le temps d'attente à un guichet a pour espérance $\mu = 8$ min et variance $V = 4$ min\textsuperscript{2}. Quelle est la probabilité d'attendre plus de 14 minutes ?

$P(|X - 8| \geqslant 6) \leqslant \dfrac{4}{36} = \dfrac{1}{9} \approx 11{,}1\,\%$

Donc $P(X \geqslant 14) \leqslant P(|X - 8| \geqslant 6) \leqslant 11{,}1\,\%$.
\end{exemple}

% ── II. Inégalité de concentration ──
\section{Inégalité de concentration}

\begin{propriete}[Inégalité de concentration pour la moyenne]
Soit $(X_1, \ldots, X_n)$ un échantillon de variables aléatoires indépendantes de même loi, d'espérance $\mu$ et de variance $\sigma^2$. Soit $M_n = \frac{1}{n}(X_1 + \cdots + X_n)$ la moyenne d'échantillon. Pour tout $\delta > 0$ :
\[ \boxed{P\!\left(\left|M_n - \mu\right| \geqslant \delta\right) \leqslant \frac{\sigma^2}{n\delta^2}} \]
\end{propriete}

\begin{demonstration}
On applique l'inégalité de Bienaymé-Tchebychev à $M_n$, dont l'espérance est $\mu$ et la variance est $\frac{\sigma^2}{n}$ (chapitre~14) :
\[ P(|M_n - \mu| \geqslant \delta) \leqslant \frac{V(M_n)}{\delta^2} = \frac{\sigma^2}{n\delta^2} \qed \]
\end{demonstration}

\begin{methode}[Interprétation]
Le membre de droite $\frac{\sigma^2}{n\delta^2}$ tend vers \textbf{0} quand $n \to +\infty$. Cela signifie que pour tout écart $\delta > 0$ fixé, aussi petit soit-il, la probabilité que $M_n$ s'écarte de $\mu$ de plus de $\delta$ devient \textbf{arbitrairement petite} quand $n$ est assez grand.

La moyenne d'échantillon se \textbf{concentre} autour de l'espérance théorique.
\end{methode}

\begin{exemple}
On lance un dé $n$ fois et on note $M_n$ la moyenne des résultats. On a $\mu = 3{,}5$ et $\sigma^2 = 35/12$.

Quelle est la probabilité que $M_n$ s'écarte de 3,5 de plus de 0,1 ?

$P(|M_n - 3{,}5| \geqslant 0{,}1) \leqslant \dfrac{35/12}{n \times 0{,}01} = \dfrac{35}{0{,}12\,n} \approx \dfrac{291{,}7}{n}$

\begin{center}
\begin{tabular}{c|cccc}
\toprule
$n$ & 1\,000 & 10\,000 & 100\,000 & 1\,000\,000 \\
\midrule
Borne & $29{,}2\,\%$ & $2{,}9\,\%$ & $0{,}29\,\%$ & $0{,}029\,\%$ \\
\bottomrule
\end{tabular}
\end{center}
\end{exemple}

% ── III. Loi des grands nombres ──
\section{Loi des grands nombres}

\begin{theoreme}[Loi (faible) des grands nombres]
Soit $(X_n)$ une suite de variables aléatoires \textbf{indépendantes et de même loi}, d'espérance $\mu$ et de variance $\sigma^2$ finie. La moyenne $M_n = \frac{1}{n}(X_1 + \cdots + X_n)$ \textbf{converge en probabilité} vers $\mu$ :
\[ \boxed{\forall\, \delta > 0, \quad \lim_{n \to +\infty} P\!\left(\left|M_n - \mu\right| \geqslant \delta\right) = 0} \]
\end{theoreme}

\begin{demonstration}
C'est une conséquence immédiate de l'inégalité de concentration :
\[ 0 \leqslant P(|M_n - \mu| \geqslant \delta) \leqslant \frac{\sigma^2}{n\delta^2} \xrightarrow[n \to +\infty]{} 0 \]
Par le théorème des gendarmes, $P(|M_n - \mu| \geqslant \delta) \to 0$. \qed
\end{demonstration}

\begin{methode}[Ce que dit (et ne dit pas) la loi des grands nombres]
\textbf{Ce qu'elle dit :}
\begin{itemize}
  \item La moyenne empirique se rapproche de l'espérance théorique quand $n$ augmente.
  \item La fréquence observée d'un événement converge vers sa probabilité.
  \item Les sondages, les moyennes expérimentales, les estimations statistiques sont \textbf{fiables} si $n$ est assez grand.
\end{itemize}

\textbf{Ce qu'elle ne dit pas :}
\begin{itemize}
  \item Elle ne donne \textbf{pas} la vitesse exacte de convergence (seulement une borne via Bienaymé-Tchebychev).
  \item Elle ne dit \textbf{pas} que chaque réalisation se rapproche de $\mu$, mais que la \textbf{moyenne} le fait.
  \item Elle ne dit \textbf{pas} qu'après une série de résultats extrêmes, les suivants « compensent » (c'est le \textbf{sophisme du joueur}).
\end{itemize}
\end{methode}

\begin{attention}[Le sophisme du joueur]
Si une pièce équilibrée donne 10 fois pile de suite, la probabilité d'obtenir pile au 11\ieme{} lancer reste $\frac{1}{2}$. La loi des grands nombres ne dit \textbf{pas} que face « doit » arriver pour compenser. Elle dit que l'effet des 10 premiers lancers est \textbf{noyé} dans la moyenne quand $n$ devient très grand, pas que les lancers futurs compensent les lancers passés.
\end{attention}

\subsection{Lien fréquence -- probabilité}

\begin{propriete}
Soit $A$ un événement de probabilité $p$. On répète $n$ fois l'expérience de manière indépendante et on note $F_n = \frac{X}{n}$ la fréquence observée de $A$, où $X \sim \mathcal{B}(n, p)$. Alors :
\[ \forall\, \delta > 0, \quad P(|F_n - p| \geqslant \delta) \leqslant \frac{p(1-p)}{n\delta^2} \leqslant \frac{1}{4n\delta^2} \]
car $p(1-p) \leqslant \frac{1}{4}$ pour tout $p \in [0,1]$.
\end{propriete}

\textit{La fréquence observée converge vers la probabilité théorique : c'est le \textbf{lien fondamental} entre statistique et probabilité.}

% ── IV. Application à la taille d'échantillon ──
\section{Application à la taille d'échantillon}

\begin{methode}[Déterminer la taille d'échantillon pour estimer une proportion]
On veut estimer une proportion $p$ inconnue avec une \textbf{précision} $\delta$ et un \textbf{niveau de confiance} $1 - \alpha$, c'est-à-dire :
\[ P(|F_n - p| < \delta) \geqslant 1 - \alpha \]

En utilisant l'inégalité de concentration avec $p(1-p) \leqslant \frac{1}{4}$ :
\[ P(|F_n - p| \geqslant \delta) \leqslant \frac{1}{4n\delta^2} \leqslant \alpha \]

On obtient la condition :
\[ \boxed{n \geqslant \frac{1}{4\alpha\delta^2}} \]
\end{methode}

\begin{exemple}[Sondage électoral]
On veut estimer une proportion $p$ à $\delta = 0{,}03$ (3 points de pourcentage) près, avec un niveau de confiance de 95\,\% ($\alpha = 0{,}05$).

$n \geqslant \dfrac{1}{4 \times 0{,}05 \times 0{,}03^2} = \dfrac{1}{4 \times 0{,}05 \times 0{,}0009} = \dfrac{1}{0{,}00018} \approx 5\,556$

Il faut interroger au moins $\mathbf{5\,556}$ personnes.

\medskip

\textit{Remarque :} cette borne est pessimiste (car l'inégalité de Bienaymé-Tchebychev est universelle). En utilisant l'approximation normale, on obtient $n \approx 1\,068$, ce qui est le résultat usuel des instituts de sondage.
\end{exemple}

\begin{exemple}[Contrôle qualité]
Un fabricant veut estimer le taux de défaut $p$ de sa production à 1\,\% près ($\delta = 0{,}01$), avec 99\,\% de confiance ($\alpha = 0{,}01$).

$n \geqslant \dfrac{1}{4 \times 0{,}01 \times 0{,}0001} = \dfrac{1}{0{,}000004} = 250\,000$

Avec la borne universelle, il faudrait 250\,000 pièces. C'est considérable, ce qui montre les limites de l'approche par Bienaymé-Tchebychev pour les applications pratiques.
\end{exemple}

% ── V. Méthode de Monte-Carlo ──
\section{Méthode de Monte-Carlo}

\begin{methode}[Estimation par simulation]
La \textbf{méthode de Monte-Carlo} exploite la loi des grands nombres pour estimer des quantités difficiles à calculer. Le principe est :
\begin{enumerate}
  \item Formuler la quantité à estimer comme une espérance $\mu = E(X)$.
  \item Simuler $n$ réalisations indépendantes $x_1, x_2, \ldots, x_n$ de $X$.
  \item Estimer $\mu$ par la moyenne $\bar{x} = \frac{1}{n}(x_1 + x_2 + \cdots + x_n)$.
\end{enumerate}
Par la loi des grands nombres, $\bar{x} \to \mu$ quand $n \to +\infty$.
\end{methode}

\begin{exemple}[Estimation de $\pi$]
On tire un point $(x, y)$ au hasard dans le carré $[0, 1]^2$. La probabilité qu'il tombe dans le quart de disque de rayon 1 est :
\[ p = \frac{\text{aire du quart de disque}}{\text{aire du carré}} = \frac{\pi/4}{1} = \frac{\pi}{4} \]

On définit $X = 1$ si $x^2 + y^2 \leqslant 1$, et $X = 0$ sinon. Alors $E(X) = p = \pi/4$.

En simulant $n$ tirages et en comptant la fréquence $F_n$ des points dans le quart de disque :
\[ \pi \approx 4 F_n \]

Par exemple, avec $n = 100\,000$ tirages, on obtient typiquement $\pi \approx 3{,}14$ avec 2 décimales correctes.

\medskip

La précision est de l'ordre de $\frac{1}{\sqrt{n}}$ : pour gagner une décimale, il faut multiplier $n$ par 100.
\end{exemple}

% ── Bilan ──
\section*{Bilan du chapitre}

\begin{retenir}
\begin{itemize}[leftmargin=*]
  \item \textbf{Bienaymé-Tchebychev :} $P(|X - \mu| \geqslant \delta) \leqslant \dfrac{V(X)}{\delta^2}$ \quad (toujours valable)
  \item \textbf{Concentration :} $P(|M_n - \mu| \geqslant \delta) \leqslant \dfrac{\sigma^2}{n\delta^2}$ \quad (tend vers 0)
  \item \textbf{Loi des grands nombres :} $M_n$ converge en probabilité vers $\mu$
  \item \textbf{Fréquence $\to$ probabilité :} $P(|F_n - p| \geqslant \delta) \leqslant \dfrac{1}{4n\delta^2}$
  \item \textbf{Taille d'échantillon :} $n \geqslant \dfrac{1}{4\alpha\delta^2}$ pour estimer $p$ à $\delta$ près avec confiance $1-\alpha$
  \item \textbf{Monte-Carlo :} simuler pour estimer une espérance, précision en $1/\sqrt{n}$
\end{itemize}
\end{retenir}

\end{document}
